\modulechapter{3}{Mô hình dữ liệu phân tích: thiết kế Gold cho Olist}
\moduleinfo{Cơ bản tới trung cấp}{30 giờ}{Module 1}

\begin{xltabular}{\textwidth}{@{}P{3.4cm} L@{}}
\toprule
\thd{thẻ module} & \thd{nội dung} \\
\midrule
\endhead
\bottomrule
\endfoot
Giờ & 30 giờ lõi: ba tuần 9 giờ và một tuần 3 giờ, theo nhịp 9 giờ module mỗi tuần; thêm khoảng 2,5 giờ cho các bài trong Đọc thêm, và khoảng 30 giờ nếu học cả khoá 4 của DeepLearning.AI \\
Dùng trong repo & 9 file CSV Olist trong \repo{source\_\allowbreak{}sim/\allowbreak{}data/}; \repo{docs/\allowbreak{}data-model.md} mục 5, 6, 9, 10 và 11; server \repo{source\_\allowbreak{}sim/\allowbreak{}app.py} chạy như một process ở cổng 8010 trong Lab 3.4, không cần Docker; tên bảng và tên cột của \repo{core.orders}, \repo{core.customers}, \repo{core.order\_\allowbreak{}items}, \repo{core.products}, \repo{core.sellers} và \repo{semantic.zip\_\allowbreak{}geo} cho Lab 3.5 \\
Phiên bản & DuckDB 1.5.6, pyarrow 25.0.1, FastAPI 0.146.0 và uvicorn 0.54.0 trong \repo{lakehouse/.venv}, Python 3.10.20; code cũng chạy trên DuckDB 2.0.0. Docs: \href{https://duckdb.org/docs/1.5/sql/statements/merge_into}{DuckDB 1.5: MERGE INTO}, \href{https://duckdb.org/docs/2.0/sql/statements/merge_into.html}{DuckDB 2.0: MERGE INTO}, \href{https://arrow.apache.org/docs/25.0/python/generated/pyarrow.Table.html}{pyarrow 25.0: Table} \\
Lướt được & Không nên bỏ bài nào. Bài 1 và Bài 10 đọc nhanh được nếu bạn đã quen chuẩn hoá schema và TDD. Kiểm 15 phút: chạy khối Làm thử của Bài 1 và khối Làm thử của Bài 10, rồi giải thích vì sao tiền kiểu \code{DOUBLE} làm hàng trăm nhóm lệch, và vì sao test đối chiếu đỏ khi chỉ thêm hai dòng. Giải thích được cả hai thì đọc nhanh hai bài đó \\
\end{xltabular}
% Kiểm: lakehouse/.venv/bin/python -c "import duckdb, pyarrow, fastapi, uvicorn" in 1.5.6, 25.0.1, 0.146.0, 0.54.0; Python 3.10.20; 06/11/2026.
% Kiểm: trang docs DuckDB 1.5 và 2.0 cho MERGE INTO trả 200 ngày 06/11/2026; chỉ trang 2.0 có câu "cannot affect the same target row more than once".
% =====================================================================
%  Bản 2 của Module 3 theo kiểu tutorial.
%  Mỗi khối code có dòng "% chạy: <id>"; output thật nằm sau "% kết quả: <id>".
%  Output chạy lại ngày 06/11/2026 bằng lakehouse/.venv/bin/python, Python 3.10.20,
%  DuckDB 1.5.6, từ một thư mục chứa symlink lakehouse trỏ tới repo. Mọi khối cũng chạy
%  được trên DuckDB 2.0.0 với Python 3.11; output giống hệt, trừ số byte ở Bài 8.
%  Khối ghi "cần-lời-giải" chỉ chạy được sau khi người đọc viết code của mình; lời giải
%  tham khảo dùng để kiểm output nằm ở solutions/m03/.
% =====================================================================

Module 1 dạy bạn hỏi dữ liệu bằng SQL trên chính schema của app. Module này dạy bạn dựng lại dữ liệu đó thành một dạng khác, gọi là \textit{star schema}, để câu hỏi phân tích dễ viết và số liệu không sai.

Bạn học bằng cách làm trên bộ Olist. Mỗi khái niệm đi kèm một lỗi thật đo được trên dữ liệu, một đoạn code chạy được, và chỗ nó giống hoặc khác thứ bạn đã quen khi làm phần mềm. Sản phẩm cuối là một tầng Gold chạy thử cho Olist trên DuckDB: một star schema có bộ test tự động, một accumulating snapshot dựng từ API của repo, và năm mart cho năm câu hỏi của \repo{docs/\allowbreak{}data-model.md}, kèm phương trình đối soát. Đó là mốc capstone 1.

\begin{ghichu}
\textbf{Cách chạy code trong module này.}
\begin{itemize}
  \item Lời giải mẫu nằm trong \repo{solutions/m03/}. Đọc sau khi tự làm, để đối chiếu.
  \item Mọi khối code chạy từ thư mục cha của \repo{lakehouse/} bằng \repo{lakehouse/.venv/bin/python}, trong đó có DuckDB 1.5.6. DuckDB là database phân tích chạy ngay trong process Python, không cần server, nên module này không cần Docker.
  \item Lộ trình ghim DuckDB 2.0.0, nhưng bản này cần Python 3.11 trở lên. Venv của \repo{lakehouse} dùng Python 3.10 nên có DuckDB 1.5.6, và mọi output trong module là của 1.5.6. Theo thông báo ra mắt bản 2.0, DuckDB 1.5 được hỗ trợ tới tháng 3/2027. Trên 2.0.0, mọi khối cho cùng output, trừ số byte ở Bài 8 lệch một chút.
  % Nguồn: https://pypi.org/pypi/duckdb/2.0.0/json (requires_python >=3.11.0); https://duckdb.org/2026/10/21/announcing-duckdb-200
  \item Dữ liệu là 9 file CSV Olist trong \repo{lakehouse/source\_\allowbreak{}sim/data/}. Nếu bạn để CSV ở chỗ khác, ví dụ \code{olist/} như ở Module 1, hãy sửa biến \code{D} ở đầu mỗi khối.
  \item Lab 3.4 chạy server giả lập nguồn của repo, \repo{source\_\allowbreak{}sim}, như một process Python ở cổng 8010, không cần Docker. Mọi file mà các lab tạo ra nằm trong \code{lab\_\allowbreak{}m03/}; không khối nào ghi vào \repo{lakehouse/}.
  \item Phần ``Kết quả'' dưới mỗi khối là output thật, chạy lại ngày 06/11/2026. Nếu máy bạn in ra số khác, dừng lại tìm lý do trước khi đọc tiếp; số byte ở Bài 8 là ngoại lệ, như mục trên.
\end{itemize}
\end{ghichu}

\begin{ghichu}
\textbf{Cú pháp DuckDB mà Module 1 chưa dạy.} Module 1 đã dạy \code{NOT EXISTS}, \code{FILTER} và \code{QUALIFY}. Các khối code ở đây dùng thêm:
\begin{itemize}
  \item \code{A ANTI JOIN B USING (k)} giữ các dòng của \code{A} không có dòng khớp ở \code{B}, giống \code{WHERE NOT EXISTS}.
  \item \code{GROUP BY ALL} nhóm theo mọi cột của \code{SELECT} không nằm trong aggregate. \code{ORDER BY ALL} sắp theo mọi cột, từ trái sang phải.
  \item \code{A EXCEPT B} trả các dòng có ở \code{A} mà không có ở \code{B}, sau khi bỏ dòng trùng. \code{SELECT * EXCLUDE (c)} chọn mọi cột trừ cột \code{c}.
  \item \code{range(a, b, bước) r(d)} sinh một bảng tên \code{r} có cột \code{d}, chạy từ \code{a} tới trước \code{b}. Giống \code{range()} của Python, điểm cuối không được tính.
\end{itemize}
% Nguồn: https://duckdb.org/docs/2.0/sql/query_syntax/from (mục Semi and Anti Joins); https://duckdb.org/docs/2.0/sql/query_syntax/groupby; https://duckdb.org/docs/2.0/sql/query_syntax/orderby; https://duckdb.org/docs/2.0/sql/query_syntax/setops
% Nguồn: https://duckdb.org/docs/2.0/sql/functions/timestamp (range: half open range); https://docs.python.org/3/library/stdtypes.html (range: r[i] < stop)
\end{ghichu}

% =====================================================================
\section{Vấn đề thật: doanh thu tăng 26,9\% mà không ai bán thêm gì}

Bạn là kỹ sư phần mềm vừa chuyển sang team dữ liệu. Một BA, tức người phân tích nghiệp vụ, nhắn: ``Cho mình tổng tiền khách đã trả, chia theo danh mục sản phẩm.''

Tiền thanh toán nằm ở bảng \code{order\_\allowbreak{}payments}, còn danh mục gắn với sản phẩm của từng dòng hàng trong \code{order\_\allowbreak{}items}. Bạn làm việc quen tay: join hai bảng theo \code{order\_\allowbreak{}id} rồi cộng \code{payment\_\allowbreak{}value}. Trước khi chia theo danh mục, bạn thử xem con số tổng.

% chạy: A
\begin{codeblock}
lakehouse/.venv/bin/python - <<'EOF'
import duckdb
D = 'lakehouse/source_sim/data/'
con = duckdb.connect()
for t in ['order_items', 'order_payments']:
    con.sql(f"CREATE VIEW {t} AS SELECT * FROM '{D}olist_{t}_dataset.csv'")
one = lambda q: con.sql(q).fetchall()[0][0]
dung = one("SELECT sum(CAST(payment_value AS DECIMAL(14,2))) FROM order_payments")
sai = one("""SELECT sum(CAST(p.payment_value AS DECIMAL(14,2)))
             FROM order_payments p JOIN order_items i USING (order_id)""")
print('chỉ cộng thanh toán:', dung)
print('join rồi mới cộng  :', sai, f'(cao hơn {sai / dung - 1:.1%})')
con.sql("""SELECT p.payment_value, count(*) AS so_dong_sau_join,
                  sum(CAST(p.payment_value AS DECIMAL(14,2))) AS tong_sau_join
           FROM order_payments p JOIN order_items i USING (order_id)
           WHERE order_id = '8272b63d03f5f79c56e9e4120aec44ef'
           GROUP BY 1""").show()
EOF
\end{codeblock}
\textit{Kết quả:}
% kết quả: A
\begin{codeblock}
chỉ cộng thanh toán: 16008872.12
join rồi mới cộng  : 20308134.71 (cao hơn 26.9%)
┌───────────────┬──────────────────┬───────────────┐
│ payment_value │ so_dong_sau_join │ tong_sau_join │
│    double     │      int64       │ decimal(38,2) │
├───────────────┼──────────────────┼───────────────┤
│        196.11 │               21 │       4118.31 │
└───────────────┴──────────────────┴───────────────┘
\end{codeblock}

Con số sau join cao hơn con số đúng 4.299.262,59 BRL. Không có lỗi, không có cảnh báo nào: database làm đúng điều bạn viết.

Bảng nhỏ ở cuối cho thấy lý do. Đơn \code{8272b63d...} trả tiền một lần, 196,11 BRL, nhưng có 21 dòng hàng. Join lặp khoản tiền đó cho từng dòng hàng, nên nó bị cộng 21 lần thành 4.118,31 BRL.

Lỗi này có tên là \textit{fan-out}, và nó là một lỗi rất hay gặp khi đọc thẳng schema của app để phân tích. Module này dạy cách thiết kế bảng để lỗi này khó xảy ra, và cách viết test để bắt nó khi nó vẫn xảy ra. Câu hỏi của cả module là: thiết kế tầng Gold cho Olist thế nào để mỗi con số chỉ có một định nghĩa, và một test đỏ ngay khi tổng tiền lệch dù chỉ 0,01 BRL?

% =====================================================================
\section{Từ SWE sang DE}

Bảng dưới nối thứ bạn đã biết với tên gọi trong data engineering, theo thứ tự các bài học. Cột cuối quan trọng nhất: so sánh nào cũng có chỗ gãy, và mỗi bài học sẽ chỉ ra chỗ gãy đó trên dữ liệu thật.

\begin{xltabular}{\textwidth}{@{}P{3.2cm} P{2.8cm} L L@{}}
\toprule
\thd{bạn đã biết (SWE)} & \thd{trong DE gọi là} & \thd{giống ở đâu} & \thd{khác ở đâu} \\
\midrule
\endhead
\bottomrule
\endfoot
schema do ORM và migration tạo ra & schema OLTP, đã chuẩn hoá & mỗi sự thật nằm ở một chỗ, ghi nhanh và an toàn & câu hỏi phân tích phải join nhiều bảng và dễ nhân dòng \\
read model trong CQRS & star schema ở tầng gold & một mô hình riêng cho việc đọc, dựng từ dữ liệu ghi & dựng theo lô từ nhiều nguồn, trả lời cả câu hỏi chưa ai đặt ra \\
quyết định một resource của API là gì & grain & khoá suy ra từ định nghĩa & grain là một câu nghiệp vụ viết trước khi chọn cột, và trong lakehouse không ai ép khoá giúp bạn \\
lỗi gộp nhiều aggregation trong \code{annotate()} của Django & fan-out & cùng cơ chế: join nhân dòng rồi mới cộng & \code{Count(distinct=True)} chữa được phép đếm; với tổng tiền thì \code{sum(DISTINCT)} sai, phải gộp từng bảng về grain trước khi join \\
bảng \code{events} hay audit log chỉ thêm & bảng fact & mỗi dòng là một chuyện đã xảy ra & có measure để cộng và khoá tới dimension; loại accumulating snapshot lại được sửa \\
counter, gauge, tỉ lệ lỗi trong monitoring & measure cộng được, cộng được một phần, không cộng được & cùng luật cộng, và luật đó phải ghi trong đặc tả & counter là số cộng dồn, nên phải lấy phần tăng giữa hai lần đo rồi mới cộng; measure của fact đã là con số của từng sự kiện \\
enum dùng chung giữa các service, ma trận service và capability & conformed dimension, bus matrix & một định nghĩa cho mọi nơi; nhìn một bảng biết ai dùng gì & dimension là bảng có dữ liệu và lịch sử, đổi nó làm đổi số liệu quá khứ; bus matrix còn quyết định thứ tự xây, mỗi lần một quy trình \\
tách bảng cho khỏi lặp chuỗi, hay nhúng sẵn như một document & star, snowflake, One Big Table & cùng câu hỏi: tách ra hay chép vào & Parquet đã lưu mỗi giá trị một lần trong từ điển của cột, nên tách không tiết kiệm; chép hết vào một bảng thì đổi tên \code{bed\_\allowbreak{}bath\_\allowbreak{}table} phải ghi lại 11.115 dòng (Bài 8) \\
câu \code{UPDATE} ghi đè; S3 Versioning & SCD Type 1, SCD Type 2 & giá trị mới nhất; hoặc nhiều version với một version hiện hành & Type 1 làm sai số liệu quá khứ; Type 2 chỉ thấy thay đổi mà pipeline kịp nhìn, và fact phải trỏ đúng version \\
bảng nối của \code{ManyToManyField}, bảng lịch sử chỉ thêm & Data Vault: hub, link, satellite & tách định danh, quan hệ và thuộc tính ra các bảng riêng & vault phục vụ việc nạp và kiểm toán; người phân tích vẫn cần một star dựng bên trên (Bài 11) \\
state machine của đơn: cột \code{status} bị \code{UPDATE} qua \code{created}, \code{approved}, \code{shipped} & accumulating snapshot & một dòng mỗi đơn, đi dần qua các mốc & app chỉ giữ trạng thái hiện tại, snapshot giữ mỗi mốc một cột; engine nào cũng cần lọc version trùng trong lô và chặn version cũ; trạng thái trung gian của mock là giả lập (Lab 3.4) \\
unit test, checksum, \code{Collectors.\allowbreak{}partitioningBy} của Java & data test, phương trình đối soát & assert tự động; hai nửa của một phép chia cộng lại bằng tổng & chạy trên dữ liệu thật nên đỏ có thể do nguồn; phương trình chỉ chứng minh không mất đồng nào, không chứng minh bộ lọc đúng (Bài 10, Bài 12) \\
\end{xltabular}

\textbf{Thói quen SWE gây hại nhất ở module này là tin rằng câu SQL chạy không lỗi thì con số đúng.} Fan-out, join làm rơi dòng, version cũ ghi đè version mới đều chạy mà không báo gì; chỉ test đối chiếu với nguồn mới bắt được chúng.

% =====================================================================
\section{Bạn sẽ làm được}

\begin{itemize}
  \item Viết grain của mỗi fact trong một câu, suy ra khoá, chứng minh khoá không trùng bằng một câu SQL, và phân loại measure kèm lý do.
  \item Dựng star schema Olist trên DuckDB gồm bốn dimension và hai fact, với bộ test viết trước trả về 0 dòng lỗi, rồi chứng minh fan-out trên chính star đó.
  \item Cài SCD Type 2 bằng SQL với khoảng hiệu lực nửa mở, chạy lại không sinh thêm dòng, và gắn mỗi dòng hàng với đúng một version người bán.
  \item Dựng accumulating snapshot của đơn hàng từ 391.551 version mà \repo{source\_\allowbreak{}sim} phát qua API, sao cho áp lại hay áp lộn thứ tự đều không đổi dòng nào.
  \item Định nghĩa metric doanh thu trong mart cho năm câu hỏi của repo, và chứng minh bằng phương trình đối soát rằng không đồng nào bị mất hay bị cộng hai lần.
  \item Tự chọn grain cho fact review và viết phương trình đối soát cho nó, không cần hướng dẫn từng bước.
\end{itemize}

% =====================================================================
\section{Bài 1: Schema của app và schema cho phân tích}

\textbf{Vấn đề.} Bộ Olist có 9 bảng, đúng kiểu một app thương mại điện tử: đơn, dòng hàng, thanh toán, đánh giá, khách, sản phẩm, người bán, bảng dịch tên danh mục và toạ độ theo zip. Câu hỏi ``doanh thu theo danh mục tiếng Anh và tháng'' cần 3 join qua 4 bảng và hai luật: dịch tên danh mục, cắt thời điểm mua về đầu tháng. Mỗi người viết câu này sẽ tự viết hai luật đó, và số của hai người có thể lệch nhau.

\textbf{Giải thích.} Schema của app là \textit{OLTP}, viết tắt của online transaction processing. Đó là database phục vụ app đang chạy: nhận rất nhiều giao dịch nhỏ, ưu tiên ghi nhanh và ghi đúng.

Schema OLTP được \textit{chuẩn hoá}: dữ liệu được tách ra nhiều bảng để mỗi sự thật chỉ nằm ở một chỗ, sửa một chỗ là xong. Mức hay gặp là 3NF, trong đó mỗi cột chỉ phụ thuộc vào khoá của bảng chứa nó.

Phân tích cần thứ ngược lại: quét hàng triệu dòng rồi gộp, với ít join và các luật được viết sẵn một lần. \textit{Star schema} làm đúng việc đó. Nó có một bảng \textit{fact} ở giữa, mỗi dòng là một sự kiện kèm các con số đo được. Xung quanh là các bảng \textit{dimension}, mỗi bảng mô tả một thứ như khách, sản phẩm hay ngày.

Mỗi dimension cách fact đúng một join. Dimension còn được làm phẳng sẵn: \code{dim\_\allowbreak{}product} chứa luôn tên danh mục tiếng Anh, nên người đọc không cần biết bảng dịch tồn tại.

\begin{codeblock}
Schema của app (OLTP):  ──<  nghĩa là một dòng bên trái, nhiều dòng bên phải

customers ──< orders ──< order_items >── products >── category_translation
                │             │
                ├──< order_payments        sellers ──< order_items
                └──< order_reviews

Star schema:  fact ở giữa, mỗi dimension cách fact đúng một join

                       dim_date
                          │
      dim_customer ── fct_order_items ── dim_product
                          │
                      dim_seller
\end{codeblock}

\begin{ghichu}
\textbf{So với SWE.} Schema OLTP chính là thứ ORM và migration của bạn tạo ra. Tài liệu Django gọi migration là cách đưa các thay đổi trên model vào schema của database.
% Nguồn: https://docs.djangoproject.com/en/6.1/topics/migrations/

Star schema giống \textit{read model} trong CQRS: Martin Fowler mô tả CQRS là dùng một model để cập nhật và một model khác để đọc. Fowler còn có khái niệm \textit{reporting database}: một bản sao dữ liệu vận hành đặt trong schema khác, không cần chuẩn hoá vì chỉ để đọc, và thường được nạp bằng một job chạy đêm. Star schema là một reporting database có luật thiết kế rõ ràng.
% Nguồn: https://martinfowler.com/bliki/CQRS.html
% Nguồn: https://martinfowler.com/bliki/ReportingDatabase.html

\textbf{Chỗ so sánh gãy.} Read model của CQRS thường được cập nhật sau mỗi lệnh ghi, với độ trễ nhỏ. Star schema phục vụ cả những câu hỏi chưa ai đặt ra, được dựng theo lô từ nhiều nguồn, và giữ lại lịch sử mà app đã ghi đè.
\end{ghichu}

\subsection{Làm thử}

Khối dưới dựng một star tối giản gồm \code{dim\_\allowbreak{}product}, \code{dim\_\allowbreak{}date} và \code{fct\_\allowbreak{}order\_\allowbreak{}items}, rồi chạy cùng một câu hỏi trên hai schema. Nó chạy hai lần: lần đầu tiền có kiểu \code{DECIMAL}, lần sau có kiểu \code{DOUBLE}. Khoá của \code{dim\_\allowbreak{}date} ở đây tạm là chính giá trị ngày; Bài 6 chuyển sang khoá số dạng \code{yyyymmdd}.

% chạy: B
\begin{codeblock}
lakehouse/.venv/bin/python - <<'EOF'
import duckdb
D = 'lakehouse/source_sim/data/'
con = duckdb.connect()
for t in ['orders', 'order_items', 'products']:
    con.sql(f"CREATE VIEW {t} AS SELECT * FROM '{D}olist_{t}_dataset.csv'")
con.sql(f"CREATE VIEW category_translation AS "
        f"SELECT * FROM '{D}product_category_name_translation.csv'")
one = lambda q: con.sql(q).fetchall()[0][0]

for money in ['DECIMAL(14,2)', 'DOUBLE']:
    # Star tối giản: luật dịch tên và luật cắt tháng nằm trong dimension
    con.sql("""CREATE OR REPLACE TABLE dim_product AS
      SELECT p.product_id, coalesce(t.product_category_name_english,
                                    p.product_category_name, 'unknown') AS category
      FROM products p
      LEFT JOIN category_translation t USING (product_category_name)""")
    con.sql("""CREATE OR REPLACE TABLE dim_date AS
      SELECT CAST(d AS DATE) AS date_key,
             CAST(date_trunc('month', d) AS DATE) AS month
      FROM range(DATE '2016-01-01', DATE '2019-01-01', INTERVAL 1 DAY) r(d)""")
    con.sql(f"""CREATE OR REPLACE TABLE fct_order_items AS
      SELECT i.order_id, i.order_item_id, i.product_id,
             CAST(o.order_purchase_timestamp AS DATE) AS date_key,
             CAST(i.price AS {money}) AS price
      FROM order_items i JOIN orders o USING (order_id)""")
    # Cùng một câu hỏi: doanh thu theo danh mục tiếng Anh và tháng
    oltp = f"""SELECT coalesce(t.product_category_name_english,
                              p.product_category_name, 'unknown') AS category,
                     CAST(date_trunc('month', o.order_purchase_timestamp) AS DATE)
                       AS month,
                     sum(CAST(i.price AS {money})) AS revenue
              FROM order_items i
              JOIN orders o ON o.order_id = i.order_id
              JOIN products p ON p.product_id = i.product_id
              LEFT JOIN category_translation t
                ON t.product_category_name = p.product_category_name
              GROUP BY 1, 2"""
    star = """SELECT p.category, d.month, sum(f.price) AS revenue
              FROM fct_order_items f
              JOIN dim_product p USING (product_id)
              JOIN dim_date d USING (date_key)
              GROUP BY 1, 2"""
    rows = one(f"SELECT count(*) FROM ({star})")
    a_not_b = one(f"SELECT count(*) FROM (({oltp}) EXCEPT ({star}))")
    b_not_a = one(f"SELECT count(*) FROM (({star}) EXCEPT ({oltp}))")
    print(money, '| số dòng:', rows, '| có nhóm lệch:', a_not_b + b_not_a > 0)
con.sql("""SELECT p.category, d.month,
                  sum(CAST(f.price AS DECIMAL(14,2))) AS revenue
           FROM fct_order_items f JOIN dim_product p USING (product_id)
           JOIN dim_date d USING (date_key)
           GROUP BY 1, 2 ORDER BY revenue DESC LIMIT 3""").show()
EOF
\end{codeblock}
\textit{Kết quả:}
% kết quả: B
\begin{codeblock}
DECIMAL(14,2) | số dòng: 1283 | có nhóm lệch: False
DOUBLE | số dòng: 1283 | có nhóm lệch: True
┌───────────────┬────────────┬───────────────┐
│   category    │   month    │    revenue    │
│    varchar    │    date    │ decimal(38,2) │
├───────────────┼────────────┼───────────────┤
│ watches_gifts │ 2018-05-01 │     123872.66 │
│ health_beauty │ 2018-08-01 │     120803.94 │
│ health_beauty │ 2018-06-01 │     107908.82 │
└───────────────┴────────────┴───────────────┘
\end{codeblock}

Với \code{DECIMAL}, hai cách cho cùng 1.283 nhóm, và \code{EXCEPT} theo cả hai chiều không còn dòng nào: hai schema trả lời giống hệt nhau. Câu trên star schema chỉ có 2 join và không còn luật nào, vì luật đã nằm trong dimension.

Đó là nghĩa của chữ ``nhất quán'' trong phân tích: mọi BA nhận cùng một định nghĩa, vì định nghĩa chỉ được viết một lần.

\begin{luuy}
\textbf{Bẫy thường gặp.}
\begin{itemize}
  \item \textbf{Để tiền ở kiểu \code{DOUBLE}.} Khoảng 680 trên 1.283 nhóm lệch ở vài chữ số cuối, vì cộng số thực theo thứ tự khác nhau cho kết quả hơi khác. Luật ``không dùng float cho tiền'' của SWE vẫn đúng.
  % Kiểm: chạy lại khối B với tiền DOUBLE, 10 luồng, DuckDB 1.5.6, 06/11/2026: bốn lần chạy cho 689, 674, 689, 689 nhóm lệch; tổng star của watches_gifts 2018-05 là 123872.65999999979 ở ba lần, 123872.65999999955 ở một lần; OLTP luôn 123872.66000000009. Ngày 05/11/2026 với SET threads = 1: cả hai 123872.65999999955, lệch 0.
  \item \textbf{Tin số nhóm lệch là cố định.} Nó đổi theo số luồng và giữa các lần chạy, nên khối ở trên chỉ in có lệch hay không. Với \code{SET threads = 1}, thứ tự cộng cố định nên hết lệch.
  % Kiểm: khối B thêm SET threads = 1 ngay sau connect, chạy ba lần ngày 06/11/2026, DuckDB 1.5.6: dòng DOUBLE in "có nhóm lệch: False" cả ba lần.
  \item \textbf{Nghĩ star schema thay thế OLTP.} App vẫn ghi vào OLTP. Star schema được dựng lại từ dữ liệu đó, ở một tầng riêng cho phân tích.
  \item \textbf{Coi việc lặp tên danh mục trong \code{dim\_\allowbreak{}product} là lỗi.} Lặp là cố ý. Bài 8 đo xem nó tốn bao nhiêu byte.
\end{itemize}
\end{luuy}

% =====================================================================
\section{Bài 2: Quy trình nghiệp vụ và grain}

\textbf{Vấn đề.} Ai đó yêu cầu ``một bảng doanh thu''. Một dòng của bảng đó là gì: một đơn, một món hàng, hay một lần trả tiền? Nếu không trả lời câu này trước, cột nào cũng có vẻ hợp lý để thêm vào, và \code{payment\_\allowbreak{}value} sẽ nằm cạnh \code{price} trên cùng một dòng. Đó chính là lỗi ở phần mở đầu.

\textbf{Giải thích.} Phương pháp dùng trong module này là của Ralph Kimball, thường gọi tắt là Kimball. Kimball chia việc thiết kế thành bốn bước, làm đúng thứ tự:
% Nguồn: https://www.kimballgroup.com/wp-content/uploads/2013/08/2013.09-Kimball-Dimensional-Modeling-Techniques11.pdf (trang 1 và 2)
\begin{enumerate}
  \item Chọn \textit{quy trình nghiệp vụ}: một hoạt động có thật mà mỗi lần xảy ra lại sinh ra con số, như nhận một đơn hàng. Đừng chọn theo phòng ban hay theo báo cáo.
  \item Khai báo \textit{grain}: một câu nói chính xác một dòng fact là gì.
  \item Chọn dimension: ai, cái gì, ở đâu, khi nào.
  \item Chọn fact: các con số đo đúng ở grain đó.
\end{enumerate}

Grain là bước then chốt. Kimball gọi câu grain là một hợp đồng ràng buộc cả thiết kế: mọi dimension và mọi fact phải khớp với nó, nên nó phải được viết ra trước.

Hãy bắt đầu ở \textit{atomic grain}, tức mức chi tiết nhất mà quy trình ghi nhận. Bảng đã gộp sẵn, như doanh thu theo ngày, chỉ trả lời được những câu hỏi đã đoán trước. Mỗi grain là một bảng riêng; hai grain không bao giờ nằm chung một fact.

Trong Olist, module này chọn bốn quy trình. Mỗi quy trình cho một fact với grain riêng:

\begin{xltabular}{\textwidth}{@{}P{2.1cm} L P{2.6cm} r P{4.2cm}@{}}
\toprule
\thd{quy trình} & \thd{một sự kiện là} & \thd{nguồn} & \thd{số dòng} & \thd{fact và khoá} \\
\midrule
\endhead
\bottomrule
\endfoot
đặt hàng & một dòng hàng được mua & \code{order\_\allowbreak{}items} & 112.650 & \code{fct\_\allowbreak{}order\_\allowbreak{}items}: \code{order\_\allowbreak{}id}, \code{order\_\allowbreak{}item\_\allowbreak{}id} \\
thanh toán & một lần trả tiền cho một đơn & \code{order\_\allowbreak{}payments} & 103.886 & \code{fct\_\allowbreak{}payments}: \code{order\_\allowbreak{}id}, \code{payment\_\allowbreak{}sequential} \\
giao hàng & một đơn đi qua các mốc từ lúc đặt tới lúc giao & \code{orders} & 99.441 & \code{fct\_\allowbreak{}orders}: \code{order\_\allowbreak{}id} \\
đánh giá & một review của khách & \code{order\_\allowbreak{}reviews} & 99.224 & \code{fct\_\allowbreak{}reviews}: bạn tự xác định ở thử thách kỹ năng \\
\end{xltabular}

\begin{ghichu}
\textbf{So với SWE.} Grain giống lúc bạn thiết kế một resource của REST API. \code{GET /orders/\{id\}} trả về một đơn, còn \code{GET /orders/\{id\}/items/\{n\}} trả về một dòng hàng. Khi đã quyết định một resource là gì, id của nó tự suy ra. Grain cũng vậy: câu ``một dòng là một dòng hàng trong một đơn'' cho ra khoá gồm \code{order\_\allowbreak{}id} và \code{order\_\allowbreak{}item\_\allowbreak{}id}.

\textbf{Chỗ so sánh gãy.} Primary key chỉ nói dòng nào là duy nhất; grain nói dòng đó đo sự kiện gì. Ngoài ra, trong app thì database chặn dòng trùng khoá giúp bạn. Trong lakehouse thì không: spec của Iceberg ghi rõ tính duy nhất theo khoá định danh không được Iceberg bảo đảm hay đòi hỏi. Bạn phải tự kiểm bằng truy vấn, như khối dưới.
% Nguồn: https://iceberg.apache.org/spec/#identifier-field-ids
\end{ghichu}

\subsection{Làm thử}

Mỗi câu grain được viết thành một khoá. Nếu câu grain đúng, số dòng của bảng phải bằng số khoá khác nhau.

% chạy: C
\begin{codeblock}
lakehouse/.venv/bin/python - <<'EOF'
import duckdb
D = 'lakehouse/source_sim/data/'
con = duckdb.connect()
for t in ['orders', 'order_items', 'order_payments']:
    con.sql(f"CREATE VIEW {t} AS SELECT * FROM '{D}olist_{t}_dataset.csv'")
# grain viết thành khoá: số dòng phải bằng số khoá khác nhau
checks = [('order_items', 'order_id, order_item_id'),
          ('order_payments', 'order_id, payment_sequential'),
          ('orders', 'order_id'),
          ('order_items', 'order_id')]   # thử sai: grain "một dòng mỗi đơn"
for table, key in checks:
    rows, keys = con.sql(f"SELECT count(*), count(DISTINCT ({key})) "
                         f"FROM {table}").fetchall()[0]
    print(f'{table}, khoá ({key}): {rows} dòng, {keys} khoá')
# Dạng test: trả về các khoá bị trùng; 0 dòng là đạt
con.sql("""SELECT order_id, order_item_id, count(*) AS n
           FROM order_items GROUP BY ALL HAVING count(*) > 1""").show()
EOF
\end{codeblock}
\textit{Kết quả:}
% kết quả: C
\begin{codeblock}
order_items, khoá (order_id, order_item_id): 112650 dòng, 112650 khoá
order_payments, khoá (order_id, payment_sequential): 103886 dòng, 103886 khoá
orders, khoá (order_id): 99441 dòng, 99441 khoá
order_items, khoá (order_id): 112650 dòng, 98666 khoá
┌──────────┬───────────────┬───────┐
│ order_id │ order_item_id │   n   │
│ varchar  │     int64     │ int64 │
└──────────┴───────────────┴───────┘
               0 rows
\end{codeblock}

Ba câu grain đầu đúng trên dữ liệu nguồn. Dòng thứ tư là một grain sai cố ý: nếu bạn tuyên bố \code{order\_\allowbreak{}items} có ``một dòng mỗi đơn'', 112.650 dòng chỉ có 98.666 khoá khác nhau.

Câu \code{GROUP BY} kèm \code{HAVING count(*) > 1} ở cuối là dạng chuẩn của một \textit{test grain}. Nó trả về các khoá bị trùng, và 0 dòng nghĩa là đạt. Bạn sẽ gặp lại đúng dạng này ở Lab 3.2.

\begin{luuy}
\textbf{Bẫy thường gặp.}
\begin{itemize}
  \item \textbf{Thiết kế theo báo cáo.} Mỗi dashboard một bảng, mỗi bảng tính ``doanh thu'' một kiểu, rồi số giữa các dashboard lệch nhau. Kimball làm ngược lại: một fact gắn với một sự kiện có thật, không gắn với một báo cáo.
  \item \textbf{Nhầm grain với primary key.} Hai bảng có thể cùng dùng \code{order\_\allowbreak{}id} trong khoá mà đo hai thứ khác nhau: một bảng ghi các mốc giao của mỗi đơn, một bảng ghi trạng thái của mỗi đơn vào cuối mỗi ngày.
  \item \textbf{Gộp quá sớm.} Bảng doanh thu theo ngày chạy nhanh, nhưng không trả lời được câu ``doanh thu theo người bán''. Giữ fact ở atomic grain, gộp sau.
\end{itemize}
\end{luuy}

% =====================================================================
\section{Bài 3: Fan-out, lỗi nhân dòng}

\textbf{Vấn đề.} Phần mở đầu cho thấy join thanh toán với dòng hàng rồi cộng làm tổng tiền tăng 26,9\%. Bài này trả lời hai câu: sửa thế nào cho đúng, và vì sao vài cách sửa trông có vẻ đúng vẫn sai.

\textbf{Giải thích.} Khi join theo quan hệ một nhiều, mỗi dòng bên ``một'' bị lặp lại cho mỗi dòng khớp bên ``nhiều''. Nếu sau đó bạn cộng một cột của bên ``một'', mỗi giá trị bị cộng nhiều lần.

Luật sửa rất ngắn: không cộng qua một join làm đổi grain. Gộp từng bảng về cùng một grain trước, ví dụ một dòng mỗi đơn, rồi mới join các kết quả đã gộp.

\begin{ghichu}
\textbf{So với SWE.} Nếu bạn từng dùng Django ORM, bạn đã gặp lỗi này dưới một cái tên khác. Tài liệu Django ghi rằng gộp nhiều aggregation trong \code{annotate()} sẽ cho kết quả sai, vì ORM dùng join thay cho subquery. Ví dụ của họ là một cuốn sách có 2 tác giả và 3 cửa hàng. Gọi \code{annotate()} với hai \code{Count}, một trên \code{authors} và một trên \code{store}, thì cả hai số đều ra 6, vì 2 nhân 3 bằng 6. Đó chính là fan-out.
% Nguồn: https://docs.djangoproject.com/en/6.1/topics/db/aggregation/#combining-multiple-aggregations

\textbf{Chỗ so sánh gãy.} Django có cách chữa riêng cho phép đếm là tham số \code{distinct=True}, và chính tài liệu của họ viết rằng với phần lớn aggregate khác thì không có cách nào tránh vấn đề này trong \code{annotate()}. Django cũng không cảnh báo lúc chạy, lỗi chỉ được ghi trong tài liệu. Ở kho dữ liệu cũng vậy, và cách chữa kiểu distinct không dùng được cho tổng tiền: khối dưới cho thấy \code{sum(DISTINCT)} ra 7.700.442,27 BRL.
\end{ghichu}

\subsection{Làm thử}

Khối dưới thử cách sửa đúng và hai cách sửa trông có vẻ đúng.

% chạy: D
\begin{codeblock}
lakehouse/.venv/bin/python - <<'EOF'
import duckdb
D = 'lakehouse/source_sim/data/'
con = duckdb.connect()
for t in ['order_items', 'order_payments']:
    con.sql(f"CREATE VIEW {t} AS SELECT * FROM '{D}olist_{t}_dataset.csv'")
# Gộp từng bảng về grain "một dòng mỗi đơn" trước, rồi mới join
con.sql("""CREATE TABLE p AS SELECT order_id,
             sum(CAST(payment_value AS DECIMAL(14,2))) AS paid
           FROM order_payments GROUP BY 1""")
con.sql("""CREATE TABLE i AS SELECT order_id, count(*) AS n_items
           FROM order_items GROUP BY 1""")
one = lambda q: con.sql(q).fetchall()[0][0]
print('gộp trước, LEFT JOIN :',
      one("SELECT sum(paid) FROM p LEFT JOIN i USING (order_id)"))
print('gộp trước, INNER JOIN:',
      one("SELECT sum(paid) FROM p JOIN i USING (order_id)"))
print('đơn có tiền mà không có hàng:',
      one("SELECT count(*) FROM p ANTI JOIN i USING (order_id)"))
print('sum(DISTINCT) sau join:', one("""
  SELECT sum(DISTINCT CAST(p.payment_value AS DECIMAL(14,2)))
  FROM order_payments p JOIN order_items i USING (order_id)"""))
EOF
\end{codeblock}
\textit{Kết quả:}
% kết quả: D
\begin{codeblock}
gộp trước, LEFT JOIN : 16008872.12
gộp trước, INNER JOIN: 15846280.17
đơn có tiền mà không có hàng: 775
sum(DISTINCT) sau join: 7700442.27
\end{codeblock}

\begin{itemize}
  \item Gộp về grain đơn trước rồi \code{LEFT JOIN} cho đúng 16.008.872,12 BRL.
  \item Cùng cách đó nhưng dùng \code{INNER JOIN} cho 15.846.280,17 BRL, tức mất 162.591,95 BRL của 775 đơn có thanh toán mà không có dòng hàng nào. Đây là nửa còn lại của lỗi join: join không chỉ nhân dòng mà còn có thể làm rơi dòng, và cũng không báo gì.
  \item \code{sum(DISTINCT ...)} sau join cho 7.700.442,27 BRL, chưa tới một nửa số đúng. \code{DISTINCT} bỏ luôn những khoản tiền của các đơn khác nhau nhưng trùng số tiền.
\end{itemize}

\begin{luuy}
\textbf{Bẫy thường gặp.}
\begin{itemize}
  \item \textbf{Thêm \code{DISTINCT} cho ``hết trùng''.} Nó che lỗi ở số đếm và làm sai số tổng theo một cách khác.
  \item \textbf{Dùng \code{INNER JOIN} khi bên kia có thể thiếu dòng.} Hãy hỏi trước: có đơn nào không có dòng hàng không? Ở Olist có 775 đơn như vậy.
  \item \textbf{Nghĩ fan-out chỉ làm sai tổng.} Nó làm sai cả trung bình, vì dòng bị lặp được tính thêm trọng số. Bài 7 cho thấy điểm review trung bình của bang SP thành 4,13 thay vì 4,17 sau một join như vậy.
\end{itemize}
\end{luuy}

% =====================================================================
\section{Bài 4: Bảng fact, loại fact và loại measure}

\textbf{Vấn đề.} Không phải con số nào cũng cộng được. Tỉ lệ giao trễ của cả nước có phải là trung bình tỉ lệ của 27 bang không? Số đơn đang chờ giao vào cuối mỗi ngày cộng qua một tuần có nghĩa gì không? Và với câu hỏi ``mỗi bước giao hàng mất bao lâu'', một dòng fact nên là gì?

\textbf{Giải thích.} Bảng fact chứa các con số đo được, gọi là \textit{measure}, cùng với khoá trỏ sang các dimension. Kimball chia fact thành các loại theo cách một dòng được sinh ra:
% Nguồn: Kimball techniques PDF, trang 4 và 5

\begin{xltabular}{\textwidth}{@{}P{2.7cm} L P{2.6cm} L@{}}
\toprule
\thd{loại fact} & \thd{một dòng là} & \thd{dòng có bị sửa} & \thd{ví dụ Olist} \\
\midrule
\endhead
\bottomrule
\endfoot
transaction & một sự kiện tại một thời điểm & không, chỉ thêm & \code{fct\_\allowbreak{}order\_\allowbreak{}items}, \code{fct\_\allowbreak{}payments} \\
periodic snapshot & trạng thái của một thứ vào cuối mỗi kỳ, như mỗi ngày & không, mỗi kỳ thêm một lượt dòng & số đơn đang chờ giao của mỗi người bán vào cuối mỗi ngày \\
accumulating snapshot & một quy trình có các mốc biết trước, mỗi mốc một cột ngày & có, sửa khi quy trình qua từng mốc & \code{fct\_\allowbreak{}orders}: đặt, duyệt, giao cho hãng vận chuyển, giao tới khách, ngày dự kiến \\
factless & các dimension gặp nhau tại một thời điểm, không có con số & không & Olist không có loại này; ví dụ của Kimball là một học sinh có mặt ở một lớp vào một ngày \\
\end{xltabular}

Measure cũng có ba loại, theo việc nó cộng được theo những chiều nào:

\begin{xltabular}{\textwidth}{@{}P{3.3cm} P{3.0cm} L L@{}}
\toprule
\thd{loại measure} & \thd{cộng được theo} & \thd{ví dụ Olist} & \thd{cách gộp đúng} \\
\midrule
\endhead
\bottomrule
\endfoot
cộng được, \textit{additive} & mọi dimension & \code{price}, \code{freight\_\allowbreak{}value} & \code{sum} \\
cộng được một phần, \textit{semi-additive} & mọi dimension trừ thời gian & số đơn đang chờ giao vào cuối ngày & lấy giá trị cuối kỳ, hoặc trung bình theo thời gian \\
không cộng được, \textit{non-additive} & không chiều nào & \code{review\_\allowbreak{}score}, tỉ lệ giao trễ & cộng tử số và mẫu số riêng, chia ở bước cuối \\
\end{xltabular}

\begin{ghichu}
\textbf{So với SWE.} Fact transaction giống bảng \code{events} hay audit log chỉ thêm mà bạn từng viết: mỗi dòng là một chuyện đã xảy ra, và không ai sửa nó.

Ba loại measure giống ba loại số liệu bạn gặp khi làm monitoring. Measure cộng được giống từng lần tăng của một counter: cộng theo service hay theo giờ đều có nghĩa. Measure cộng được một phần giống một gauge, loại metric mà Prometheus định nghĩa là một giá trị có thể tăng giảm tuỳ ý. Cộng bộ nhớ của mọi pod tại một thời điểm thì có nghĩa, còn cộng bộ nhớ của một pod qua 24 giờ thì không. Measure không cộng được giống tỉ lệ lỗi: tỉ lệ lỗi của cả hệ thống là tổng số lỗi chia tổng số request, không phải trung bình tỉ lệ lỗi của từng service.
% Nguồn: https://prometheus.io/docs/concepts/metric_types/

\textbf{Chỗ so sánh gãy.} Accumulating snapshot không chỉ thêm: dòng của một đơn được sửa mỗi khi đơn qua một mốc, và Kimball nói đây là loại fact duy nhất được cập nhật như vậy. Ngoài ra, bảng events của app hay chứa một payload JSON; fact thì chứa measure có kiểu rõ ràng và khoá tới dimension.
\end{ghichu}

\subsection{Làm thử}

Khối đầu dựng \code{fct\_\allowbreak{}orders} dạng accumulating snapshot: một dòng mỗi đơn, đủ năm mốc ở bảng trên, mỗi mốc một cột ngày. Bảng còn có một cột độ trễ và điểm của review mới nhất. Khối này trả lời câu hỏi thứ nhất và thứ ba ở phần vấn đề; khối thứ hai trả lời câu hỏi thứ hai.

% chạy: E
\begin{codeblock}
lakehouse/.venv/bin/python - <<'EOF'
import duckdb
D = 'lakehouse/source_sim/data/'
con = duckdb.connect()
for t in ['orders', 'order_reviews', 'customers']:
    con.sql(f"CREATE VIEW {t} AS SELECT * FROM '{D}olist_{t}_dataset.csv'")
# Accumulating snapshot: một dòng mỗi đơn, mỗi mốc một cột ngày
con.sql("""CREATE TABLE fct_orders AS
  SELECT o.order_id, c.customer_state,
         CAST(o.order_purchase_timestamp AS DATE)      AS purchase_date,
         CAST(o.order_approved_at AS DATE)             AS approved_date,
         CAST(o.order_delivered_carrier_date AS DATE)  AS carrier_date,
         CAST(o.order_delivered_customer_date AS DATE) AS delivered_date,
         CAST(o.order_estimated_delivery_date AS DATE) AS estimated_date,
         CAST(o.order_delivered_customer_date AS DATE)
           - CAST(o.order_purchase_timestamp AS DATE)  AS days_to_customer,
         o.order_delivered_customer_date > o.order_estimated_delivery_date
           AS late_by_timestamp,
         r.review_score
  FROM orders o JOIN customers c USING (customer_id)
  LEFT JOIN (SELECT order_id, review_score FROM order_reviews
             QUALIFY row_number() OVER (PARTITION BY order_id
               ORDER BY review_creation_date DESC, review_answer_timestamp DESC,
                        review_id DESC) = 1) r USING (order_id)""")
con.sql("""SELECT count(*) AS da_giao,
                  count(*) FILTER (WHERE late_by_timestamp) AS tre_theo_gio,
                  count(*) FILTER (WHERE delivered_date > estimated_date)
                    AS tre_theo_ngay
           FROM fct_orders WHERE delivered_date IS NOT NULL""").show()
con.sql("""SELECT delivered_date > estimated_date AS tre, count(*) AS so_don,
                  round(avg(review_score), 2) AS diem_tb
           FROM fct_orders WHERE delivered_date IS NOT NULL
           GROUP BY 1 ORDER BY 1""").show()
# Tỉ lệ là measure không cộng được: tính từ tổng, đừng lấy trung bình tỉ lệ con
con.sql("""WITH s AS (
             SELECT customer_state,
                    count(*) FILTER (WHERE delivered_date > estimated_date) AS late,
                    count(*) AS delivered
             FROM fct_orders WHERE delivered_date IS NOT NULL GROUP BY 1)
           SELECT round(sum(late) / sum(delivered), 4) AS ti_le_tu_tong,
                  round(avg(late / delivered), 4)      AS trung_binh_27_bang
           FROM s""").show()
# Mỗi bước mất bao lâu: các mốc nằm trên cùng một dòng, nên chỉ cần trừ hai cột
con.sql("""SELECT
             round(avg(approved_date - purchase_date), 1) AS dat_toi_duyet,
             round(avg(carrier_date - approved_date), 1)  AS duyet_toi_van_chuyen,
             round(avg(delivered_date - carrier_date), 1) AS van_chuyen_toi_khach,
             round(avg(days_to_customer), 1)              AS ca_quy_trinh
           FROM fct_orders WHERE delivered_date IS NOT NULL""").show()
EOF
\end{codeblock}
\textit{Kết quả:}
% kết quả: E
\begin{codeblock}
┌─────────┬──────────────┬───────────────┐
│ da_giao │ tre_theo_gio │ tre_theo_ngay │
│  int64  │    int64     │     int64     │
├─────────┼──────────────┼───────────────┤
│   96476 │         7827 │          6535 │
└─────────┴──────────────┴───────────────┘

┌─────────┬────────┬─────────┐
│   tre   │ so_don │ diem_tb │
│ boolean │ int64  │ double  │
├─────────┼────────┼─────────┤
│ false   │  89941 │    4.29 │
│ true    │   6535 │    2.27 │
└─────────┴────────┴─────────┘

┌───────────────┬────────────────────┐
│ ti_le_tu_tong │ trung_binh_27_bang │
│    double     │       double       │
├───────────────┼────────────────────┤
│        0.0677 │             0.0915 │
└───────────────┴────────────────────┘

┌───────────────┬──────────────────────┬──────────────────────┬──────────────┐
│ dat_toi_duyet │ duyet_toi_van_chuyen │ van_chuyen_toi_khach │ ca_quy_trinh │
│    double     │        double        │        double        │    double    │
├───────────────┼──────────────────────┼──────────────────────┼──────────────┤
│           0.5 │                  2.7 │                  9.3 │         12.5 │
└───────────────┴──────────────────────┴──────────────────────┴──────────────┘
\end{codeblock}

\begin{itemize}
  \item So theo giờ cho 7.827 đơn trễ; so theo ngày cho 6.535 đơn. Chênh lệch là 1.292 đơn giao đúng ngày dự kiến nhưng sau 00:00, vì \code{order\_\allowbreak{}estimated\_\allowbreak{}delivery\_\allowbreak{}date} của Olist luôn là 00:00:00. Repo \repo{lakehouse} cũng so theo ngày vì lý do này, ở mục 6 của \repo{docs/data-model.md}.
  % Kiểm: count(*) FILTER (WHERE CAST(order_estimated_delivery_date AS TIME) <> '00:00:00') = 0 trên 99.441 đơn.
  \item Đơn giao trễ được 2,27 điểm review trung bình, đơn đúng hẹn được 4,29 điểm. Accumulating snapshot làm câu hỏi này dễ, vì mọi mốc của một đơn nằm trên cùng một dòng.
  \item Tỉ lệ trễ tính từ tổng là 6,77\%; trung bình tỉ lệ của 27 bang là 9,15\%. Trung bình của các tỉ lệ cho mỗi bang cùng một trọng số, nên bang AL, với 397 đơn đã giao và tỉ lệ trễ 21,4\%, kéo số chung lên như thể AL lớn bằng SP.
  \item Mỗi bước là phép trừ hai cột trên cùng một dòng: trung bình 0,5 ngày từ lúc đặt tới lúc duyệt, 2,7 ngày từ lúc duyệt tới lúc giao cho hãng vận chuyển, và 9,3 ngày từ đó tới tay khách. Nếu mỗi mốc là một dòng trong bảng events, bạn phải ghép các dòng của cùng một đơn trước khi trừ.
  \item Cột \code{days\_\allowbreak{}to\_\allowbreak{}customer} là một \textit{lag fact}. Kimball gợi ý lưu cho mỗi mốc một cột độ trễ tính từ mốc đầu tiên; khi đó độ trễ giữa hai mốc bất kỳ chỉ là hiệu của hai cột. Cả quy trình mất trung bình 12,5 ngày.
  % Nguồn: Kimball techniques PDF, trang 15, mục Lag/Duration Facts
\end{itemize}

Khối đầu của phần Làm thử dựng bảng một lần, từ trạng thái cuối của mọi đơn trong CSV. Lab 3.4 dựng chính bảng này theo cách pipeline làm: từ từng version của đơn, sửa dòng mỗi khi đơn qua một mốc.

Khối thứ hai dựng một periodic snapshot nhỏ cho cả sàn: mỗi ngày từ 07/05 tới 13/05/2018 một dòng, đếm số đơn đã duyệt mà khách chưa nhận vào cuối ngày đó. Đơn không có ngày giao tới khách không được tính, vì không biết nó hết chờ lúc nào.

% chạy: E2
\begin{codeblock}
lakehouse/.venv/bin/python - <<'EOF'
import duckdb
D = 'lakehouse/source_sim/data/'
con = duckdb.connect()
con.sql(f"CREATE VIEW orders AS SELECT * FROM '{D}olist_orders_dataset.csv'")
# Periodic snapshot: một dòng mỗi ngày. Đơn đang chờ vào cuối ngày d là đơn
# đã duyệt trong hoặc trước ngày d, và tới tay khách sau ngày d
con.sql("""CREATE TABLE fct_pending_daily AS
  SELECT CAST(r.d AS DATE) AS snapshot_date, count(o.order_id) AS pending_orders
  FROM range(DATE '2018-05-07', DATE '2018-05-14', INTERVAL 1 DAY) r(d)
  LEFT JOIN orders o
    ON CAST(o.order_approved_at AS DATE) <= r.d
   AND CAST(o.order_delivered_customer_date AS DATE) > r.d
  GROUP BY 1""")
con.sql("SELECT * FROM fct_pending_daily ORDER BY 1").show()
# arg_max(x, ngày): giá trị x ở ngày lớn nhất, tức giá trị cuối kỳ
con.sql("""SELECT sum(pending_orders) AS cong_7_ngay,
                  round(avg(pending_orders), 1) AS trung_binh_ngay,
                  arg_max(pending_orders, snapshot_date) AS cuoi_tuan
           FROM fct_pending_daily""").show()
EOF
\end{codeblock}
\textit{Kết quả:}
% kết quả: E2
\begin{codeblock}
┌───────────────┬────────────────┐
│ snapshot_date │ pending_orders │
│     date      │     int64      │
├───────────────┼────────────────┤
│ 2018-05-07    │           2746 │
│ 2018-05-08    │           2791 │
│ 2018-05-09    │           2782 │
│ 2018-05-10    │           2722 │
│ 2018-05-11    │           2626 │
│ 2018-05-12    │           2730 │
│ 2018-05-13    │           2901 │
└───────────────┴────────────────┘

┌─────────────┬─────────────────┬───────────┐
│ cong_7_ngay │ trung_binh_ngay │ cuoi_tuan │
│   int128    │     double      │   int64   │
├─────────────┼─────────────────┼───────────┤
│       19298 │          2756.9 │      2901 │
└─────────────┴─────────────────┴───────────┘
\end{codeblock}

\begin{itemize}
  \item Cộng 7 ngày được 19.298, nhưng ngày cao nhất của tuần chỉ có 2.901 đơn đang chờ. Một đơn chờ cả tuần bị đếm 7 lần, giống cộng bộ nhớ của một pod qua 24 giờ ở ghi chú trên.
  \item Measure cộng được một phần phải gộp theo thời gian bằng giá trị cuối kỳ, 2.901 đơn vào cuối ngày 13/05/2018, hoặc bằng trung bình, 2.756,9 đơn mỗi ngày. Ví dụ quen thuộc của Kimball là số dư tài khoản: cộng được theo mọi dimension trừ thời gian.
  % Nguồn: Kimball techniques PDF, trang 4, mục Additive, Semi-Additive, and Non-Additive Facts
\end{itemize}

\begin{luuy}
\textbf{Bẫy thường gặp.}
\begin{itemize}
  \item \textbf{So một timestamp với một ngày lưu dạng 00:00:00.} Đọc dữ liệu trước khi viết điều kiện: 1.292 đơn bị tính trễ oan.
  \item \textbf{Lấy trung bình của các tỉ lệ con.} Lưu tử số và mẫu số, cộng chúng, rồi mới chia.
  \item \textbf{Cộng periodic snapshot theo thời gian.} Số đơn đang chờ giao của 7 ngày cộng lại không phải số đơn đang chờ của tuần đó.
  \item \textbf{Để mốc chưa xảy ra thành NULL trong khoá.} Dùng dòng khoá 0 ``chưa biết'' của \code{dim\_\allowbreak{}date}, ở Bài 6.
\end{itemize}
\end{luuy}

% =====================================================================
\section{Bài 5: Dimension và khoá}

\textbf{Vấn đề.} Bạn được hỏi ``bao nhiêu phần trăm khách mua lại lần hai''. Bảng \code{customers} có cột \code{customer\_\allowbreak{}id}, nên bạn đếm số khách có hơn một đơn theo cột đó và được 0. Không ai mua lại sao? Thật ra Olist cấp một \code{customer\_\allowbreak{}id} mới cho mỗi đơn. Người mua thật được định danh bằng \code{customer\_\allowbreak{}unique\_\allowbreak{}id}.

\textbf{Giải thích.} Dimension là bảng mô tả, trả lời các câu ai, cái gì, ở đâu, khi nào. Cột của nó được dùng trong \code{WHERE} và \code{GROUP BY}. Bảng dimension ngắn và rộng: mỗi thực thể một dòng, nhiều cột chữ dễ đọc. Một dimension có ba loại khoá:
% Nguồn: Kimball techniques PDF, trang 4, 7 và 8
\begin{itemize}
  \item \textbf{Khoá tự nhiên}, \textit{natural key}: khoá đến từ hệ nguồn, như \code{product\_\allowbreak{}id}.
  \item \textbf{Khoá thay thế}, \textit{surrogate key}: số nguyên không mang nghĩa do kho dữ liệu tự cấp, bắt đầu từ 1. Kimball khuyên mọi dimension dùng loại khoá này làm khoá chính, vì khoá tự nhiên có thể đổi, có thể trùng giữa hai hệ nguồn, và sẽ lặp lại khi dimension giữ lịch sử ở Bài 9. Dimension ngày là ngoại lệ: nó dùng được khoá dạng \code{yyyymmdd}.
  \item \textbf{Khoá bền}, \textit{durable key}: khoá không đổi của một thực thể qua mọi version của nó. Ở Olist, \code{customer\_\allowbreak{}unique\_\allowbreak{}id} đóng vai này.
\end{itemize}

Mỗi dimension còn có một dòng đặc biệt, thường mang khoá 0, cho giá trị ``chưa biết''. Kimball viết rằng khoá ngoại NULL trên fact vi phạm toàn vẹn tham chiếu; trong thực tế nó còn làm \code{INNER JOIN} rơi dòng. Thay cho khoá NULL, dimension phải có một dòng mặc định cho trường hợp chưa biết hoặc không áp dụng.

Trong star schema chuẩn, các dimension không trỏ tới nhau. Quan hệ giữa hai dimension là một sự kiện, nên nó nằm trong fact. Câu ``người mua \code{2e875ea5} mua của người bán \code{7299e27e}'' chỉ có nghĩa khi gắn với một đơn, nên nó là một dòng của \code{fct\_\allowbreak{}order\_\allowbreak{}items}. Có ba ngoại lệ có tên riêng. Kimball khuyên tránh snowflake, dùng outrigger hạn chế, và dùng bridge khi một dimension có nhiều giá trị cho một dòng fact:
% Nguồn: Kimball techniques PDF, trang 9 và 18
\begin{itemize}
  \item \textbf{Snowflake}: dimension tách tiếp thành dimension con, như \code{dim\_\allowbreak{}product} trỏ sang \code{dim\_\allowbreak{}category}. Bài 8 đo nó.
  \item \textbf{Outrigger}: dimension trỏ sang một dimension khác, như ngày mua đầu tiên của người mua trỏ sang \code{dim\_\allowbreak{}date}.
  \item \textbf{Bridge}: bảng nối cho quan hệ nhiều nhiều, như một đơn trả bằng nhiều kiểu thanh toán.
\end{itemize}

Đây là cách đọc một đơn thật qua fact và dimension. Đơn \code{002f98c0f7efd42638ed6100ca699b42} có 2 dòng hàng; các id được cắt còn 8 ký tự. Fact chỉ chứa khoá và con số:

\begin{xltabular}{\textwidth}{@{}l l l l L r r@{}}
\toprule
\thd{dòng} & \thd{ngày đặt} & \thd{người mua} & \thd{sản phẩm} & \thd{người bán} & \thd{price} & \thd{freight} \\
\midrule
\endhead
\bottomrule
\endfoot
1 & 2017-08-04 & \code{2e875ea5} & \code{d41dc2f2} & \code{7299e27e} & 8,99 & 32,57 \\
2 & 2017-08-04 & \code{2e875ea5} & \code{880be32f} & \code{fa40cc5b} & 44,90 & 7,16 \\
\end{xltabular}

Nghĩa của các khoá nằm ở dimension. Sản phẩm \code{d41dc2f2} thuộc danh mục \code{consoles\_\allowbreak{}games}, nặng 450 g, có 1 ảnh; sản phẩm \code{880be32f} thuộc \code{toys}, nặng 200 g, có 4 ảnh. Người bán \code{7299e27e} ở araguari, bang MG; người bán \code{fa40cc5b} ở sao paulo, bang SP. Người mua \code{2e875ea5} ở porto alegre, bang RS.

Ghép lại, dòng 1 nghĩa là: một người mua ở porto alegre, bang RS, mua một món \code{consoles\_\allowbreak{}games} của người bán ở araguari, bang MG, ngày 2017-08-04, giá 8,99 BRL, phí vận chuyển 32,57 BRL.

\begin{ghichu}
\textbf{So với SWE.} Khoá thay thế có cùng lý do với việc bạn không lấy email làm primary key của bảng \code{users}: email có thể đổi. Cặp \code{customer\_\allowbreak{}id} và \code{customer\_\allowbreak{}unique\_\allowbreak{}id} giống một id gắn với từng đơn so với user id thật.

Dòng ``unknown'' là Null Object pattern: thay vì trả về NULL rồi buộc mọi nơi phải kiểm, bạn trả về một đối tượng mặc định có hành vi an toàn. Ở đây join không rơi dòng, và báo cáo hiện chữ ``unknown'' thay vì một ô trống.

\textbf{Chỗ so sánh gãy.} Trong app, một user có một id suốt đời. Trong dimension có lịch sử, một người mua có nhiều khoá thay thế, mỗi version một khoá, và chỉ khoá bền là không đổi. Bài 9 làm rõ điều này.

Null Object cũng gãy ở một chỗ: dòng unknown có thể nuốt lỗi âm thầm. Nếu bước tra khoá hỏng, fact vẫn nạp đủ dòng, chỉ là trỏ về khoá 0, nên cần thêm một test theo dõi tỉ lệ dòng mang khoá 0.
\end{ghichu}

\subsection{Làm thử}

% chạy: F
\begin{codeblock}
lakehouse/.venv/bin/python - <<'EOF'
import duckdb
D = 'lakehouse/source_sim/data/'
con = duckdb.connect()
for t in ['orders', 'customers']:
    con.sql(f"CREATE VIEW {t} AS SELECT * FROM '{D}olist_{t}_dataset.csv'")
# Ai là "một người mua"? Đếm số đơn theo từng loại id
con.sql("""WITH n AS (
             SELECT c.customer_id, c.customer_unique_id, count(*) OVER
                      (PARTITION BY c.customer_unique_id) AS orders_per_buyer
             FROM orders o JOIN customers c USING (customer_id))
           SELECT count(DISTINCT customer_id)        AS customer_id,
                  count(DISTINCT customer_unique_id) AS customer_unique_id,
                  count(DISTINCT customer_unique_id)
                    FILTER (WHERE orders_per_buyer > 1) AS mua_tu_2_lan
           FROM n""").show()
# dim_customer: khoá thay thế từ 1, dòng "unknown" mang khoá 0
con.sql("""CREATE TABLE dim_customer AS
  SELECT 0 AS customer_key, 'unknown' AS customer_unique_id,
         'unknown' AS customer_city, 'unknown' AS customer_state
  UNION ALL
  SELECT row_number() OVER (ORDER BY customer_unique_id), customer_unique_id,
         customer_city, customer_state
  FROM (SELECT c.customer_unique_id, c.customer_city, c.customer_state
        FROM orders o JOIN customers c USING (customer_id)
        QUALIFY row_number() OVER (PARTITION BY c.customer_unique_id
                ORDER BY o.order_purchase_timestamp DESC,
                         c.customer_id DESC) = 1)""")
con.sql("""SELECT count(*) AS so_dong, count(DISTINCT customer_key) AS so_khoa,
                  min(customer_key) AS khoa_min, max(customer_key) AS khoa_max
           FROM dim_customer""").show()
EOF
\end{codeblock}
\textit{Kết quả:}
% kết quả: F
\begin{codeblock}
┌─────────────┬────────────────────┬──────────────┐
│ customer_id │ customer_unique_id │ mua_tu_2_lan │
│    int64    │       int64        │    int64     │
├─────────────┼────────────────────┼──────────────┤
│       99441 │              96096 │         2997 │
└─────────────┴────────────────────┴──────────────┘

┌─────────┬─────────┬──────────┬──────────┐
│ so_dong │ so_khoa │ khoa_min │ khoa_max │
│  int64  │  int64  │  int64   │  int64   │
├─────────┼─────────┼──────────┼──────────┤
│   96097 │   96097 │        0 │    96096 │
└─────────┴─────────┴──────────┴──────────┘
\end{codeblock}

\begin{itemize}
  \item Có 99.441 \code{customer\_\allowbreak{}id} cho 99.441 đơn, nên mỗi \code{customer\_\allowbreak{}id} chỉ có một đơn, và tỉ lệ mua lại tính theo cột này luôn là 0.
  \item Theo \code{customer\_\allowbreak{}unique\_\allowbreak{}id} có 96.096 người mua, trong đó 2.997 người mua từ 2 lần trở lên, khoảng 3,1\%.
  \item \code{dim\_\allowbreak{}customer} có 96.097 dòng: 96.096 người mua cộng dòng khoá 0. Địa chỉ được lấy ở đơn mới nhất; Bài 9 cho thấy cách này làm sai thành phố của 130 đơn và sai bang của 43 đơn.
\end{itemize}

\begin{luuy}
\textbf{Bẫy thường gặp.}
\begin{itemize}
  \item \textbf{Dùng \code{customer\_\allowbreak{}id} của Olist làm người mua.} Kiểm bằng cách đếm giá trị khác nhau của hai cột trước khi chọn khoá.
  \item \textbf{Dùng khoá tự nhiên làm khoá chính của dimension có lịch sử.} Khi một người mua có hai version, khoá tự nhiên xuất hiện hai lần.
  \item \textbf{Để NULL trong khoá của fact.} Join với dimension sẽ làm rơi dòng đó khỏi mọi báo cáo. Dùng khoá 0.
  \item \textbf{Cho các dimension trỏ sang nhau ``cho tiện''.} Mỗi lần như vậy, người đọc phải join thêm một bảng, và quan hệ đáng lẽ nằm trong fact bị giấu đi.
\end{itemize}
\end{luuy}

% =====================================================================
\section{Bài 6: Dimension ngày, vai, degenerate và junk}

\textbf{Vấn đề.} Vài dimension không hề có trong app. App tính thứ trong tuần bằng code mỗi khi cần, còn phân tích lại cần nhóm theo tuần, tháng, ngày lễ ở hầu hết câu hỏi. Một đơn có năm mốc thời gian, và mốc nào cũng là một ngày. \code{order\_\allowbreak{}id} không còn thuộc tính nào sau khi khách, sản phẩm và người bán đã thành dimension riêng. Đơn còn vài cờ nhỏ như trạng thái, kiểu thanh toán, giao đúng hẹn hay không, có review hay không. Làm mỗi cờ một dimension thì fact sẽ có cả chục khoá.

\textbf{Giải thích.}
% Nguồn: Kimball techniques PDF, trang 7 tới 9
\begin{itemize}
  \item \textbf{Dimension ngày.} Mỗi ngày một dòng, với các cột như năm, tháng, thứ, cờ ngày lễ. Kimball viết rằng bạn không bao giờ muốn tính ngày lễ Phục sinh bằng SQL, mà tra nó trong dimension ngày. Khoá có thể là số \code{yyyymmdd}. Bảng cần một dòng khoá 0 cho ngày chưa biết, như ngày giao của đơn chưa giao.
  \item \textbf{Role-playing.} Một bảng \code{dim\_\allowbreak{}date} được fact tham chiếu nhiều lần, mỗi lần với một vai: ngày đặt, ngày giao, ngày dự kiến. Kimball khuyên mỗi vai là một view riêng với tên cột riêng, để hai vai dùng được trong cùng một câu truy vấn.
  \item \textbf{Degenerate dimension.} \code{order\_\allowbreak{}id} nằm thẳng trên fact và không có bảng \code{dim\_\allowbreak{}order}, vì sau khi tách hết thuộc tính thì đơn không còn gì để mô tả. Kimball dùng ví dụ số hoá đơn trên các dòng của hoá đơn.
  \item \textbf{Junk dimension.} Gom các cờ nhỏ vào một dimension, mỗi dòng là một tổ hợp cờ, và fact chỉ giữ một khoá. Kimball viết rằng bảng này chỉ nên chứa các tổ hợp thật sự xuất hiện trong dữ liệu. Hướng dẫn Power BI của Microsoft lại nói bảng này thường là tích Descartes của mọi giá trị, tức kết quả \code{CROSS JOIN} của tất cả các cờ. Bảng \code{CROSS JOIN} không bao giờ thiếu tổ hợp; cách của Kimball nhỏ hơn, nhưng pipeline phải thêm dòng khi gặp tổ hợp mới.
  % Nguồn: https://learn.microsoft.com/en-us/power-bi/guidance/star-schema (mục Junk dimensions)
\end{itemize}

\begin{ghichu}
\textbf{So với SWE.} Dimension ngày giống một hàm tiện ích về lịch mà cả team dùng chung, nhưng được lưu thành dữ liệu để mọi công cụ SQL nhận cùng một câu trả lời. Role-playing giống một class \code{Date} dùng cho nhiều field như \code{created\_\allowbreak{}at}, \code{shipped\_\allowbreak{}at}, \code{delivered\_\allowbreak{}at}.

Degenerate dimension giống \code{trace\_\allowbreak{}id} trong log: bạn nhóm theo nó để gom các dòng của cùng một request, nhưng không có bảng \code{traces} nào chứa thuộc tính của nó. Junk dimension giống flyweight pattern hay string interning: rất nhiều đối tượng dùng chung vài tổ hợp giá trị, nên mỗi tổ hợp được lưu một lần và đối tượng chỉ giữ một tham chiếu nhỏ.

\textbf{Chỗ so sánh gãy.} Hàm tiện ích chạy lúc được gọi; dimension ngày phải được nạp trước và có hạn. \code{dim\_\allowbreak{}date} của module này dừng ở 2018-12-31, nên một đơn đặt sau ngày đó sẽ không tìm thấy ngày của nó. Junk dimension theo cách của Kimball cũng vậy: tổ hợp mới phải được thêm vào trước khi fact dùng tới.

Với \code{trace\_\allowbreak{}id}, chỗ gãy nằm ở lý do không có bảng. Đơn hàng vốn có bảng \code{orders} ở nguồn; star bỏ \code{dim\_\allowbreak{}order} vì ngày đặt, người mua và các thuộc tính khác của đơn đã thành khoá tới các dimension khác trên từng dòng fact.
% Nguồn: Kimball techniques PDF, trang 7, mục Degenerate Dimensions
\end{ghichu}

\subsection{Làm thử}

Khối đầu dựng \code{dim\_\allowbreak{}date} có dòng khoá 0, tạo hai vai ngày đặt và ngày giao, rồi dùng cả hai vai trong một câu: đơn đặt trong tháng 12/2017 được giao vào tháng nào?

% chạy: G
\begin{codeblock}
lakehouse/.venv/bin/python - <<'EOF'
import duckdb
D = 'lakehouse/source_sim/data/'
con = duckdb.connect()
con.sql(f"CREATE VIEW orders AS SELECT * FROM '{D}olist_orders_dataset.csv'")
# dim_date: khoá yyyymmdd, thêm dòng khoá 0 cho ngày chưa biết
con.sql("""CREATE TABLE dim_date AS
  SELECT CAST(strftime(d, '%Y%m%d') AS INTEGER) AS date_key, CAST(d AS DATE) AS day,
         year(d) AS year, month(d) AS month, dayname(d) AS weekday
  FROM range(DATE '2016-01-01', DATE '2019-01-01', INTERVAL 1 DAY) r(d)
  UNION ALL SELECT 0, NULL, NULL, NULL, 'unknown'""")
# Mỗi vai là một view riêng, tên cột riêng
con.sql("""CREATE VIEW dim_purchase_date AS SELECT date_key AS purchase_date_key,
             year AS purchase_year, month AS purchase_month FROM dim_date""")
con.sql("""CREATE VIEW dim_delivered_date AS SELECT date_key AS delivered_date_key,
             year AS delivered_year, month AS delivered_month FROM dim_date""")
con.sql("""CREATE TABLE fct_orders AS
  SELECT order_id,
         CAST(strftime(order_purchase_timestamp, '%Y%m%d') AS INTEGER)
           AS purchase_date_key,
         coalesce(CAST(strftime(order_delivered_customer_date, '%Y%m%d')
                       AS INTEGER), 0) AS delivered_date_key
  FROM orders""")
print('dim_date:', con.sql("SELECT count(*) FROM dim_date").fetchall()[0][0], 'dòng')
# Đơn đặt trong tháng 12/2017 được giao vào tháng nào?
con.sql("""SELECT d.delivered_year, d.delivered_month, count(*) AS so_don
           FROM fct_orders f
           JOIN dim_purchase_date p USING (purchase_date_key)
           JOIN dim_delivered_date d USING (delivered_date_key)
           WHERE p.purchase_year = 2017 AND p.purchase_month = 12
           GROUP BY ALL ORDER BY ALL""").show()
EOF
\end{codeblock}
\textit{Kết quả:}
% kết quả: G
\begin{codeblock}
dim_date: 1097 dòng
┌────────────────┬─────────────────┬────────┐
│ delivered_year │ delivered_month │ so_don │
│     int64      │      int64      │ int64  │
├────────────────┼─────────────────┼────────┤
│           2017 │              12 │   3332 │
│           2018 │               1 │   2112 │
│           2018 │               2 │     58 │
│           2018 │               3 │      8 │
│           2018 │               4 │      3 │
│           NULL │            NULL │    160 │
└────────────────┴─────────────────┴────────┘
\end{codeblock}

\code{dim\_\allowbreak{}date} có 1.097 dòng: 1.096 ngày từ đầu 2016 tới hết 2018, cộng dòng khoá 0. Phần lớn đơn đặt tháng 12/2017 được giao trong tháng 12 và tháng 1. 160 đơn chưa giao nhận khoá 0, nên hiện NULL ở cột năm và tháng thay vì biến mất khỏi kết quả.

Khối thứ hai đo degenerate dimension và dựng junk dimension cho các cờ của đơn.

% chạy: G2
\begin{codeblock}
lakehouse/.venv/bin/python - <<'EOF'
import duckdb
D = 'lakehouse/source_sim/data/'
con = duckdb.connect()
for t in ['orders', 'order_items', 'order_payments', 'order_reviews']:
    con.sql(f"CREATE VIEW {t} AS SELECT * FROM '{D}olist_{t}_dataset.csv'")
# Degenerate dimension: order_id nằm trên fact, không có bảng dim_order
n_rows, n_orders = con.sql("SELECT count(*), count(DISTINCT order_id) "
                           "FROM order_items").fetchall()[0]
print('dòng hàng:', n_rows, '| đơn khác nhau:', n_orders)
# Junk dimension: gom các cờ nhỏ của đơn, chỉ giữ tổ hợp thật sự xuất hiện
con.sql("""CREATE TABLE flags AS
  SELECT o.order_status,
         coalesce(p.pay_types, 'none') AS pay_types,
         CASE WHEN o.order_delivered_customer_date IS NULL THEN 'chưa giao'
              WHEN CAST(o.order_delivered_customer_date AS DATE)
                   <= CAST(o.order_estimated_delivery_date AS DATE) THEN 'đúng hẹn'
              ELSE 'trễ' END AS delivery,
         r.order_id IS NOT NULL AS has_review
  FROM orders o
  LEFT JOIN (SELECT order_id, string_agg(DISTINCT payment_type, '+'
                    ORDER BY payment_type) AS pay_types
             FROM order_payments GROUP BY 1) p USING (order_id)
  LEFT JOIN (SELECT DISTINCT order_id FROM order_reviews) r USING (order_id)""")
con.sql("""SELECT count(DISTINCT order_status) AS n_status,
                  count(DISTINCT pay_types)    AS n_pay,
                  count(DISTINCT delivery)     AS n_delivery,
                  count(DISTINCT has_review)   AS n_review,
                  n_status * n_pay * n_delivery * n_review AS co_the,
                  (SELECT count(*) FROM (SELECT DISTINCT * FROM flags)) AS xuat_hien
           FROM flags""").show()
con.sql("""CREATE TABLE dim_order_profile AS
  SELECT row_number() OVER (ORDER BY order_status, pay_types, delivery, has_review)
           AS profile_key, *
  FROM (SELECT DISTINCT * FROM flags)""")
con.sql("""SELECT p.*, count(*) AS so_don FROM flags f
           JOIN dim_order_profile p
             USING (order_status, pay_types, delivery, has_review)
           GROUP BY ALL ORDER BY so_don DESC LIMIT 3""").show()
EOF
\end{codeblock}
\textit{Kết quả:}
% kết quả: G2
\begin{codeblock}
dòng hàng: 112650 | đơn khác nhau: 98666
┌──────────┬───────┬────────────┬──────────┬────────┬───────────┐
│ n_status │ n_pay │ n_delivery │ n_review │ co_the │ xuat_hien │
│  int64   │ int64 │   int64    │  int64   │ int64  │   int64   │
├──────────┼───────┼────────────┼──────────┼────────┼───────────┤
│        8 │     8 │          3 │        2 │    384 │        70 │
└──────────┴───────┴────────────┴──────────┴────────┴───────────┘

┌─────────────┬──────────────┬─────────────┬──────────┬────────────┬────────┐
│ profile_key │ order_status │  pay_types  │ delivery │ has_review │ so_don │
│    int64    │   varchar    │   varchar   │ varchar  │  boolean   │ int64  │
├─────────────┼──────────────┼─────────────┼──────────┼────────────┼────────┤
│          24 │ delivered    │ credit_card │ đúng hẹn │ true       │  66922 │
│          19 │ delivered    │ boleto      │ đúng hẹn │ true       │  17690 │
│          22 │ delivered    │ credit_card │ trễ      │ true       │   4714 │
└─────────────┴──────────────┴─────────────┴──────────┴────────────┴────────┘
\end{codeblock}

\begin{itemize}
  \item 112.650 dòng hàng thuộc 98.666 đơn. \code{order\_\allowbreak{}id} nhóm các dòng hàng của cùng một đơn mà không cần bảng riêng.
  \item Các cờ có 8 trạng thái, 8 tổ hợp kiểu thanh toán, 3 trạng thái giao và 2 giá trị có review, tức 384 tổ hợp có thể. Chỉ 70 tổ hợp thật sự xuất hiện trên 99.441 đơn, và ba tổ hợp lớn nhất đã chiếm 89.326 đơn.
  \item Kiểu thanh toán được gộp thành một chuỗi như \code{credit\_\allowbreak{}card+voucher}, vì 2.246 đơn dùng hơn một kiểu. Nếu không gộp, cờ này không có một giá trị duy nhất cho mỗi đơn.
  % Kiểm: count(*) FROM (SELECT order_id FROM order_payments GROUP BY 1 HAVING count(DISTINCT payment_type) > 1) = 2246.
\end{itemize}

\begin{luuy}
\textbf{Bẫy thường gặp.}
\begin{itemize}
  \item \textbf{Lọc theo khoá \code{yyyymmdd} thay vì theo cột năm, tháng.} Kimball khuyên lọc và nhóm theo thuộc tính của dimension ngày, không theo khoá.
  \item \textbf{Mỗi cờ một khoá ngoại.} Fact phình ra thành một bảng có cả chục khoá, mỗi khoá trỏ tới một dimension chỉ vài dòng.
  \item \textbf{Quên dòng khoá 0.} Đơn chưa giao rơi khỏi mọi báo cáo theo ngày giao mà không ai biết.
  \item \textbf{Hai vai dùng chung tên cột.} Nếu cả hai view đều có cột \code{year}, câu truy vấn và biểu đồ sẽ không cho biết đó là năm đặt hay năm giao.
\end{itemize}
\end{luuy}

% =====================================================================
\section{Bài 7: Conformed dimension, bus matrix và drill-across}

\textbf{Vấn đề.} Giám đốc muốn một bảng: doanh thu và điểm review trung bình theo bang của khách. Doanh thu nằm ở \code{fct\_\allowbreak{}order\_\allowbreak{}items}, điểm nằm ở \code{fct\_\allowbreak{}reviews}, và hai fact khác grain: một đơn có nhiều dòng hàng, và có đơn có nhiều review. Join thẳng hai fact rồi gộp là gặp lại fan-out, lần này ở cả tổng lẫn trung bình.

\textbf{Giải thích.}
% Nguồn: Kimball techniques PDF, trang 10 và 11
\begin{itemize}
  \item \textbf{Conformed dimension}: một dimension được nhiều fact dùng chung, với cùng tên cột và cùng miền giá trị. Nhờ đó ``bang SP'' có cùng một nghĩa ở mọi fact.
  \item \textbf{Drill-across}: truy vấn riêng từng fact, gộp theo cùng một thuộc tính của conformed dimension, rồi ghép các kết quả theo thuộc tính đó. Hai fact không bao giờ được join với nhau ở mức dòng.
  \item \textbf{Bus matrix}: bảng có dòng là quy trình nghiệp vụ, cột là dimension, và đánh dấu ô nào áp dụng. Đọc theo dòng để kiểm một dimension có đúng một giá trị cho mỗi sự kiện không. Đọc theo cột để thấy dimension nào phải conform giữa nhiều quy trình. Kimball khuyên xây từng dòng một.
\end{itemize}

Bus matrix của Olist trong module này:

\begin{xltabular}{\textwidth}{@{}P{2.3cm} L P{1.9cm} P{1.9cm} P{1.9cm} P{2.3cm}@{}}
\toprule
\thd{quy trình} & \thd{\code{dim\_\allowbreak{}date}} & \thd{\code{dim\_\allowbreak{}customer}} & \thd{\code{dim\_\allowbreak{}product}} & \thd{\code{dim\_\allowbreak{}seller}} & \thd{\code{dim\_\allowbreak{}order\_\allowbreak{}profile}} \\
\midrule
\endhead
\bottomrule
\endfoot
đặt hàng & ngày đặt, ngày giao & x & x & x & x \\
thanh toán & ngày đặt & x & & & \\
giao hàng & năm mốc của đơn & x & & & x \\
đánh giá & ngày tạo review & x & & & \\
\end{xltabular}

Ô người bán của dòng giao hàng bỏ trống có lý do: ở grain một dòng mỗi đơn, 1.278 đơn có từ hai người bán trở lên, nên người bán không có một giá trị duy nhất. Muốn phân tích giao hàng theo người bán, grain phải là một dòng mỗi đơn và mỗi người bán.
% Kiểm: count(*) FROM (SELECT order_id FROM order_items GROUP BY 1 HAVING count(DISTINCT seller_id) > 1) = 1278.

Một hệ OLTP thường sinh ra nhiều star, mỗi quy trình một star, dùng chung các conformed dimension; tập các star đó đôi khi được gọi là \textit{fact constellation}.
% Nguồn: https://en.wikipedia.org/wiki/Fact_constellation
Chiều ngược lại cũng đúng: một star có thể lấy dữ liệu từ nhiều service, như thanh toán từ service thanh toán và dòng hàng từ service đơn hàng.

\begin{ghichu}
\textbf{So với SWE.} Conformed dimension giống một enum hay một kiểu dùng chung mà mọi service cùng import, như một file protobuf định nghĩa mã bang. Nếu mỗi service tự đặt mã bang của mình, không ai ghép được số liệu giữa các service. Bus matrix giống ma trận service và capability mà team kiến trúc hay vẽ: nhìn một bảng là biết ai dùng gì, và đổi một cột thì ảnh hưởng tới những dòng nào.

\textbf{Chỗ so sánh gãy.} Enum dùng chung là hằng số trong code, đổi bằng một lần deploy. Conformed dimension là một bảng có dữ liệu và lịch sử, có pipeline và người chịu trách nhiệm. Đổi nó có thể làm đổi số liệu quá khứ của mọi fact dùng nó.

Bus matrix còn là kế hoạch xây: Kimball dùng nó để xếp thứ tự ưu tiên với phía nghiệp vụ, và mỗi lần chỉ làm một dòng, tức một quy trình.
% Nguồn: Kimball techniques PDF, trang 11, mục Enterprise Data Warehouse Bus Matrix
\end{ghichu}

\subsection{Làm thử}

Trong khối dưới, \code{dim\_\allowbreak{}customer} tạm dùng bảng \code{customers} của Olist, mỗi dòng là địa chỉ của khách tại một đơn. Bài 9 dựng bản đúng bằng SCD Type 2 và cho ra cùng các con số này.

% chạy: H
\begin{codeblock}
lakehouse/.venv/bin/python - <<'EOF'
import duckdb
D = 'lakehouse/source_sim/data/'
con = duckdb.connect()
for t in ['orders', 'order_items', 'order_reviews', 'customers']:
    con.sql(f"CREATE VIEW {t} AS SELECT * FROM '{D}olist_{t}_dataset.csv'")
# Hai fact khác grain, cùng dùng một dimension khách hàng
con.sql("""CREATE TABLE fct_order_items AS
  SELECT i.order_id, o.customer_id, CAST(i.price AS DECIMAL(12,2)) AS price
  FROM order_items i JOIN orders o USING (order_id)""")
con.sql("""CREATE TABLE fct_reviews AS
  SELECT r.review_id, r.order_id, o.customer_id, r.review_score
  FROM order_reviews r JOIN orders o USING (order_id)""")
con.sql("""CREATE TABLE dim_customer AS
  SELECT customer_id, customer_state FROM customers""")
print('Sai: join hai fact với nhau rồi mới gộp')
con.sql("""SELECT c.customer_state, sum(f.price) AS revenue,
                  round(avg(r.review_score), 2) AS avg_review
           FROM fct_order_items f JOIN dim_customer c USING (customer_id)
           LEFT JOIN fct_reviews r USING (order_id)
           GROUP BY 1 ORDER BY revenue DESC LIMIT 3""").show()
print('Đúng: drill-across, mỗi fact gộp riêng theo customer_state rồi ghép')
con.sql("""WITH rev AS (
             SELECT c.customer_state, sum(f.price) AS revenue
             FROM fct_order_items f JOIN dim_customer c USING (customer_id)
             GROUP BY 1),
           rv AS (
             SELECT c.customer_state, round(avg(r.review_score), 2) AS avg_review
             FROM fct_reviews r JOIN dim_customer c USING (customer_id)
             GROUP BY 1)
           SELECT * FROM rev JOIN rv USING (customer_state)
           ORDER BY revenue DESC LIMIT 3""").show()
EOF
\end{codeblock}
\textit{Kết quả:}
% kết quả: H
\begin{codeblock}
Sai: join hai fact với nhau rồi mới gộp
┌────────────────┬───────────────┬────────────┐
│ customer_state │    revenue    │ avg_review │
│    varchar     │ decimal(38,2) │   double   │
├────────────────┼───────────────┼────────────┤
│ SP             │    5228869.17 │       4.13 │
│ RJ             │    1831678.85 │       3.81 │
│ MG             │    1591518.47 │       4.09 │
└────────────────┴───────────────┴────────────┘

Đúng: drill-across, mỗi fact gộp riêng theo customer_state rồi ghép
┌────────────────┬───────────────┬────────────┐
│ customer_state │    revenue    │ avg_review │
│    varchar     │ decimal(38,2) │   double   │
├────────────────┼───────────────┼────────────┤
│ SP             │    5202955.05 │       4.17 │
│ RJ             │    1824092.67 │       3.87 │
│ MG             │    1585308.03 │       4.14 │
└────────────────┴───────────────┴────────────┘
\end{codeblock}

\begin{itemize}
  \item Join thẳng hai fact làm doanh thu của SP thành 5.228.869,17 BRL thay vì 5.202.955,05 BRL, vì đơn có nhiều review làm dòng hàng bị lặp. Điểm trung bình thành 4,13 thay vì 4,17, vì mỗi review bị đếm một lần cho mỗi dòng hàng của đơn.
  \item Drill-across cho số đúng: SP được 5.202.955,05 BRL và 4,17 điểm, RJ được 1.824.092,67 BRL và 3,87 điểm, MG được 1.585.308,03 BRL và 4,14 điểm.
\end{itemize}

\begin{luuy}
\textbf{Bẫy thường gặp.}
\begin{itemize}
  \item \textbf{Join hai fact ở mức dòng.} Luôn gộp riêng từng fact về thuộc tính chung trước.
  \item \textbf{Ghép theo hai vai khác nhau của một thuộc tính.} Bang của khách và bang của người bán đều là ``bang'', nhưng ghép doanh thu theo bang người bán với điểm review theo bang khách là sai nghĩa.
  \item \textbf{Quên rằng trung bình phụ thuộc grain.} Bang AM được 4,18 điểm khi lấy trung bình 147 dòng review, nhưng được 4,21 điểm khi mỗi đơn chỉ giữ điểm của review mới nhất, trên 146 đơn. Ghi rõ định nghĩa ngay cạnh con số.
  % Kiểm: AM, 615 điểm trên 147 dòng review = 4,1837; review mới nhất theo thứ tự ở khối E, 614 điểm trên 146 đơn = 4,2055. AM lệch nhiều nhất trong 27 bang; RJ chỉ lệch từ 3,8750 lên 3,8771.
\end{itemize}
\end{luuy}

% =====================================================================
\section{Bài 8: Star, snowflake hay One Big Table}

\textbf{Vấn đề.} Phản xạ SWE của bạn nhìn \code{dim\_\allowbreak{}product} và thấy tên danh mục bị lặp trên 32.951 dòng, trong khi chỉ có 74 danh mục. Bạn muốn chuẩn hoá: tách ra \code{dim\_\allowbreak{}category} và để \code{dim\_\allowbreak{}product} chỉ giữ khoá. Một đồng nghiệp lại muốn đi hướng ngược lại: chép mọi thuộc tính của dimension lên từng dòng fact, để BA khỏi phải join. Ai đúng?

\textbf{Giải thích.} Tách một dimension thành các bảng con như vậy gọi là \textit{snowflake}. Kimball khuyên tránh snowflake: nó khó đọc với người dùng nghiệp vụ, có thể làm truy vấn chậm hơn, và một dimension phẳng chứa đúng cùng thông tin.
% Nguồn: Kimball techniques PDF, trang 9

Còn lý do tiết kiệm chỗ thì sao? Trên file dạng cột như Parquet, một cột có ít giá trị khác nhau gần như không tốn chỗ, nhờ \textit{dictionary encoding}. Mỗi giá trị khác nhau được lưu một lần trong một trang từ điển của cột, còn mỗi dòng chỉ lưu một số nguyên nhỏ trỏ vào từ điển. Module 4 sẽ giải thích kỹ cách mã hoá này.
% Nguồn: https://parquet.apache.org/docs/file-format/data-pages/encodings/

Hướng ngược lại có tên là \textit{One Big Table}, viết tắt OBT: một bảng rộng ở grain của fact, mang theo mọi thuộc tính của các dimension. Năm 2022, Michael Kaminsky so hai cách trên Redshift, Snowflake và BigQuery với dữ liệu TPC-DS. Các câu kiểu dashboard trên OBT chạy nhanh hơn star khoảng 25 tới 50\%, nhưng trên Redshift OBT chiếm khoảng gấp đôi chỗ của bảng fact. Ông gợi ý dựng star trước, rồi join lại thành OBT cho người dùng cuối.
% Nguồn: https://www.fivetran.com/blog/star-schema-vs-obt (bài viết của Michael Kaminsky, 16/08/2022): "25%-50% faster depending on the warehouse"; one_big_table khoảng 60.250 MB so với store_sales 29.778 MB; "staging your ELT process such that the data all get transformed into something like a star schema before everything gets re-joined back together for end-user querying".

\begin{ghichu}
\textbf{So với SWE.} Dictionary encoding giống string interning, hay giống việc bạn lưu một enum thay vì chuỗi: mỗi giá trị có một bản, phần còn lại chỉ là tham chiếu. Nói cách khác, Parquet đã tự ``chuẩn hoá'' cột đó ở tầng lưu trữ, nên bạn không cần chuẩn hoá ở tầng mô hình.

OBT giống một document nhúng sẵn dữ liệu liên quan. Tài liệu MongoDB khuyên nhúng khi hai thực thể có quan hệ ``chứa'', như document \code{contacts} chứa \code{address}, để đọc mọi thứ trong một lần truy cập.
% Nguồn: https://www.mongodb.com/docs/manual/data-modeling/embedding/

\textbf{Chỗ so sánh gãy.} Khác string interning, từ điển của Parquet chỉ nằm trong từng column chunk, và bị bỏ khi cột có quá nhiều giá trị khác nhau, như cột \code{product\_\allowbreak{}id} ở khối dưới.

Lý do chuẩn hoá trong OLTP không phải là chỗ lưu, mà là tránh sửa sót: sửa tên danh mục ở một bảng rồi quên bảng khác. Dimension do code của pipeline cập nhật, không ai sửa tay, nên lý do đó cũng không còn.

OBT thì làm lý do đó quay lại. Một dòng hàng không ``chứa'' danh mục của sản phẩm; OBT chỉ chép tên danh mục lên mọi dòng hàng của nó, nên đổi tên một danh mục là phải ghi lại hàng nghìn dòng, như khối thứ hai cho thấy.
\end{ghichu}

\subsection{Làm thử}

Khối đầu ghi bản star và bản snowflake ra Parquet trong một thư mục tạm, với nén mặc định của DuckDB, rồi so kích thước.

% chạy: I
\begin{codeblock}
lakehouse/.venv/bin/python - <<'EOF'
import duckdb, os, tempfile
D = 'lakehouse/source_sim/data/'
con = duckdb.connect()
con.sql(f"CREATE VIEW products AS SELECT * FROM '{D}olist_products_dataset.csv'")
con.sql(f"CREATE VIEW category_translation AS "
        f"SELECT * FROM '{D}product_category_name_translation.csv'")
# Star: tên danh mục nằm thẳng trong dim_product
con.sql("""CREATE TABLE dim_product AS
  SELECT p.product_id, coalesce(t.product_category_name_english,
                                p.product_category_name, 'unknown') AS category
  FROM products p LEFT JOIN category_translation t USING (product_category_name)""")
# Snowflake: tách danh mục ra bảng riêng, dim_product chỉ giữ khoá
con.sql("""CREATE TABLE dim_category AS
  SELECT row_number() OVER (ORDER BY category) AS category_id, category
  FROM (SELECT DISTINCT category FROM dim_product)""")
con.sql("""CREATE TABLE dim_product_sf AS
  SELECT p.product_id, c.category_id
  FROM dim_product p JOIN dim_category c USING (category)""")
with tempfile.TemporaryDirectory() as tmp:
    for t in ['dim_product', 'dim_product_sf', 'dim_category']:
        f = f'{tmp}/{t}.parquet'
        con.sql(f"COPY {t} TO '{f}' (FORMAT parquet)")
        n = con.sql(f"SELECT count(*) FROM {t}").fetchall()[0][0]
        print(f'{t:15} {n:6} dòng {os.path.getsize(f):10,} byte')
    con.sql(f"""SELECT path_in_schema AS cot, encodings,
                       total_compressed_size AS byte
                FROM parquet_metadata('{tmp}/dim_product.parquet')""").show()
EOF
\end{codeblock}
\textit{Kết quả:}
% kết quả: I
\begin{codeblock}
dim_product      32951 dòng  1,092,443 byte
dim_product_sf   32951 dòng  1,091,566 byte
dim_category        74 dòng      1,820 byte
┌────────────┬──────────────────┬─────────┐
│    cot     │    encodings     │  byte   │
│  varchar   │     varchar      │  int64  │
├────────────┼──────────────────┼─────────┤
│ product_id │ PLAIN            │ 1061589 │
│ category   │ PLAIN_DICTIONARY │   30231 │
└────────────┴──────────────────┴─────────┘
\end{codeblock}

\begin{itemize}
  \item Bản star: \code{dim\_\allowbreak{}product} nặng 1.092.443 byte.
  \item Bản snowflake: \code{dim\_\allowbreak{}product\_\allowbreak{}sf} nặng 1.091.566 byte, cộng \code{dim\_\allowbreak{}category} nặng 1.820 byte, tổng 1.093.386 byte. Bản snowflake nặng hơn 943 byte.
  \item Cột \code{category} chỉ chiếm 30.231 byte, khoảng 2,8\% file, vì nó được ghi bằng dictionary encoding. DuckDB in ra \code{PLAIN\_\allowbreak{}DICTIONARY}, giá trị cũ mà spec Parquet đã đánh dấu deprecated; Module 4 dùng giá trị hiện tại là \code{RLE\_\allowbreak{}DICTIONARY}. Phần lớn file là cột \code{product\_\allowbreak{}id}, 1.061.589 byte.
\end{itemize}

Snowflake không tiết kiệm được gì ở đây, lại thêm một join cho mọi câu hỏi theo danh mục.

Khối thứ hai dựng OBT cho dòng hàng, rồi đo hai cái giá của nó: chỗ lưu, và số dòng phải sửa khi một dimension đổi.

% chạy: I2
\begin{codeblock}
lakehouse/.venv/bin/python - <<'EOF'
import duckdb, tempfile
D = 'lakehouse/source_sim/data/'
con = duckdb.connect()
for t in ['orders', 'order_items', 'customers', 'products', 'sellers']:
    con.sql(f"CREATE VIEW {t} AS SELECT * FROM '{D}olist_{t}_dataset.csv'")
con.sql(f"CREATE VIEW category_translation AS "
        f"SELECT * FROM '{D}product_category_name_translation.csv'")
con.sql("""CREATE TABLE dim_product AS
  SELECT p.product_id, coalesce(t.product_category_name_english,
                                p.product_category_name, 'unknown') AS category
  FROM products p LEFT JOIN category_translation t USING (product_category_name)""")
# OBT: thuộc tính của dimension chép lên từng dòng hàng
con.sql("""CREATE TABLE obt_order_items AS
  SELECT i.order_id, i.order_item_id, i.product_id, i.seller_id, o.customer_id,
         CAST(o.order_purchase_timestamp AS DATE) AS purchase_date,
         CAST(i.price AS DECIMAL(12,2)) AS price, p.category,
         s.seller_city, s.seller_state, c.customer_city, c.customer_state
  FROM order_items i JOIN orders o USING (order_id)
  JOIN dim_product p USING (product_id) JOIN sellers s USING (seller_id)
  JOIN customers c USING (customer_id) ORDER BY i.order_id, i.order_item_id""")
with tempfile.TemporaryDirectory() as tmp:
    con.sql(f"COPY obt_order_items TO '{tmp}/obt.parquet' (FORMAT parquet)")
    total, copied = con.sql(f"""SELECT sum(total_compressed_size),
        sum(total_compressed_size) FILTER (WHERE path_in_schema IN ('category',
          'seller_city', 'seller_state', 'customer_city', 'customer_state'))
        FROM parquet_metadata('{tmp}/obt.parquet')""").fetchall()[0]
print(f'OBT trên Parquet: {total:,} byte; 5 cột chép từ dimension: '
      f'{copied:,} byte ({copied / total:.1%})')
# Đổi tên một danh mục: star sửa dimension, OBT sửa mọi dòng hàng của danh mục
print('dòng phải sửa khi đổi tên bed_bath_table, ở star và ở OBT:', con.sql("""
  SELECT (SELECT count(*) FROM dim_product WHERE category = 'bed_bath_table'),
         (SELECT count(*) FROM obt_order_items
          WHERE category = 'bed_bath_table')""").fetchall()[0])
EOF
\end{codeblock}
\textit{Kết quả:}
% kết quả: I2
\begin{codeblock}
OBT trên Parquet: 10,467,714 byte; 5 cột chép từ dimension: 585,411 byte (5.6%)
dòng phải sửa khi đổi tên bed_bath_table, ở star và ở OBT: (3029, 11115)
\end{codeblock}

\begin{itemize}
  \item Năm cột chép từ dimension chỉ tốn 585.411 byte, khoảng 5,6\% file, cũng nhờ dictionary encoding. Trên Parquet, chép chữ lặp lại gần như không tốn chỗ.
  \item Kaminsky dựng bảng bằng một công cụ không đặt được kiểu mã hoá cột trên Redshift, và chính ông ghi rằng số chỗ lưu đo được có lẽ là trường hợp xấu nhất.
  % Nguồn: https://www.fivetran.com/blog/star-schema-vs-obt: "Because dbt doesn't have the ability to specify column compression or encoding style in Redshift, this is probably the worst-possible-case in terms of disk storage size."
  \item Cái giá thật nằm ở cập nhật. Đổi tên danh mục \code{bed\_\allowbreak{}bath\_\allowbreak{}table} là sửa 3.029 dòng của \code{dim\_\allowbreak{}product}, hoặc 1 dòng của \code{dim\_\allowbreak{}category} nếu dùng snowflake, nhưng là ghi lại 11.115 dòng của OBT.
\end{itemize}

Câu hỏi theo danh mục và bang trên OBT không cần join nào, vì mọi cột đã nằm trên cùng một dòng. Kết luận cho Olist: giữ star làm mô hình chính; nếu một dashboard cần OBT, dựng nó từ star như một mart, và dựng lại nó mỗi khi dimension đổi.

\begin{luuy}
\textbf{Bẫy thường gặp.}
\begin{itemize}
  \item \textbf{Tối ưu trước khi đo.} Khối trên chạy chưa tới một giây trên máy này. Hãy đo trước khi tách bảng.
  \item \textbf{Snowflake vì nguồn đã chuẩn hoá.} Nguồn chuẩn hoá vì nó phục vụ việc ghi; star schema phục vụ việc đọc.
  \item \textbf{Đem kết luận này sang dimension rất lớn mà không đo.} Dimension lớn nhất ở đây chỉ có 99.441 dòng. Với dimension lớn hơn nhiều, hãy đo như trên rồi mới quyết định có tách hay không.
  \item \textbf{Nhét cột của grain khác vào OBT.} Thêm \code{payment\_\allowbreak{}value} vào OBT ở grain dòng hàng là lặp lại đúng lỗi ở phần mở đầu: tổng tiền tăng 26,9\%.
  \item \textbf{Coi OBT là nơi giữ dữ liệu gốc.} OBT là bản sao để đọc. Sửa trên OBT mà không sửa star thì lần dựng lại sau sẽ xoá mất chỗ sửa đó.
\end{itemize}
\end{luuy}

% =====================================================================
\section{Bài 9: Slowly changing dimension}

\textbf{Vấn đề.} Người mua \code{37266d5c6bf5cb0700372d16599a9b21} đặt đơn đầu tiên vào ngày 2017-10-08 ở timoteo, zip 35180, rồi đặt đơn thứ hai ngày 2018-06-17 ở coronel fabriciano, zip 35170. Nếu \code{dim\_\allowbreak{}customer} chỉ giữ địa chỉ mới nhất, như \repo{core.customers} của repo \repo{lakehouse} hiện nay, đơn năm 2017 bị tính cho coronel fabriciano. Trên cả Olist, 250 người mua có hơn một zip.
% Nguồn: lakehouse/docs/data-model.md, mục 6, core.customers lấy địa chỉ ở đơn mới nhất.

\textbf{Giải thích.} \textit{Slowly changing dimension}, viết tắt SCD, là cách một dimension xử lý khi thuộc tính của một thực thể đổi theo thời gian. Kimball đánh số các cách làm từ Type 0 tới Type 7. Module này dùng năm cách dưới; Type 4, 5 và 7 nằm ngoài phạm vi module.
% Nguồn: Kimball techniques PDF, trang 12 và 13

\begin{xltabular}{\textwidth}{@{}P{1.6cm} L P{4.6cm}@{}}
\toprule
\thd{kiểu} & \thd{làm gì} & \thd{đơn 2017 của người mua trên được tính vào} \\
\midrule
\endhead
\bottomrule
\endfoot
Type 0 & giữ giá trị gốc, không bao giờ đổi & timoteo \\
Type 1 & ghi đè bằng giá trị mới, mất lịch sử & coronel fabriciano \\
Type 2 & thêm một dòng cho mỗi version, có ngày bắt đầu hiệu lực, ngày hết hiệu lực và cờ dòng hiện hành & timoteo \\
Type 3 & ghi đè cột chính, thêm một cột giữ giá trị ngay trước đó & coronel fabriciano theo cột chính, timoteo theo cột giá trị trước; chỉ nhớ một lần đổi \\
Type 6 & Type 2, cộng thêm các cột giá trị hiện hành được ghi đè trên mọi version & cả hai: timoteo theo version lúc đặt, coronel fabriciano theo giá trị hiện tại \\
\end{xltabular}

Type 2 cần khoá thay thế cho mỗi version, ví dụ \code{customer\_\allowbreak{}sk}, vì \code{customer\_\allowbreak{}unique\_\allowbreak{}id} giờ lặp lại trên nhiều dòng. Fact trỏ tới version còn hiệu lực lúc sự kiện xảy ra.

Khoảng hiệu lực là khoảng nửa mở \code{[valid\_\allowbreak{}from, valid\_\allowbreak{}to)}: tính cả điểm đầu, không tính điểm cuối. Hai version liền nhau dùng chung một mốc ranh giới nhưng không bao giờ chồng nhau, nên mỗi thời điểm thuộc đúng một version. Dòng hiện hành để \code{valid\_\allowbreak{}to} là NULL, hoặc là một ngày rất xa như 9999-12-31; Microsoft nhắc tới cả hai cách.
% Nguồn: https://learn.microsoft.com/en-us/power-bi/guidance/star-schema (mục Type 2 SCD)

Một trường hợp liên quan là \textit{late arriving dimension}: fact tới trước dimension của nó. Kimball tạo một dòng dimension giữ chỗ chứa khoá tự nhiên và các thuộc tính ``chưa biết'', rồi ghi đè kiểu Type 1 khi thông tin tới. Repo \repo{lakehouse} chọn cách khác: \repo{core.orders} để \code{customer\_\allowbreak{}key} là NULL khi customer chưa tới, rồi tính lại đơn khi customer tới.
% Nguồn: Kimball techniques PDF, trang 20; lakehouse/docs/data-model.md, mục 6, core.orders.

Lịch sử Type 2 và các data test ở Bài 10 cũng viết được bằng dbt, công cụ mà nhánh tuỳ chọn T1 dạy; module này tự viết bằng SQL để bạn thấy từng bước.
% Nguồn: https://docs.getdbt.com/docs/build/snapshots; https://docs.getdbt.com/docs/build/data-tests

\begin{ghichu}
\textbf{So với SWE.} Type 1 là câu \code{UPDATE customers SET city = ...} mà app của bạn vẫn chạy mỗi ngày.

Type 2 giống S3 Versioning. Bucket bật versioning giữ một version hiện hành cùng không hoặc nhiều version cũ của mỗi object; mỗi version có version ID riêng, còn key của object giữ nguyên. Ở đây key là khoá bền \code{customer\_\allowbreak{}unique\_\allowbreak{}id}, còn version ID là khoá thay thế \code{customer\_\allowbreak{}sk}.
% Nguồn: https://docs.aws.amazon.com/AmazonS3/latest/userguide/Versioning.html

Khoảng nửa mở giống lịch đặt phòng: hai lượt \code{[14:30, 15:30)} và \code{[15:30, 16:30)} nối nhau mà không chồng. Hàm tạo range của PostgreSQL mặc định dùng đúng quy ước này. PostgreSQL 18 còn có ràng buộc \code{UNIQUE (id, valid\_\allowbreak{}at WITHOUT OVERLAPS)}, đúng luật ``một thực thể không có hai version chồng thời gian'' của Type 2. Cột \code{valid\_\allowbreak{}at} phải có kiểu range như \code{tstzrange}, và cột \code{id}, nếu là số hay chuỗi, cần extension \code{btree\_\allowbreak{}gist}.
% Nguồn: https://www.postgresql.org/docs/18/rangetypes.html; https://www.postgresql.org/docs/18/sql-createtable.html

\textbf{Chỗ so sánh gãy.} S3 tạo một version cho mỗi lệnh ghi. Type 2 chỉ thấy các thay đổi mà pipeline kịp nhìn thấy giữa hai lần nạp: nếu địa chỉ đổi hai lần trong một ngày mà pipeline chạy mỗi ngày một lần, version ở giữa mất vĩnh viễn. Fact cũng phải giữ đúng khoá của version, hoặc tra theo thời gian lúc nạp; S3 không có khái niệm này.
\end{ghichu}

\subsection{Làm thử}

Bảng \code{customers} của Olist có một dòng cho mỗi đơn, nên nó chụp lại địa chỉ của người mua tại lúc đặt. Khối dưới dùng các ảnh chụp đó để dựng Type 1 và Type 2, rồi đếm sai số. Thay đổi được theo dõi theo zip.

% chạy: J
\begin{codeblock}
lakehouse/.venv/bin/python - <<'EOF'
import duckdb
D = 'lakehouse/source_sim/data/'
con = duckdb.connect()
for t in ['orders', 'order_items', 'customers']:
    con.sql(f"CREATE VIEW {t} AS SELECT * FROM '{D}olist_{t}_dataset.csv'")
one = lambda q: con.sql(q).fetchall()[0]
# Mỗi đơn chụp lại địa chỉ của người mua tại lúc đặt
con.sql("""CREATE TABLE addr AS
  SELECT o.order_id, o.order_purchase_timestamp AS ts, c.customer_id,
         c.customer_unique_id AS buyer, c.customer_zip_code_prefix AS zip,
         c.customer_city AS city, c.customer_state AS state
  FROM orders o JOIN customers c USING (customer_id)""")

# Type 1: chỉ giữ địa chỉ mới nhất, giống một câu UPDATE ghi đè
con.sql("""CREATE TABLE dim_customer_t1 AS
  SELECT buyer, zip, city, state FROM addr
  QUALIFY row_number() OVER (PARTITION BY buyer
                             ORDER BY ts DESC, customer_id DESC) = 1""")
print('Type 1, số đơn sai thành phố, sai bang:', one("""
  SELECT count(*) FILTER (WHERE a.city <> t.city),
         count(*) FILTER (WHERE a.state <> t.state)
  FROM addr a JOIN dim_customer_t1 t USING (buyer)"""))
print('Type 1, tiền hàng bị tính sang bang khác:', one("""
  SELECT sum(CAST(i.price AS DECIMAL(14,2)))
  FROM order_items i JOIN addr a USING (order_id)
  JOIN dim_customer_t1 t USING (buyer)
  WHERE a.state <> t.state""")[0])

# Type 2: mỗi lần zip đổi là một version, hiệu lực trong [valid_from, valid_to)
con.sql("""CREATE TABLE dim_customer_t2 AS
  WITH marked AS (
    SELECT *, zip IS DISTINCT FROM lag(zip) OVER (PARTITION BY buyer ORDER BY ts)
                AS is_new_version
    FROM addr)
  SELECT row_number() OVER (ORDER BY buyer, ts) AS customer_sk,
         buyer, zip, city, state,
         ts AS valid_from,
         lead(ts) OVER (PARTITION BY buyer ORDER BY ts) AS valid_to,
         lead(ts) OVER (PARTITION BY buyer ORDER BY ts) IS NULL AS is_current
  FROM marked WHERE is_new_version""")
print('Type 2, số version, số người mua, số dòng hiện hành:', one("""
  SELECT count(*), count(DISTINCT buyer), count(*) FILTER (WHERE is_current)
  FROM dim_customer_t2"""))

# Mỗi đơn phải khớp đúng một version
match = """SELECT a.order_id, count(t.customer_sk) AS n
           FROM addr a LEFT JOIN dim_customer_t2 t ON t.buyer = a.buyer AND {cond}
           GROUP BY 1"""
half_open = "a.ts >= t.valid_from AND (t.valid_to IS NULL OR a.ts < t.valid_to)"
closed = ("a.ts BETWEEN t.valid_from "
          "AND coalesce(t.valid_to, TIMESTAMP '9999-12-31')")
for name, cond in [('nửa mở [from, to)', half_open),
                   ('đóng, dùng BETWEEN', closed)]:
    print(name, '(số version khớp, số đơn):', con.sql(
        f"SELECT n, count(*) FROM ({match.format(cond=cond)}) GROUP BY 1 ORDER BY 1"
    ).fetchall())

example = con.sql("""
  SELECT customer_sk, zip, city, valid_from, valid_to, is_current
  FROM dim_customer_t2 WHERE buyer = '37266d5c6bf5cb0700372d16599a9b21'""")
for row in example.fetchall():
    print(*row, sep=' | ')
EOF
\end{codeblock}
\textit{Kết quả:}
% kết quả: J
\begin{codeblock}
Type 1, số đơn sai thành phố, sai bang: (130, 43)
Type 1, tiền hàng bị tính sang bang khác: 6164.99
Type 2, số version, số người mua, số dòng hiện hành: (96352, 96096, 96096)
nửa mở [from, to) (số version khớp, số đơn): [(1, 99441)]
đóng, dùng BETWEEN (số version khớp, số đơn): [(1, 99185), (2, 256)]
20825 | 35180 | timoteo | 2017-10-08 21:38:21 | 2018-06-17 23:51:02 | False
20826 | 35170 | coronel fabriciano | 2018-06-17 23:51:02 | None | True
\end{codeblock}

\begin{itemize}
  \item Type 1 tính sai thành phố của 130 đơn và sai bang của 43 đơn; 6.164,99 BRL tiền hàng bị tính sang bang khác. Doanh thu của SP theo Type 1 là 5.205.024,43 BRL thay vì 5.202.955,05 BRL.
  % Kiểm: tổng price theo bang của bản Type 1, so với bang lúc đặt; chạy ngày 05/11/2026.
  \item Type 2 có 96.352 version cho 96.096 người mua, và mỗi người mua có đúng một dòng hiện hành. Với khoảng nửa mở, cả 99.441 đơn khớp đúng một version. Doanh thu theo bang qua bản Type 2 bằng đúng các con số ở Bài 7.
  % Kiểm: bang theo version Type 2 khác bang lúc đặt ở 0 đơn.
  \item Với \code{BETWEEN}, tức khoảng đóng ở cả hai đầu, 256 đơn khớp hai version. Đó là 256 đơn nằm đúng ranh giới, mỗi lần đổi zip một đơn. Fact join theo cách này sẽ bị fan-out như ở Bài 3.
  \item Hai dòng cuối là hai version của người mua ở phần vấn đề, mỗi version một \code{customer\_\allowbreak{}sk}.
\end{itemize}

\begin{luuy}
\textbf{Bẫy thường gặp.}
\begin{itemize}
  \item \textbf{Dùng \code{BETWEEN} cho khoảng hiệu lực.} Dùng \code{>=} ở đầu và \code{<} ở cuối.
  \item \textbf{Quên dòng hiện hành có \code{valid\_\allowbreak{}to} là NULL.} Điều kiện \code{ts < valid\_\allowbreak{}to} loại mọi dòng hiện hành; hãy thêm \code{valid\_\allowbreak{}to IS NULL OR}, hoặc dùng một ngày rất xa.
  \item \textbf{Theo dõi cột nhiễu.} Nếu so cả tên thành phố, bạn được 96.355 version thay vì 96.352, vì 2 người mua có cùng zip nhưng tên thành phố ghi khác nhau, như planaltina và planaltina de goias. Chỉ theo dõi cột mà thay đổi của nó có nghĩa với nghiệp vụ.
  % Kiểm: dựng Type 2 theo (zip, city, state) cho 96.355 version.
  \item \textbf{Dựng Type 2 khi không ai cần lịch sử.} Type 2 làm bảng to hơn và join khó hơn. Hỏi BA trước: báo cáo cần địa chỉ lúc đặt hay địa chỉ hiện tại?
\end{itemize}
\end{luuy}

% =====================================================================
\section{Bài 10: Mô hình trong lakehouse và test dữ liệu}

\textbf{Vấn đề.} Trên PostgreSQL ở Module 1, ràng buộc \code{PRIMARY KEY} và \code{FOREIGN KEY} chặn dòng sai ngay lúc ghi. Trong lakehouse, không ai chặn giúp bạn. Spec của Iceberg ghi rằng tính duy nhất theo khoá định danh không được Iceberg bảo đảm hay đòi hỏi, và việc đó thuộc trách nhiệm của engine hoặc bên cung cấp dữ liệu. Databricks có cú pháp khoá chính và khoá ngoại, nhưng tài liệu của họ viết rằng các ràng buộc này chỉ để thông tin và không được thực thi. Vậy ai giữ cho khoá của star schema đúng?
% Nguồn: https://iceberg.apache.org/spec/#identifier-field-ids; https://docs.databricks.com/aws/en/tables/constraints

\textbf{Giải thích.} Star schema nằm ở tầng gold. Databricks mô tả gold là nơi mô hình hoá dữ liệu bằng dimensional model, thiết lập quan hệ và định nghĩa measure. Star schema được dựng từ tầng silver đã làm sạch, không đọc thẳng từ OLTP.
% Nguồn: https://docs.databricks.com/aws/en/lakehouse/medallion

Trong repo \repo{lakehouse}, tầng silver \repo{core} đã gần như một star: dimension \code{customers}, \code{sellers}, \code{products}, fact \code{orders}, \code{order\_\allowbreak{}items}, và fact không chép cột của dimension. Tầng gold, gồm \code{analytics} và \code{marts}, còn nằm trong kế hoạch. Chi tiết ở mục 6 và mục 11 của \repo{docs/data-model.md}; Lab 3.5 dựng thử \code{marts} trên DuckDB.

Thứ giữ cho khoá đúng là \textit{data test}: một câu \code{SELECT} đi tìm các dòng ``lỗi'', tức các dòng bác bỏ điều bạn khẳng định; trả về 0 dòng nghĩa là đạt. Một star schema cần ít nhất các loại sau:
% Nguồn: https://docs.getdbt.com/docs/build/data-tests

\begin{xltabular}{\textwidth}{@{}P{2.4cm} L L@{}}
\toprule
\thd{loại test} & \thd{khẳng định} & \thd{bắt được lỗi gì} \\
\midrule
\endhead
\bottomrule
\endfoot
grain & khoá của fact và dimension không trùng & nạp một lô hai lần \\
quan hệ & mọi khoá trên fact có dòng ở dimension, không có NULL & dimension thiếu dòng, khoá quên đổi về 0 \\
đối chiếu & số dòng và tổng tiền ở fact bằng ở nguồn & fan-out, join làm rơi dòng \\
SCD Type 2 & mỗi khoá bền có đúng một dòng hiện hành, các version không chồng nhau & lần nạp đóng thiếu hoặc mở thừa một version \\
\end{xltabular}

\begin{ghichu}
\textbf{So với SWE.} Data test giống câu assert trong unit test cộng với checksum sau khi chép file. Test grain và test quan hệ là các assert về cấu trúc. Test đối chiếu là checksum: tổng \code{price} ở fact bằng tổng ở nguồn nghĩa là không dòng nào bị nhân lên hay rơi mất trên đường đi.

\textbf{Chỗ so sánh gãy.} Unit test chạy trên fixture cố định, và đỏ nghĩa là code sai. Data test chạy trên dữ liệu thật ở mỗi lần build, nên đỏ có thể do nguồn đổi chứ không do code của bạn. Bạn phải quyết định dừng pipeline hay chỉ cảnh báo, và đó là chủ đề của Module 10.
\end{ghichu}

\subsection{Làm thử}

Khối dưới chạy ba test trên một fact nhỏ, cố ý làm hỏng dữ liệu, rồi chạy lại. Đây là vòng đỏ và xanh quen thuộc của TDD, chỉ khác là thứ bị kiểm là dữ liệu.

% chạy: K
\begin{codeblock}
lakehouse/.venv/bin/python - <<'EOF'
import duckdb
D = 'lakehouse/source_sim/data/'
con = duckdb.connect()
for t in ['order_items', 'products']:
    con.sql(f"CREATE VIEW {t} AS SELECT * FROM '{D}olist_{t}_dataset.csv'")
con.sql("CREATE TABLE dim_product AS SELECT product_id FROM products")
con.sql("""CREATE TABLE fct_order_items AS
  SELECT order_id, order_item_id, product_id, CAST(price AS DECIMAL(12,2)) AS price
  FROM order_items""")
# Mỗi test là một câu SELECT trả về các dòng lỗi; 0 dòng nghĩa là đạt
tests = {
  'grain: (order_id, order_item_id) không trùng': """
     SELECT order_id, order_item_id FROM fct_order_items
     GROUP BY ALL HAVING count(*) > 1""",
  'khoá ngoại: product_id nào cũng có trong dim_product': """
     SELECT f.product_id FROM fct_order_items f
     ANTI JOIN dim_product p USING (product_id)""",
  'đối chiếu: tổng price bằng nguồn': """
     SELECT 1 WHERE (SELECT sum(price) FROM fct_order_items)
       IS DISTINCT FROM
       (SELECT sum(CAST(price AS DECIMAL(14,2))) FROM order_items)""",
}
def run_tests(label):
    print(label)
    for name, sql in tests.items():
        n = len(con.sql(sql).fetchall())
        print(f"  {'ĐẠT ' if n == 0 else 'LỖI '} {name} ({n} dòng lỗi)")
run_tests('Lần 1, dữ liệu sạch:')
# Cố ý làm hỏng: chép lại một dòng, và thêm một dòng trỏ tới sản phẩm không tồn tại
con.sql("INSERT INTO fct_order_items SELECT * FROM fct_order_items LIMIT 1")
con.sql("INSERT INTO fct_order_items VALUES ('x', 1, 'khong_co_san_pham', 10.00)")
run_tests('Lần 2, sau khi làm hỏng:')
EOF
\end{codeblock}
\textit{Kết quả:}
% kết quả: K
\begin{codeblock}
Lần 1, dữ liệu sạch:
  ĐẠT  grain: (order_id, order_item_id) không trùng (0 dòng lỗi)
  ĐẠT  khoá ngoại: product_id nào cũng có trong dim_product (0 dòng lỗi)
  ĐẠT  đối chiếu: tổng price bằng nguồn (0 dòng lỗi)
Lần 2, sau khi làm hỏng:
  LỖI  grain: (order_id, order_item_id) không trùng (1 dòng lỗi)
  LỖI  khoá ngoại: product_id nào cũng có trong dim_product (1 dòng lỗi)
  LỖI  đối chiếu: tổng price bằng nguồn (1 dòng lỗi)
\end{codeblock}

Lần đầu cả ba test đạt. Lần sau, dòng bị chép làm test grain đỏ, dòng trỏ tới sản phẩm không tồn tại làm test quan hệ đỏ, và cả hai dòng thừa làm tổng tiền không còn khớp nên test đối chiếu cũng đỏ. Mỗi test chỉ là một câu \code{SELECT}, rẻ hơn nhiều so với việc dò lỗi trên một dashboard đã sai.

\begin{luuy}
\textbf{Bẫy thường gặp.}
\begin{itemize}
  \item \textbf{Tin ràng buộc khoá trên bảng lakehouse.} Iceberg không ép tính duy nhất, và khoá của Databricks chỉ để thông tin.
  \item \textbf{Chỉ test cấu trúc, không test đối chiếu.} Test grain vẫn đạt khi một \code{INNER JOIN} làm rơi hẳn một nhóm dòng, như 775 đơn không có dòng hàng ở Bài 3. Chỉ test đối chiếu số dòng và tổng tiền mới bắt được.
  \item \textbf{Đối chiếu fact với chính nó.} Test đối chiếu phải so với nguồn, không so với một bảng dựng từ cùng câu SQL.
\end{itemize}
\end{luuy}

% =====================================================================
\section{Bài 11: Kimball, Inmon và Data Vault}

\textbf{Vấn đề.} Khi đọc tài liệu DE hay đi phỏng vấn, bạn sẽ gặp hai cái tên khác ngoài Kimball: Inmon và Data Vault. Một team khác có thể đã chọn một trong hai. Bạn cần biết mỗi cách đặt dữ liệu chi tiết ở đâu và người phân tích đọc bảng nào, để không tranh luận chỉ vì tên gọi.

\textbf{Giải thích.} Ba cách khác nhau chủ yếu ở chỗ dữ liệu chi tiết nhất được lưu dưới dạng nào.
\begin{itemize}
  \item \textbf{Kimball}, cách của module này: dữ liệu chi tiết được lưu ngay ở dạng dimension và fact, mỗi quy trình một star, các star ghép với nhau qua conformed dimension.
  \item \textbf{Inmon}, qua mô hình Corporate Information Factory, viết tắt CIF: dữ liệu chi tiết vào một kho trung tâm chuẩn hoá theo 3NF, và data mart của từng phòng ban chỉ được nạp từ kho đó.
  % Nguồn: https://williaminmon.substack.com/p/a-tale-of-two-architectures-kimball (bài viết của William Inmon, 23/12/2025): "the data in the corporate information factory data warehouse is stored in a normalized relational format"; "the only source of data for the data marts is the data warehouse".
  \item Margy Ross của Kimball Group đặt khác biệt cốt lõi ở cùng chỗ: Kimball đòi dữ liệu chi tiết có dạng dimension, còn CIF đặt nó trong kho chuẩn hoá; mart của CIF chứa dữ liệu đã gộp, dạng dimension.
  % Nguồn: https://decisionworks.com/2004/03/differences-of-opinion/ (bài viết của Margy Ross, 03/03/2004): "a third normal form (3NF) relational database containing atomic data"; "marts are tailored by business department/function with dimensionally structured summary data"; "The Kimball approach says atomic data must be dimensionally structured." Bài không nhắc tên Inmon, chỉ nói CIF.
  \item \textbf{Data Vault}, của Dan Linstedt, tách dữ liệu thành ba loại bảng:
  \begin{itemize}
    \item \textit{Hub} giữ danh sách khoá nghiệp vụ không trùng, kèm lúc khoá tới kho và hệ nguồn của nó.
    \item \textit{Link} giữ quan hệ nhiều nhiều giữa các hub.
    \item \textit{Satellite} giữ thuộc tính mô tả của một hub hoặc một link. Mỗi lần thuộc tính đổi, một dòng mới được thêm vào, kèm giờ nạp và hệ nguồn.
  \end{itemize}
  % Nguồn: https://dataversity.net/data-vault-series-1-data-vault-overview/5054 (bài viết của Dan Linstedt, TDAN.com, 07/2002): "a unique list of business keys"; "Load Date Time Stamp, recording when the key itself first arrived in the warehouse"; "a physical representation of a many-to-many 3NF relationship"; "Satellites may be connected to Hubs or Links"; "the new row is always inserted"; hub, link và satellite đều có "Record Source". Kiểm lại ngày 06/11/2026.
\end{itemize}

Hướng dẫn của Databricks cho Data Vault đặt mô hình dimension ở tầng gold, phía trên vault, và nhắc rằng Data Vault khuyên dùng mã băm của khoá nghiệp vụ làm khoá chính.
% Nguồn: https://www.databricks.com/blog/2022/06/24/prescriptive-guidance-for-implementing-a-data-vault-model-on-the-databricks-lakehouse-platform.html (24/06/2022): "Data Vault modeling recommends using a hash of business keys as the primary keys"; "loading a dimensional model Data Warehouse in Gold layer".

Repo \repo{lakehouse} đi gần Kimball hơn: \repo{core} đã có dạng dimension và fact, và Gold sẽ đọc thẳng \repo{core}. Nhưng nó có một chi tiết giống Data Vault: \code{customer\_\allowbreak{}key} của \repo{semantic.customer\_\allowbreak{}identity} là mã băm SHA-256 của \code{customer\_\allowbreak{}unique\_\allowbreak{}id}, như khoá băm của một hub.
% Nguồn: lakehouse/docs/data-model.md, mục 5, semantic.customer_identity: customer_key = sha2(customer_unique_id, 256).

Phác thảo Data Vault cho đơn hàng và người mua của Olist. Như sơ đồ ở Bài 1, dấu \code{<} hay \code{>} nằm ở phía có nhiều dòng:

\begin{codeblock}[samepage=true]
hub_customer               link_order_customer        hub_order
┌────────────────────┐     ┌────────────────────┐     ┌────────────────────┐
│ customer_hk        │────<│ customer_hk        │     │ order_hk           │
│ customer_unique_id │     │ order_hk           │>────│ order_id           │
│ load_ts            │     │ load_ts            │     │ load_ts            │
│ record_source      │     │ record_source      │     │ record_source      │
└─────────┬──────────┘     └────────────────────┘     └─────────┬──────────┘
          │ một hub, nhiều dòng satellite                       │
┌─────────┴──────────┐                                ┌─────────┴──────────┐
│ customer_hk        │                                │ order_hk           │
│ load_ts            │                                │ load_ts            │
│ record_source      │                                │ record_source      │
│ zip, city, state   │                                │ trạng thái, mốc    │
└────────────────────┘                                └────────────────────┘
sat_customer_address                                  sat_order_status
\end{codeblock}

\begin{ghichu}
\textbf{So với SWE.} Ba loại bảng của Data Vault đều quen với bạn. Hub giống một bảng chỉ có id và khoá nghiệp vụ. Link giống bảng nối mà Django tự tạo cho một \code{ManyToManyField}: tài liệu Django ghi rằng phía sau, Django tạo một bảng trung gian để biểu diễn quan hệ nhiều nhiều. Satellite giống một bảng lịch sử chỉ thêm, như audit log của một bảng.
% Nguồn: https://docs.djangoproject.com/en/6.1/ref/models/fields/ (ManyToManyField, Database Representation): "Behind the scenes, Django creates an intermediary join table to represent the many-to-many relationship."

\textbf{Chỗ so sánh gãy.} Trong app, bảng nối và bảng lịch sử là chi tiết nằm sau ORM; trong vault, chúng là toàn bộ mô hình. Một câu đơn giản như ``doanh thu theo bang của người mua'' phải đi qua hub đơn, link, hub người mua, rồi chọn đúng dòng satellite theo thời gian, trong khi star chỉ cần một join. Vault được thiết kế để nạp và kiểm toán, không để người phân tích đọc trực tiếp.
\end{ghichu}

Một bản dựng thử từ CSV, lấy giờ nạp là giờ đặt đơn như cách \repo{source\_\allowbreak{}sim} phát customer, có 96.096 dòng \code{hub\_\allowbreak{}customer}, đúng bằng số người mua, và 99.441 dòng cho \code{hub\_\allowbreak{}order} cũng như cho link. Satellite địa chỉ có 96.355 dòng, đúng bằng số version Type 2 khi theo dõi cả tên thành phố ở phần bẫy của Bài 9. Chính Linstedt so satellite với dimension Type 2 của Kimball.
% Kiểm: dựng hub_customer (sha256(customer_unique_id), min giờ đặt), hub_order, link_order_customer và sat_customer_address (thêm dòng khi (zip, city, state) IS DISTINCT FROM lag theo giờ đặt) từ orders JOIN customers; DuckDB 1.5.6, 06/11/2026: 96.096, 99.441, 99.441, 96.355 dòng.
% Nguồn: lakehouse/source_sim/README.md dòng 39: customers có updated_at bằng order_purchase_timestamp của đơn.
% Nguồn: Linstedt 2002 như trên: "The Satellite is most closely related to a Type 2 dimension as defined by Ralph Kimball."

\begin{luuy}
\textbf{Bẫy thường gặp.}
\begin{itemize}
  \item \textbf{Chọn Data Vault cho một nguồn nhỏ.} Mọi bảng của vault ghi hệ nguồn và giờ nạp, thứ có ích khi nhiều hệ nguồn cùng đổ về một kho. Olist chỉ có một nguồn, nên vault chỉ thêm join.
  \item \textbf{Cho người phân tích truy vấn thẳng vault.} Dựng star hoặc mart bên trên, như Gold của repo đọc \repo{core}.
  \item \textbf{Tranh luận Kimball hay Inmon bằng tên gọi.} Hỏi hai câu cụ thể: dữ liệu chi tiết nằm ở dạng nào, và người phân tích đọc bảng nào.
\end{itemize}
\end{luuy}

% =====================================================================
\section{Bài 12: Metric trong mart và phương trình đối soát}

\textbf{Vấn đề.} Giám đốc hỏi doanh thu của Olist là bao nhiêu. Bốn người trả lời bốn con số, từ 13,22 tới 16,01 triệu BRL, và không ai viết sai SQL. Mỗi người dùng một định nghĩa ``doanh thu'' riêng. Khi các con số đã lên dashboard, không ai còn biết dashboard nào dùng định nghĩa nào.

\textbf{Giải thích.} Một \textit{metric} là một con số có tên mà nhiều người cùng dùng. Định nghĩa của nó phải trả lời đủ năm câu: cộng cột nào, ở grain nào, lọc bỏ những gì, theo mốc thời gian nào, và tính bằng đơn vị gì. Định nghĩa được viết một lần, thành một cột trong mart, thay vì mỗi dashboard tự viết câu \code{WHERE} của mình. Bài 1 của Module 2 viết định nghĩa metric vào thẻ yêu cầu, giống tiêu chí chấp nhận của một user story; bài này thêm phần kiểm bằng số.
% Nguồn: roadmap/modules/m02_vong_doi.tex dòng 90 (cầu nối "user story có tiêu chí chấp nhận" với "định nghĩa metric") và dòng 145 (thẻ yêu cầu giống user story kèm tiêu chí chấp nhận), grep ngày 06/11/2026.

Đây là định nghĩa doanh thu mà Lab 3.5 dùng:

\begin{xltabular}{\textwidth}{@{}P{3.2cm} L@{}}
\toprule
\thd{trường} & \thd{doanh thu, cột \code{revenue} của \code{marts.revenue\_\allowbreak{}daily}} \\
\midrule
\endhead
\bottomrule
\endfoot
một câu & tổng giá hàng của các dòng hàng thuộc đơn không bị huỷ và không unavailable \\
cộng cột nào & \code{price} của \repo{core.order\_\allowbreak{}items} \\
không gồm & phí vận chuyển \code{freight\_\allowbreak{}value}, nằm ở cột \code{freight} riêng \\
bộ lọc & \code{order\_\allowbreak{}status NOT IN ('canceled', 'unavailable')} \\
mốc thời gian & ngày của \code{order\_\allowbreak{}purchase\_\allowbreak{}timestamp}, giờ địa phương Brazil \\
grain & một dòng mỗi ngày đặt, danh mục và bang người bán \\
đơn vị & BRL, kiểu \code{DECIMAL}, đúng tới 0,01 \\
phần bị loại & \code{marts.revenue\_\allowbreak{}excluded}, theo ngày và trạng thái đơn \\
\end{xltabular}
% Nguồn: lakehouse/docs/data-model.md mục 9 hàng 18 (giờ Brazil) và mục 11 (core.orders chứa cả đơn canceled và unavailable, lọc theo order_status).

Bộ lọc nào cũng bỏ đi một phần dữ liệu. \textit{Phương trình đối soát} (reconciliation) chứng minh phần bị bỏ là có chủ đích: phần giữ lại cộng phần bị loại phải bằng đúng tổng ở nguồn. Muốn kiểm được, phần bị loại phải được lưu thành một bảng có lý do, chứ không biến mất trong một câu \code{WHERE}.

\begin{ghichu}
\textbf{So với SWE.} Java có \code{Collectors.partitioningBy}: thay vì \code{filter()} bỏ đi phần không khớp, nó trả về cả hai nửa trong một \code{Map<Boolean, List<T>>}, và tài liệu Java ghi rằng map đó luôn có đủ hai khoá \code{false} và \code{true}. Mart nên làm giống vậy. Khi có đủ hai nửa, bạn kiểm được bất biến quen thuộc: hai nửa cộng lại bằng cái ban đầu.
% Nguồn: https://docs.oracle.com/en/java/javase/21/docs/api/java.base/java/util/stream/Collectors.html (partitioningBy): "The returned Map always contains mappings for both false and true keys."

\textbf{Chỗ so sánh gãy.} Với list, bất biến chỉ là đếm phần tử. Với tiền, phải so tổng bằng \code{DECIMAL} đúng tới 0,01 BRL, vì \code{DOUBLE} lệch ở chữ số cuối như Bài 1. Hai nửa cũng phải cùng grain với metric: lọc theo trạng thái của đơn, nhưng cộng theo dòng hàng. Cuối cùng, phương trình đúng chỉ chứng minh không đồng nào bị mất hay bị cộng hai lần; nó không chứng minh bỏ đơn huỷ là lựa chọn đúng. Lựa chọn đó là quyết định nghiệp vụ, và phải nằm trong định nghĩa.
\end{ghichu}

\subsection{Làm thử}

Khối dưới tính bốn con số cùng tên ``doanh thu'', dựng một mart theo định nghĩa trên cùng bảng phần bị loại, kiểm phương trình, rồi cố ý làm sai bảng phần bị loại.

% chạy: R
\begin{codeblock}
lakehouse/.venv/bin/python - <<'EOF'
import duckdb
D = 'lakehouse/source_sim/data/'
con = duckdb.connect()
for t in ['orders', 'order_items', 'order_payments']:
    con.sql(f"CREATE VIEW {t} AS SELECT * FROM '{D}olist_{t}_dataset.csv'")
con.sql("""CREATE VIEW items AS
  SELECT CAST(o.order_purchase_timestamp AS DATE) AS purchase_date, o.order_status,
         CAST(i.price AS DECIMAL(14,2)) AS price
  FROM order_items i JOIN orders o USING (order_id)""")
one = lambda q: con.sql(q).fetchall()[0][0]
print('Bốn con số cùng tên "doanh thu":')
print('  mọi dòng hàng                  :', one("SELECT sum(price) FROM items"))
print('  bỏ đơn canceled và unavailable :', one("""SELECT sum(price) FROM items
                                  WHERE order_status NOT IN ('canceled', 'unavailable')"""))
print('  chỉ đơn delivered              :', one("""SELECT sum(price) FROM items
                                  WHERE order_status = 'delivered'"""))
print('  tiền khách trả, payment_value  :', one("""SELECT sum(CAST(payment_value
                                  AS DECIMAL(14,2))) FROM order_payments"""))
# Mart: metric viết một lần; phần bị loại cũng thành một bảng, kèm lý do
con.sql("""CREATE TABLE mart_revenue_daily AS
  SELECT purchase_date, sum(price) AS revenue FROM items
  WHERE order_status NOT IN ('canceled', 'unavailable') GROUP BY 1""")
con.sql("""CREATE TABLE mart_revenue_excluded AS
  SELECT purchase_date, order_status AS reason, sum(price) AS revenue FROM items
  WHERE order_status IN ('canceled', 'unavailable') GROUP BY ALL""")
check = """SELECT k, x, k + x AS cong, s, k + x = s AS khop FROM (SELECT
  (SELECT sum(revenue) FROM mart_revenue_daily) AS k,
  (SELECT sum(revenue) FROM mart_revenue_excluded) AS x,
  (SELECT sum(CAST(price AS DECIMAL(14,2))) FROM order_items) AS s)"""
def doi_soat(label):
    k, x, c, s, ok = con.sql(check).fetchall()[0]
    print(f'{label}: {k} + {x} = {c}; nguồn {s}; khớp: {ok}')
doi_soat('giữ lại + bị loại')
# Lỗi cố ý: người viết bảng bị loại quên trạng thái unavailable
con.sql("DELETE FROM mart_revenue_excluded WHERE reason = 'unavailable'")
doi_soat('quên unavailable ')
EOF
\end{codeblock}
\textit{Kết quả:}
% kết quả: R
\begin{codeblock}
Bốn con số cùng tên "doanh thu":
  mọi dòng hàng                  : 13591643.70
  bỏ đơn canceled và unavailable : 13494400.74
  chỉ đơn delivered              : 13221498.11
  tiền khách trả, payment_value  : 16008872.12
giữ lại + bị loại: 13494400.74 + 97242.96 = 13591643.70; nguồn 13591643.70; khớp: True
quên unavailable : 13494400.74 + 95235.27 = 13589636.01; nguồn 13591643.70; khớp: False
\end{codeblock}

\begin{itemize}
  \item Mỗi con số đúng với định nghĩa của nó. Tiền khách trả còn gồm phí vận chuyển và tiền của 775 đơn không có dòng hàng ở Bài 3, nên không dùng được cho doanh thu theo danh mục.
  % Kiểm: 98.299 trên 98.665 đơn có cả dòng hàng và thanh toán có tổng payment_value bằng tổng price cộng freight_value, lệch không quá 0,01; DuckDB 1.5.6, 06/11/2026.
  \item Với định nghĩa của bài, phần giữ lại là 13.494.400,74 BRL và phần bị loại là 97.242,96 BRL. Cộng lại đúng bằng 13.591.643,70 BRL, tổng \code{price} của file CSV mà mục 10 của \repo{docs/data-model.md} dùng để kiểm chéo.
  \item Khi bảng phần bị loại quên trạng thái \code{unavailable}, phương trình lệch 2.007,69 BRL, tức khoảng 0,015\% tổng tiền. Một dashboard không bao giờ cho bạn thấy độ lệch nhỏ như vậy; phương trình thì thấy ngay.
  % Kiểm: 2.007,69 BRL là tổng price của 7 dòng hàng thuộc đơn unavailable; đếm bằng order_items JOIN orders GROUP BY order_status, DuckDB 1.5.6, 06/11/2026.
  \item Nếu chỉ tính đơn \code{delivered}, phần giữ lại là 13.221.498,11 BRL, phần bị loại là 370.145,59 BRL, và phương trình vẫn đúng. Đây là lý do phương trình không thay được định nghĩa.
  % Kiểm: sum(price) theo order_status <> 'delivered' = 370145.59, cùng lệnh, 06/11/2026.
\end{itemize}

\begin{luuy}
\textbf{Bẫy thường gặp.}
\begin{itemize}
  \item \textbf{Lọc trong dashboard thay vì trong mart.} Mỗi dashboard sẽ có một bộ lọc riêng, và bạn quay lại bốn con số ở phần vấn đề.
  \item \textbf{Chỉ đối soát tổng cả kỳ.} Tổng khớp không chứng minh tiền nằm đúng ngày. Đối soát thêm theo ngày hoặc theo tháng, nhất là khi đổi mốc thời gian hay múi giờ.
  \item \textbf{Đặt hai tên cho một con số, hoặc một tên cho hai con số.} Nếu \code{revenue} ở hai mart có hai định nghĩa, đổi tên một trong hai.
\end{itemize}
\end{luuy}

% =====================================================================
\section{Lab 3.1: Bus matrix và grain}

Lab này viết đặc tả trước khi viết SQL, giống việc viết interface trước khi cài đặt. Bạn cần giấy, một file markdown và mẫu truy vấn ở Bài 2.

\begin{enumerate}
  \item Với mỗi quy trình đặt hàng, thanh toán, giao hàng và đánh giá, viết một câu grain theo mẫu ``Một dòng là ...''.
  \item Vẽ bus matrix: dòng là quy trình, cột là \code{dim\_\allowbreak{}date} kèm các vai, \code{dim\_\allowbreak{}customer}, \code{dim\_\allowbreak{}product}, \code{dim\_\allowbreak{}seller} và \code{dim\_\allowbreak{}order\_\allowbreak{}profile}. Với mỗi ô bỏ trống, viết một câu lý do, như lý do người bán ở dòng giao hàng.
  \item Từ mỗi câu grain, viết ra khoá, rồi kiểm trên CSV bằng mẫu ở Bài 2: số dòng phải bằng số khoá khác nhau.
  \item Phân loại \code{price}, \code{freight\_\allowbreak{}value}, \code{payment\_\allowbreak{}value}, \code{payment\_\allowbreak{}installments}, \code{review\_\allowbreak{}score} thành cộng được, cộng được một phần hay không cộng được, kèm một câu lý do cho mỗi cột.
\end{enumerate}

\hoanthanh{bus matrix có đủ bốn quy trình; mỗi fact có một câu grain và một khoá; câu kiểm khoá của \code{fct\_\allowbreak{}order\_\allowbreak{}items} và \code{fct\_\allowbreak{}payments} cho số dòng bằng số khoá; năm measure đều có phân loại kèm lý do.}

% =====================================================================
\section{Lab 3.2: Dựng star schema Olist, test viết trước}

Lab này làm theo kiểu TDD: bạn nhận bộ test trước, chạy và thấy nó đỏ, rồi viết model cho tới khi mọi test xanh. Mọi thứ nằm trong thư mục \code{lab\_\allowbreak{}m03/}, và kết quả là file \code{lab\_\allowbreak{}m03/olist\_\allowbreak{}star.duckdb}.

\textbf{Bước 1: tạo runner.} Runner chạy các file \code{models/NN\_\allowbreak{}ten\_\allowbreak{}bang.sql} theo thứ tự tên, in số dòng của từng bảng, rồi chạy mọi test trong \code{tests/}. Nó trả mã lỗi 1 khi có test đỏ, nên sau này gắn vào CI được ngay.

% chạy: L1
\begin{codeblock}
mkdir -p lab_m03/models lab_m03/tests
cat > lab_m03/build.py <<'EOF'
import duckdb, pathlib, sys
D = 'lakehouse/source_sim/data/'
LAB = pathlib.Path('lab_m03')
con = duckdb.connect(str(LAB / 'olist_star.duckdb'))
# Nguồn là CSV Olist, đọc qua view tạm nên không bị chép vào file DuckDB
for t in ['orders', 'order_items', 'order_payments', 'order_reviews',
          'customers', 'products', 'sellers']:
    con.sql(f"CREATE OR REPLACE TEMP VIEW {t} AS "
            f"SELECT * FROM '{D}olist_{t}_dataset.csv'")
con.sql(f"CREATE OR REPLACE TEMP VIEW category_translation AS "
        f"SELECT * FROM '{D}product_category_name_translation.csv'")
# models/NN_ten_bang.sql tạo bảng ten_bang; chạy theo thứ tự tên file
for f in sorted((LAB / 'models').glob('*.sql')):
    con.execute(f.read_text())
    table = f.stem.split('_', 1)[1]
    rows = con.sql(f'SELECT count(*) FROM {table}').fetchall()[0][0]
    print(f'{f.name:26} {rows:>7} dòng')
# tests/*.sql: mỗi file là một câu SELECT trả về dòng lỗi; 0 dòng là đạt
tests = sorted((LAB / 'tests').glob('*.sql'))
failed = 0
for f in tests:
    try:
        n = len(con.sql(f.read_text()).fetchall())
        msg = f'{n} dòng lỗi'
    except duckdb.Error as e:
        n, msg = 1, str(e).splitlines()[0].split(': ', 1)[-1]
    failed += n > 0
    print(f"{'ĐẠT' if n == 0 else 'LỖI'}  {f.stem}: {msg}")
print(f'{len(tests)} test, {failed} lỗi')
sys.exit(1 if failed else 0)
EOF
\end{codeblock}

\textbf{Bước 2: model đầu tiên.} Mỗi file model bắt đầu bằng câu grain trong comment. \code{dim\_\allowbreak{}date} được cho sẵn làm mẫu.

% chạy: L2
\begin{codeblock}
cat > lab_m03/models/01_dim_date.sql <<'EOF'
/* Grain: một dòng mỗi ngày từ 2016-01-01 tới 2018-12-31,
   cộng một dòng khoá 0 cho ngày chưa biết hoặc chưa xảy ra. */
CREATE OR REPLACE TABLE dim_date AS
SELECT CAST(strftime(d, '%Y%m%d') AS INTEGER) AS date_key,
       CAST(d AS DATE) AS day, year(d) AS year, month(d) AS month,
       dayname(d) AS weekday, dayofweek(d) IN (0, 6) AS is_weekend
FROM range(DATE '2016-01-01', DATE '2019-01-01', INTERVAL 1 DAY) r(d)
UNION ALL
SELECT 0, NULL, NULL, NULL, 'unknown', NULL;
EOF
\end{codeblock}

\textbf{Bước 3: bộ test.} Tám test dưới là đặc tả của lab. Mỗi file là một câu \code{SELECT} trả về dòng lỗi. Tạo chúng rồi chạy runner.

% chạy: L3
\begin{codeblock}
cat > lab_m03/tests/dim_customer_grain.sql <<'EOF'
/* Mỗi người mua một dòng, mỗi customer_key một dòng */
SELECT customer_unique_id AS k FROM dim_customer GROUP BY 1 HAVING count(*) > 1
UNION ALL
SELECT CAST(customer_key AS VARCHAR) FROM dim_customer
GROUP BY 1 HAVING count(*) > 1;
EOF
cat > lab_m03/tests/dim_unknown_row.sql <<'EOF'
/* Dimension nào cũng có đúng một dòng khoá 0 cho giá trị chưa biết */
SELECT 'dim_date' AS dim FROM dim_date WHERE date_key = 0 HAVING count(*) <> 1
UNION ALL SELECT 'dim_customer' FROM dim_customer
          WHERE customer_key = 0 HAVING count(*) <> 1
UNION ALL SELECT 'dim_product' FROM dim_product
          WHERE product_key = 0 HAVING count(*) <> 1
UNION ALL SELECT 'dim_seller' FROM dim_seller
          WHERE seller_key = 0 HAVING count(*) <> 1;
EOF
cat > lab_m03/tests/fct_order_items_grain.sql <<'EOF'
/* Grain: một dòng mỗi (order_id, order_item_id) */
SELECT order_id, order_item_id FROM fct_order_items
GROUP BY ALL HAVING count(*) > 1;
EOF
cat > lab_m03/tests/fct_order_items_khoa_ngoai.sql <<'EOF'
/* Mọi khoá trên fact phải có dòng ở dimension; NULL cũng tính là lỗi */
SELECT 'customer_key' AS col, f.customer_key AS k FROM fct_order_items f
  LEFT JOIN dim_customer d USING (customer_key) WHERE d.customer_key IS NULL
UNION ALL SELECT 'product_key', f.product_key FROM fct_order_items f
  LEFT JOIN dim_product d USING (product_key) WHERE d.product_key IS NULL
UNION ALL SELECT 'seller_key', f.seller_key FROM fct_order_items f
  LEFT JOIN dim_seller d USING (seller_key) WHERE d.seller_key IS NULL
UNION ALL SELECT 'purchase_date_key', f.purchase_date_key FROM fct_order_items f
  LEFT JOIN dim_date d ON d.date_key = f.purchase_date_key WHERE d.date_key IS NULL
UNION ALL SELECT 'delivered_date_key', f.delivered_date_key FROM fct_order_items f
  LEFT JOIN dim_date d ON d.date_key = f.delivered_date_key
  WHERE d.date_key IS NULL;
EOF
cat > lab_m03/tests/fct_order_items_doi_chieu.sql <<'EOF'
/* Số dòng, tổng price, tổng freight_value phải bằng nguồn */
SELECT * FROM
  (SELECT count(*), sum(price), sum(freight_value) FROM fct_order_items)
EXCEPT
SELECT * FROM
  (SELECT count(*), sum(CAST(price AS DECIMAL(12,2))),
          sum(CAST(freight_value AS DECIMAL(12,2))) FROM order_items);
EOF
cat > lab_m03/tests/fct_payments_grain.sql <<'EOF'
/* Grain: một dòng mỗi (order_id, payment_sequential) */
SELECT order_id, payment_sequential FROM fct_payments
GROUP BY ALL HAVING count(*) > 1;
EOF
cat > lab_m03/tests/fct_payments_doi_chieu.sql <<'EOF'
/* Số dòng và tổng payment_value phải bằng nguồn */
SELECT * FROM (SELECT count(*), sum(payment_value) FROM fct_payments)
EXCEPT
SELECT * FROM (SELECT count(*), sum(CAST(payment_value AS DECIMAL(12,2)))
               FROM order_payments);
EOF
cat > lab_m03/tests/fct_payments_khoa_ngoai.sql <<'EOF'
/* Khoá người mua và khoá ngày đặt của thanh toán phải có ở dimension */
SELECT 'customer_key' AS col, f.customer_key AS k FROM fct_payments f
  LEFT JOIN dim_customer d USING (customer_key) WHERE d.customer_key IS NULL
UNION ALL SELECT 'purchase_date_key', f.purchase_date_key FROM fct_payments f
  LEFT JOIN dim_date d ON d.date_key = f.purchase_date_key WHERE d.date_key IS NULL;
EOF

lakehouse/.venv/bin/python lab_m03/build.py
\end{codeblock}
\textit{Kết quả:}
% kết quả: L3
\begin{codeblock}
01_dim_date.sql               1097 dòng
LỖI  dim_customer_grain: Table with name dim_customer does not exist!
LỖI  dim_unknown_row: Table with name dim_customer does not exist!
LỖI  fct_order_items_doi_chieu: Table with name fct_order_items does not exist!
LỖI  fct_order_items_grain: Table with name fct_order_items does not exist!
LỖI  fct_order_items_khoa_ngoai: Table with name fct_order_items does not exist!
LỖI  fct_payments_doi_chieu: Table with name fct_payments does not exist!
LỖI  fct_payments_grain: Table with name fct_payments does not exist!
LỖI  fct_payments_khoa_ngoai: Table with name fct_payments does not exist!
8 test, 8 lỗi
\end{codeblock}

Mọi test đỏ vì các bảng chưa tồn tại. Đây là điểm xuất phát bạn muốn.

\textbf{Bước 4: viết năm model.} Viết mỗi bảng thành một file trong \code{lab\_\allowbreak{}m03/models/}, theo đặc tả dưới. Bài 1, Bài 5 và Bài 6 đã có gần đủ các mảnh SQL bạn cần.

\begin{xltabular}{\textwidth}{@{}P{3.4cm} L L@{}}
\toprule
\thd{file và grain} & \thd{cột} & \thd{luật} \\
\midrule
\endhead
\bottomrule
\endfoot
\code{02\_\allowbreak{}dim\_\allowbreak{}customer.sql}: một dòng mỗi \code{customer\_\allowbreak{}unique\_\allowbreak{}id}, cộng dòng khoá 0 & \code{customer\_\allowbreak{}key}, \code{customer\_\allowbreak{}unique\_\allowbreak{}id}, \code{customer\_\allowbreak{}city}, \code{customer\_\allowbreak{}state} & khoá thay thế bằng \code{row\_\allowbreak{}number()} từ 1; địa chỉ lấy ở đơn mới nhất, hai đơn cùng lúc thì lấy \code{customer\_\allowbreak{}id} lớn hơn, như \repo{core.customers} \\
\code{03\_\allowbreak{}dim\_\allowbreak{}product.sql}: một dòng mỗi \code{product\_\allowbreak{}id}, cộng dòng khoá 0 & \code{product\_\allowbreak{}key}, \code{product\_\allowbreak{}id}, \code{category}, \code{product\_\allowbreak{}weight\_\allowbreak{}g}, \code{product\_\allowbreak{}photos\_\allowbreak{}qty} & \code{category} là tên tiếng Anh; không có bản dịch thì dùng tên gốc; không có tên gốc thì ghi \code{unknown} \\
\code{04\_\allowbreak{}dim\_\allowbreak{}seller.sql}: một dòng mỗi \code{seller\_\allowbreak{}id}, cộng dòng khoá 0 & \code{seller\_\allowbreak{}key}, \code{seller\_\allowbreak{}id}, \code{seller\_\allowbreak{}city}, \code{seller\_\allowbreak{}state} & khoá thay thế bằng \code{row\_\allowbreak{}number()} từ 1 \\
\code{05\_\allowbreak{}fct\_\allowbreak{}order\_\allowbreak{}items.sql}: một dòng mỗi \code{order\_\allowbreak{}id} và \code{order\_\allowbreak{}item\_\allowbreak{}id} & \code{order\_\allowbreak{}id}, \code{order\_\allowbreak{}item\_\allowbreak{}id}, \code{customer\_\allowbreak{}key}, \code{product\_\allowbreak{}key}, \code{seller\_\allowbreak{}key}, \code{purchase\_\allowbreak{}date\_\allowbreak{}key}, \code{delivered\_\allowbreak{}date\_\allowbreak{}key}, \code{order\_\allowbreak{}status}, \code{price}, \code{freight\_\allowbreak{}value} & tra người mua qua \code{orders.customer\_\allowbreak{}id} rồi \code{customers.customer\_\allowbreak{}unique\_\allowbreak{}id}; khoá không tìm thấy thì ghi 0; đơn chưa giao có \code{delivered\_\allowbreak{}date\_\allowbreak{}key} bằng 0; tiền có kiểu \code{DECIMAL(12,2)}; \code{order\_\allowbreak{}id} là degenerate dimension; \code{order\_\allowbreak{}status} tạm nằm thẳng trên fact thay cho khoá của \code{dim\_\allowbreak{}order\_\allowbreak{}profile} \\
\code{06\_\allowbreak{}fct\_\allowbreak{}payments.sql}: một dòng mỗi \code{order\_\allowbreak{}id} và \code{payment\_\allowbreak{}sequential} & \code{order\_\allowbreak{}id}, \code{payment\_\allowbreak{}sequential}, \code{customer\_\allowbreak{}key}, \code{purchase\_\allowbreak{}date\_\allowbreak{}key}, \code{payment\_\allowbreak{}type}, \code{payment\_\allowbreak{}installments}, \code{payment\_\allowbreak{}value} & tra khoá như \code{fct\_\allowbreak{}order\_\allowbreak{}items}; tiền có kiểu \code{DECIMAL(12,2)} \\
\end{xltabular}

\textbf{Bước 5: chạy lại cho tới khi xanh.} Khi cả năm file đúng, runner in ra:

% chạy: L5 cần-lời-giải
\begin{codeblock}
lakehouse/.venv/bin/python lab_m03/build.py
\end{codeblock}
\textit{Kết quả:}
% kết quả: L5
\begin{codeblock}
01_dim_date.sql               1097 dòng
02_dim_customer.sql          96097 dòng
03_dim_product.sql           32952 dòng
04_dim_seller.sql             3096 dòng
05_fct_order_items.sql      112650 dòng
06_fct_payments.sql         103886 dòng
ĐẠT  dim_customer_grain: 0 dòng lỗi
ĐẠT  dim_unknown_row: 0 dòng lỗi
ĐẠT  fct_order_items_doi_chieu: 0 dòng lỗi
ĐẠT  fct_order_items_grain: 0 dòng lỗi
ĐẠT  fct_order_items_khoa_ngoai: 0 dòng lỗi
ĐẠT  fct_payments_doi_chieu: 0 dòng lỗi
ĐẠT  fct_payments_grain: 0 dòng lỗi
ĐẠT  fct_payments_khoa_ngoai: 0 dòng lỗi
8 test, 0 lỗi
\end{codeblock}

\textbf{Bước 6: chứng minh fan-out trên star của bạn.} Viết hai câu truy vấn trên \code{olist\_\allowbreak{}star.duckdb}. Câu thứ nhất join \code{fct\_\allowbreak{}payments} với \code{fct\_\allowbreak{}order\_\allowbreak{}items} theo \code{order\_\allowbreak{}id} rồi cộng \code{payment\_\allowbreak{}value}. Câu thứ hai gộp mỗi fact về grain đơn trước, rồi mới join, và giữ cả đơn có thanh toán mà không có dòng hàng.

\hoanthanh{runner in \code{8 test, 0 lỗi} với đúng số dòng như trên; mỗi file model bắt đầu bằng câu grain; câu thứ nhất ở bước 6 cho 20.308.134,71 BRL và câu thứ hai cho 16.008.872,12 BRL.}

% =====================================================================
\section{Lab 3.3: SCD Type 2 cho người bán}

Olist chỉ có một ảnh chụp bảng \code{sellers}, nên bạn tự tạo thay đổi để luyện. Bài 9 dựng lịch sử một lần từ quá khứ. Lab này làm việc mà pipeline phải làm mỗi ngày: nhận thay đổi mới, đóng version cũ, mở version mới, và chạy lại thì không sinh thêm dòng.

\textbf{Bước 1: tạo dimension, dữ liệu thay đổi và câu SQL áp thay đổi.} Ba người bán có nhiều dòng hàng nhất chuyển thành phố vào 2018-01-01, và có thêm một người bán mới. Câu áp thay đổi chạy trong một transaction: bước đầu đóng version hiện hành của người bán có thay đổi, bước sau thêm version mới cho người bán không còn dòng hiện hành.

% chạy: M1
\begin{codeblock}
mkdir -p lab_m03/scd2
cat > lab_m03/scd2/00_init.sql <<'EOF'
/* Nạp lần đầu: mỗi người bán một version, hiệu lực từ 2016-01-01 */
CREATE OR REPLACE TABLE dim_seller_scd2 AS
SELECT row_number() OVER (ORDER BY seller_id) AS seller_key,
       seller_id, seller_city, seller_state,
       DATE '2016-01-01' AS valid_from, CAST(NULL AS DATE) AS valid_to,
       true AS is_current
FROM sellers;
/* Thay đổi giả lập: ba người bán chuyển thành phố, một người bán mới */
CREATE OR REPLACE TABLE seller_changes AS
SELECT * FROM (VALUES
  ('6560211a19b47992c3666cc44a7e94c0', 'campinas',       'SP', DATE '2018-01-01'),
  ('4a3ca9315b744ce9f8e9374361493884', 'curitiba',       'PR', DATE '2018-01-01'),
  ('1f50f920176fa81dab994f9023523100', 'belo horizonte', 'MG', DATE '2018-01-01'),
  ('new_seller_0001',                  'recife',         'PE', DATE '2018-01-01'))
  t(seller_id, seller_city, seller_state, change_date);
EOF
cat > lab_m03/scd2/01_apply.sql <<'EOF'
BEGIN;
/* Bước 1: đóng version hiện hành của người bán có thuộc tính đổi */
UPDATE dim_seller_scd2 d
SET valid_to = c.change_date, is_current = false
FROM seller_changes c
WHERE d.seller_id = c.seller_id AND d.is_current
  AND (d.seller_city IS DISTINCT FROM c.seller_city
       OR d.seller_state IS DISTINCT FROM c.seller_state);
/* Bước 2: thêm version mới cho người bán không còn dòng hiện hành */
INSERT INTO dim_seller_scd2
SELECT (SELECT max(seller_key) FROM dim_seller_scd2)
         + row_number() OVER (ORDER BY c.seller_id),
       c.seller_id, c.seller_city, c.seller_state,
       c.change_date, NULL, true
FROM seller_changes c
WHERE NOT EXISTS (SELECT 1 FROM dim_seller_scd2 d
                  WHERE d.seller_id = c.seller_id AND d.is_current);
COMMIT;
EOF
\end{codeblock}

\textbf{Bước 2: nạp lần đầu, rồi áp thay đổi hai lần.} Lần áp thứ hai phải không đổi gì; tính chất này gọi là \textit{idempotent}, giống một API \code{PUT} gọi lại hai lần vẫn cho cùng trạng thái.
% Nguồn: https://www.rfc-editor.org/rfc/rfc9110.txt (mục 9.2.2)

% chạy: M2
\begin{codeblock}
lakehouse/.venv/bin/python - <<'EOF'
import duckdb
D = 'lakehouse/source_sim/data/'
con = duckdb.connect('lab_m03/olist_star.duckdb')
con.sql(f"CREATE TEMP VIEW sellers AS SELECT * FROM '{D}olist_sellers_dataset.csv'")
con.execute(open('lab_m03/scd2/00_init.sql').read())
count = "SELECT count(*), count(*) FILTER (WHERE is_current) FROM dim_seller_scd2"
print('sau khi nạp lần đầu, (số dòng, số dòng hiện hành):',
      con.sql(count).fetchall()[0])
for lan in [1, 2]:
    con.execute(open('lab_m03/scd2/01_apply.sql').read())
    print(f'sau lần áp thay đổi thứ {lan}:', con.sql(count).fetchall()[0])
con.sql("""SELECT seller_key, seller_city, seller_state,
                  valid_from, valid_to, is_current
           FROM dim_seller_scd2
           WHERE seller_id = '4a3ca9315b744ce9f8e9374361493884'
           ORDER BY valid_from""").show()
EOF
\end{codeblock}
\textit{Kết quả:}
% kết quả: M2
\begin{codeblock}
sau khi nạp lần đầu, (số dòng, số dòng hiện hành): (3095, 3095)
sau lần áp thay đổi thứ 1: (3099, 3096)
sau lần áp thay đổi thứ 2: (3099, 3096)
┌────────────┬─────────────┬──────────────┬────────────┬────────────┬────────────┐
│ seller_key │ seller_city │ seller_state │ valid_from │  valid_to  │ is_current │
│   int64    │   varchar   │   varchar    │    date    │    date    │  boolean   │
├────────────┼─────────────┼──────────────┼────────────┼────────────┼────────────┤
│        882 │ ibitinga    │ SP           │ 2016-01-01 │ 2018-01-01 │ false      │
│       3097 │ curitiba    │ PR           │ 2018-01-01 │ NULL       │ true       │
└────────────┴─────────────┴──────────────┴────────────┴────────────┴────────────┘
\end{codeblock}

\textbf{Bước 3: test.} Thêm hai test vào \code{lab\_\allowbreak{}m03/tests/}: mỗi \code{seller\_\allowbreak{}id} có đúng một dòng \code{is\_\allowbreak{}current}, và không có hai version nào của cùng một người bán chồng khoảng \code{[valid\_\allowbreak{}from, valid\_\allowbreak{}to)}. Chạy lại runner.

\textbf{Bước 4: tra version cho từng dòng hàng.} Viết câu truy vấn gắn mỗi dòng của \code{fct\_\allowbreak{}order\_\allowbreak{}items} với version người bán còn hiệu lực vào ngày đặt, dùng khoảng nửa mở như Bài 9.

\hoanthanh{lần áp thứ hai vẫn giữ 3.099 dòng và 3.096 dòng hiện hành; runner in \code{10 test, 0 lỗi}; cả 112.650 dòng hàng khớp đúng một version người bán; với ba người bán đã chuyển, 3.260 dòng hàng trỏ tới version cũ và 2.691 dòng hàng trỏ tới version mới.}

% =====================================================================
\section{Lab 3.4: Accumulating snapshot từ API của repo}

Ở Bài 4, \code{fct\_\allowbreak{}orders} được dựng một lần từ CSV, tức từ trạng thái cuối của mọi đơn. Pipeline thật không có trạng thái cuối: nó nhận từng version của đơn theo thời gian, và phải sửa dòng của đơn trong snapshot mỗi khi đơn qua một mốc. Lab này làm đúng việc đó với \repo{source\_\allowbreak{}sim}, server giả lập nguồn REST của repo. Module 8 sẽ làm phần kéo dữ liệu này cho Bronze thật; ở đây mọi thứ nằm trong \code{lab\_\allowbreak{}m03/snapshot/}.

Server chạy như một process Python ở cổng 8010, vì cổng 8000 thuộc service \code{source-sim} của docker compose. Nó nạp mọi version vào DuckDB trong RAM, chỉ ghi file tạm của hệ điều hành, và không ghi gì vào \repo{lakehouse/}.
% Nguồn: lakehouse/docker-compose.yml:149 (127.0.0.1:8000:8000); lakehouse/source_sim/app.py:39-49 (DuckDB in-memory, NamedTemporaryFile).

\begin{luuy}
\textbf{Thật và giả lập trong dữ liệu của source\_sim.}
\begin{itemize}
  \item \textbf{Thật:} mọi mốc thời gian, vì chúng chép từ CSV, và trạng thái ở version cuối, vì version cuối đúng bằng dòng CSV.
  \item \textbf{Giả lập:} trạng thái của mọi version trung gian. Mock gán \code{created}, \code{approved}, \code{shipped}, \code{delivered} theo mốc xa nhất đã tới, và không cho trạng thái lùi.
  \item \textbf{Giả lập:} lúc đơn chuyển sang \code{canceled} hay \code{unavailable}. Olist không ghi mốc này, nên trạng thái thật hiện ra ở mốc cuối cùng của đơn.
  \item \textbf{Cần nhớ:} mock coi mọi mốc là UTC, nhưng chúng là giờ địa phương Brazil.
  \item \textbf{Lạ nhưng thật:} 166 đơn có giờ giao cho hãng vận chuyển sớm hơn giờ đặt ngay trong CSV, nên version đầu của chúng là \code{shipped} mà chưa có giờ đặt.
\end{itemize}
% Nguồn: lakehouse/source_sim/README.md dòng 29 và dòng 42 tới 44; lakehouse/source_sim/replay.py:59-73; lakehouse/docs/data-model.md mục 9 hàng 17 và hàng 18.
\end{luuy}

\textbf{Bước 1: chạy server.} \code{SIM\_\allowbreak{}FIXED\_\allowbreak{}NOW} đóng băng đồng hồ mô phỏng. Đặt nó sau version cuối của dữ liệu thì mọi version hiện ra ngay; README của mock dùng đúng giá trị 2018-12-31T00:00:00 cho việc này. \code{PYTHONDONTWRITEBYTECODE=1} giữ cho Python không ghi file \code{.pyc} vào repo.
% Nguồn: lakehouse/source_sim/README.md dòng 54.

% chạy: N1
\begin{codeblock}
mkdir -p lab_m03/snapshot
PYTHONDONTWRITEBYTECODE=1 SIM_FIXED_NOW=2018-12-31T00:00:00 \
  lakehouse/.venv/bin/python -m uvicorn --app-dir lakehouse source_sim.app:app \
  --port 8010 > lab_m03/snapshot/sim.log 2>&1 < /dev/null &
echo $! > lab_m03/snapshot/sim.pid
# Server cần vài giây để nạp mọi version vào RAM trước khi trả lời
for i in $(seq 60); do curl -s localhost:8010/health && break; sleep 1; done
echo
curl -s 'localhost:8010/api/v1/orders?page_size=1'; echo
\end{codeblock}
\textit{Kết quả:}
% kết quả: N1
\begin{codeblock}
{"status":"ok","sim_time":"2018-12-31T00:00:00"}
{"entity":"orders","page_size":1,"has_more":true,"sim_time":"2018-12-31T00:00:00","data":[{"order_id":"2e7a8482f6fb09756ca50c10d7bfc047","customer_id":"08c5351a6aca1c1589a38f244edeee9d","order_status":"created","order_purchase_timestamp":"2016-09-04 21:15:19","order_approved_at":null,"order_delivered_carrier_date":null,"order_delivered_customer_date":null,"order_estimated_delivery_date":"2016-10-20 00:00:00","updated_at":"2016-09-04T21:15:19","cursor":"2016-09-04T21:15:19|000000001"}]}
\end{codeblock}

Dòng đầu là \code{/health}, sau đó là trang đầu của \code{orders} với một version. Nếu vòng chờ hết 60 giây mà không in gì, mở \code{lab\_\allowbreak{}m03/snapshot/sim.log}. Dòng \code{address already in use} nghĩa là cổng 8010 đang bận: tìm process bằng \code{lsof -nP -iTCP:8010} rồi dừng nó.

\textbf{Bước 2: kéo mọi version vào bảng thô.} Mỗi version có \code{cursor} dạng \code{<updated\_\allowbreak{}at>|<seq>}. Trang sau gửi \code{after} bằng cursor của version cuối trang trước, và dừng khi \code{has\_\allowbreak{}more} là \code{false}. Client chỉ so cursor như chuỗi, không tách nó ra. Bảng \code{raw\_\allowbreak{}order\_\allowbreak{}versions} giữ mọi version kèm số trang, giống Bronze chỉ thêm và không sửa.
% Nguồn: lakehouse/source_sim/README.md dòng 13 tới 27 (cursor, after, has_more, page_size tối đa 5000).

% chạy: N2
\begin{codeblock}
lakehouse/.venv/bin/python - <<'EOF'
import duckdb, json, urllib.parse, urllib.request
import pyarrow as pa
COLS = ['order_id', 'customer_id', 'order_status', 'order_purchase_timestamp',
        'order_approved_at', 'order_delivered_carrier_date',
        'order_delivered_customer_date', 'order_estimated_delivery_date',
        'updated_at', 'cursor']
con = duckdb.connect('lab_m03/snapshot/orders.duckdb')
con.sql("CREATE OR REPLACE TABLE raw_order_versions (page INTEGER, "
        + ', '.join(f'{c} VARCHAR' for c in COLS) + ")")
page, after = 0, None
while True:
    url = 'http://localhost:8010/api/v1/orders?page_size=5000'
    if after:  # trang sau bắt đầu ngay sau cursor của version cuối trang trước
        url += '&after=' + urllib.parse.quote(after)
    body = json.load(urllib.request.urlopen(url))
    page += 1
    batch = pa.Table.from_pylist(body['data'],
                                 schema=pa.schema([(c, pa.string()) for c in COLS]))
    con.sql(f"INSERT INTO raw_order_versions SELECT {page}, * FROM batch")
    after = body['data'][-1]['cursor']
    if not body['has_more']:
        break
print('số trang:', page)
print('số version, số đơn, số cursor, version cuối:', con.sql("""SELECT count(*),
  count(DISTINCT order_id), count(DISTINCT cursor), max(updated_at)
  FROM raw_order_versions""").fetchall()[0])
print('(số version của một đơn, số đơn):', con.sql("""SELECT n, count(*) FROM
  (SELECT count(*) AS n FROM raw_order_versions GROUP BY order_id)
  GROUP BY 1 ORDER BY 1""").fetchall())
EOF
\end{codeblock}
\textit{Kết quả:}
% kết quả: N2
\begin{codeblock}
số trang: 79
số version, số đơn, số cursor, version cuối: (391551, 99441, 391551, '2018-10-17T17:30:18')
(số version của một đơn, số đơn): [(1, 177), (2, 1626), (3, 2430), (4, 95208)]
\end{codeblock}

79 trang chứa 391.551 version của 99.441 đơn, mỗi version một cursor riêng, và phân bố số version mỗi đơn khớp đúng hàng 17 ở mục 9 của \repo{docs/data-model.md}. Nếu bạn nhận ít version hơn, \code{sim\_\allowbreak{}time} ở bước 1 đang sớm hơn 2018-10-17T17:30:18, giờ của version cuối.

\textbf{Bước 3: xem version của hai đơn.} Khối dưới in các version của một đơn đã giao và một đơn bị huỷ, kèm các mốc đã có ở mỗi version.

% chạy: N3
\begin{codeblock}
lakehouse/.venv/bin/python - <<'EOF'
import duckdb
con = duckdb.connect('lab_m03/snapshot/orders.duckdb', read_only=True)
MOC = {'order_purchase_timestamp': 'đặt', 'order_approved_at': 'duyệt',
       'order_delivered_carrier_date': 'giao vận chuyển',
       'order_delivered_customer_date': 'tới khách'}
for oid in ['00010242fe8c5a6d1ba2dd792cb16214', '0166cc8756d58b4a30fb49de83527120']:
    print('đơn', oid[:8])
    rows = con.sql(f"""SELECT updated_at, order_status, {', '.join(MOC)}
                       FROM raw_order_versions WHERE order_id = '{oid}'
                       ORDER BY cursor""").fetchall()
    for r in rows:
        co = [ten for ten, v in zip(MOC.values(), r[2:]) if v]
        print(f'  {r[0]}  {r[1]:9}  mốc đã có: {", ".join(co)}')
EOF
\end{codeblock}
\textit{Kết quả:}
% kết quả: N3
\begin{codeblock}
đơn 00010242
  2017-09-13T08:59:02  created    mốc đã có: đặt
  2017-09-13T09:45:35  approved   mốc đã có: đặt, duyệt
  2017-09-19T18:34:16  shipped    mốc đã có: đặt, duyệt, giao vận chuyển
  2017-09-20T23:43:48  delivered  mốc đã có: đặt, duyệt, giao vận chuyển, tới khách
đơn 0166cc87
  2018-02-08T21:09:18  created    mốc đã có: đặt
  2018-02-08T21:27:56  approved   mốc đã có: đặt, duyệt
  2018-02-09T23:12:30  canceled   mốc đã có: đặt, duyệt, giao vận chuyển
\end{codeblock}

Ở đơn \code{00010242}, bốn giờ ở cột đầu là thật, còn ba trạng thái đầu là mock suy ra từ mốc đã tới. Đơn \code{0166cc87} thành \code{canceled} đúng vào giờ giao cho hãng vận chuyển: đó là giờ mock chọn, không phải giờ đơn bị huỷ, vì Olist không ghi giờ huỷ.

\textbf{Bước 4: áp từng trang vào snapshot bằng \code{MERGE}.} Snapshot có grain một dòng mỗi đơn. Câu \code{MERGE} thêm dòng cho đơn mới và sửa dòng của đơn đã có, như upsert ở Module 1, Bài 10. Hai chi tiết giữ cho bảng đúng. Thứ nhất, một trang có thể chứa nhiều version của cùng một đơn, nên câu nguồn chỉ giữ version có cursor lớn nhất của mỗi đơn. Thứ hai, điều kiện \code{s.last\_\allowbreak{}cursor > t.last\_\allowbreak{}cursor} chỉ cho version mới hơn ghi đè, nên áp lại một trang hay áp lộn thứ tự đều vô hại. \code{UPDATE BY NAME} và \code{INSERT BY NAME} ghép cột theo tên.

% chạy: N4
\begin{codeblock}
cat > lab_m03/snapshot/00_create.sql <<'EOF'
/* Grain: một dòng mỗi đơn. Mỗi mốc một cột, được điền dần khi đơn qua mốc đó. */
CREATE OR REPLACE TABLE fct_order_snapshot (
  order_id VARCHAR, customer_id VARCHAR, order_status VARCHAR,
  purchase_ts TIMESTAMP, approved_ts TIMESTAMP, carrier_ts TIMESTAMP,
  delivered_ts TIMESTAMP, estimated_date DATE, last_cursor VARCHAR);
EOF
cat > lab_m03/snapshot/01_apply_page.sql <<'EOF'
/* Áp các version của một trang vào snapshot; $page là số trang */
MERGE INTO fct_order_snapshot t
USING (
  SELECT order_id, customer_id, order_status,
         CAST(order_purchase_timestamp AS TIMESTAMP)      AS purchase_ts,
         CAST(order_approved_at AS TIMESTAMP)             AS approved_ts,
         CAST(order_delivered_carrier_date AS TIMESTAMP)  AS carrier_ts,
         CAST(order_delivered_customer_date AS TIMESTAMP) AS delivered_ts,
         CAST(order_estimated_delivery_date AS DATE)      AS estimated_date,
         cursor AS last_cursor
  FROM raw_order_versions WHERE page = $page
  /* Một trang có thể chứa nhiều version của cùng đơn: giữ version mới nhất */
  QUALIFY row_number() OVER (PARTITION BY order_id ORDER BY cursor DESC) = 1) s
ON t.order_id = s.order_id
/* Chỉ ghi đè bằng version mới hơn: áp lại hay áp lộn thứ tự đều không làm lùi */
WHEN MATCHED AND s.last_cursor > t.last_cursor THEN UPDATE BY NAME
WHEN NOT MATCHED THEN INSERT BY NAME;
EOF
lakehouse/.venv/bin/python - <<'EOF'
import duckdb
con = duckdb.connect('lab_m03/snapshot/orders.duckdb')
con.execute(open('lab_m03/snapshot/00_create.sql').read())
apply = open('lab_m03/snapshot/01_apply_page.sql').read()
def run(pages):  # số dòng snapshot mà MERGE đã thêm hoặc sửa
    return sum(con.execute(apply, {'page': p}).fetchall()[0][0] for p in pages)
print('áp trang 1 tới 40, số dòng đổi:', run(range(1, 41)))
# Đơn chưa tới version cuối của nó đang mang trạng thái do mock giả lập
con.sql("""SELECT s.order_status, count(*) FILTER (WHERE s.last_cursor = f.c) AS that,
                  count(*) FILTER (WHERE s.last_cursor < f.c) AS gia_lap
           FROM fct_order_snapshot s JOIN (SELECT order_id, max(cursor) AS c
                FROM raw_order_versions GROUP BY 1) f USING (order_id)
           GROUP BY 1 ORDER BY count(*) DESC""").show()
print('áp trang 41 tới 79, số dòng đổi:', run(range(41, 80)))
print('áp lại cả 79 trang, số dòng đổi:', run(range(1, 80)))
print('áp trang 79 rồi trang 1, số dòng đổi:', run([79, 1]))
csv = """SELECT order_id, customer_id, order_status, order_purchase_timestamp,
           order_approved_at, order_delivered_carrier_date,
           order_delivered_customer_date, CAST(order_estimated_delivery_date AS DATE)
         FROM 'lakehouse/source_sim/data/olist_orders_dataset.csv'"""
snap = "SELECT * EXCLUDE (last_cursor) FROM fct_order_snapshot"
print('số dòng, lệch với CSV hai chiều:', con.sql(f"""SELECT
  (SELECT count(*) FROM fct_order_snapshot),
  (SELECT count(*) FROM ({snap} EXCEPT {csv})),
  (SELECT count(*) FROM ({csv} EXCEPT {snap}))""").fetchall()[0])
EOF
\end{codeblock}
\textit{Kết quả:}
% kết quả: N4
\begin{codeblock}
áp trang 1 tới 40, số dòng đổi: 105619
┌──────────────┬───────┬─────────┐
│ order_status │ that  │ gia_lap │
│   varchar    │ int64 │  int64  │
├──────────────┼───────┼─────────┤
│ delivered    │ 47136 │       0 │
│ shipped      │   601 │    2262 │
│ approved     │     2 │     754 │
│ unavailable  │   509 │       0 │
│ canceled     │   308 │       0 │
│ processing   │   271 │       0 │
│ invoiced     │   208 │       0 │
│ created      │     4 │     122 │
└──────────────┴───────┴─────────┘

áp trang 41 tới 79, số dòng đổi: 111064
áp lại cả 79 trang, số dòng đổi: 0
áp trang 79 rồi trang 1, số dòng đổi: 0
số dòng, lệch với CSV hai chiều: (99441, 0, 0)
\end{codeblock}

\begin{itemize}
  \item Sau 40 trang, snapshot có 52.177 đơn, và 3.138 đơn đang mang trạng thái giả lập: 2.262 \code{shipped}, 754 \code{approved}, 122 \code{created}. Đếm ``số đơn đang chờ duyệt'' trên bảng này là đếm số do mock tạo ra.
  \item Áp lại cả 79 trang, rồi áp trang 79 trước trang 1, đều đổi 0 dòng. Đó là tính idempotent bạn đã gặp ở Lab 3.3.
  \item Snapshot cuối có 99.441 dòng và trùng khớp với CSV ở cả tám cột, theo cả hai chiều \code{EXCEPT}.
\end{itemize}

Nếu snapshot có hơn 99.441 dòng, câu nguồn thiếu bước giữ version mới nhất; nếu lần áp lại đổi dòng nào, điều kiện so cursor đang thiếu hoặc sai chiều.

\begin{luuy}
\textbf{Bẫy thường gặp.}
\begin{itemize}
  \item \textbf{Bỏ bước giữ version mới nhất trong trang.} Engine nào cũng cần bước này; chỉ cách hỏng là khác nhau:
  \begin{itemize}
    \item DuckDB 1.5.6 không báo lỗi: trang 1 thêm 5.000 dòng cho 1.562 đơn, và snapshot cuối có 257.365 dòng cho 99.441 đơn.
    \item DuckDB 2.0.0 dừng ở trang 2 với lỗi \code{MERGE INTO command cannot affect the same target row more than once}.
    \item PostgreSQL từ chối một lệnh sửa cùng một dòng hai lần, với lỗi \code{MERGE command cannot affect row a second time} mà bạn đã gặp ở Module 1, Bài 10.
  \end{itemize}
  % Kiểm: chạy 01_apply_page.sql không có dòng QUALIFY trên bản chép của raw_order_versions, DuckDB 1.5.6 và 2.0.0, 06/11/2026. PostgreSQL: roadmap/modules/m01_sql.tex, Bài 10, output "ERROR:  MERGE command cannot affect row a second time"; https://www.postgresql.org/docs/18/sql-merge.html ("a repeated UPDATE or DELETE will cause a cardinality violation").
  \item \textbf{Viết \code{UPDATE} hay \code{INSERT} không có \code{BY NAME}.} DuckDB khi đó ghép cột theo vị trí, cả với \code{UPDATE SET *} và \code{INSERT *}; bảng nguồn đổi thứ tự cột là dữ liệu vào sai cột.
  % Kiểm: MERGE với nguồn có cột (b, k, a) vào bảng (k, a, b): UPDATE và INSERT * báo Conversion Error do ghép theo vị trí, BY NAME ghép đúng; DuckDB 1.5.6 và 2.0.0, 06/11/2026.
  \item \textbf{Tin trạng thái trung gian.} Với dữ liệu thật, trạng thái là thứ hệ nguồn ghi; với mock, chỉ các mốc thời gian là thật. Metric về thời gian xử lý nên tính từ các cột mốc, như Bài 4.
\end{itemize}
\end{luuy}

\textbf{Bước 5: dựng lại snapshot mà không dùng \code{MERGE}.} Viết file \code{lab\_\allowbreak{}m03/snapshot/02\_\allowbreak{}as\_\allowbreak{}of.sql} tạo temp view \code{order\_\allowbreak{}snapshot\_\allowbreak{}as\_\allowbreak{}of}: trạng thái mọi đơn sau trang 40, dựng thẳng từ \code{raw\_\allowbreak{}order\_\allowbreak{}versions} bằng cách giữ version có cursor lớn nhất của mỗi đơn, với đúng chín cột của snapshot. Khối dưới dựng lại snapshot bằng \code{MERGE} tới trang 40 rồi so hai cách. Cách dựng thẳng chậm hơn khi dữ liệu lớn, nhưng là thước đo để kiểm cách tăng dần.

% chạy: N5 cần-lời-giải
\begin{codeblock}
lakehouse/.venv/bin/python - <<'EOF'
import duckdb
con = duckdb.connect()  # làm trong RAM, không đụng snapshot của bước 4
con.sql("ATTACH 'lab_m03/snapshot/orders.duckdb' AS lab (READ_ONLY)")
con.sql("CREATE TEMP VIEW raw_order_versions AS SELECT * FROM lab.raw_order_versions")
con.execute(open('lab_m03/snapshot/00_create.sql').read())
apply = open('lab_m03/snapshot/01_apply_page.sql').read()
for p in range(1, 41):
    con.execute(apply, {'page': p})
con.execute(open('lab_m03/snapshot/02_as_of.sql').read())  # file bạn viết
a, b = 'fct_order_snapshot', 'order_snapshot_as_of'
print('số dòng của hai cách:', con.sql(f"""SELECT (SELECT count(*) FROM {a}),
  (SELECT count(*) FROM {b})""").fetchall()[0])
print('lệch hai chiều:', con.sql(f"""SELECT
  (SELECT count(*) FROM (SELECT * FROM {a} EXCEPT SELECT * FROM {b})),
  (SELECT count(*) FROM (SELECT * FROM {b} EXCEPT SELECT * FROM {a}))
  """).fetchall()[0])
EOF
\end{codeblock}
\textit{Kết quả:}
% kết quả: N5
\begin{codeblock}
số dòng của hai cách: (52177, 52177)
lệch hai chiều: (0, 0)
\end{codeblock}

Hai cách cùng ra 52.177 dòng và không lệch dòng nào. Nếu lệch, so tên và kiểu của chín cột trước, rồi kiểm điều kiện \code{page <= 40}.

\textbf{Bước 6: dừng server.}

% chạy: N6
\begin{codeblock}
kill "$(cat lab_m03/snapshot/sim.pid)"
while kill -0 "$(cat lab_m03/snapshot/sim.pid)" 2>/dev/null; do sleep 0.2; done
grep -c 'Application shutdown complete' lab_m03/snapshot/sim.log
curl -s localhost:8010/health || echo "không còn server ở cổng 8010, curl trả mã $?"
\end{codeblock}
\textit{Kết quả:}
% kết quả: N6
\begin{codeblock}
1
không còn server ở cổng 8010, curl trả mã 7
\end{codeblock}

Số 1 là số dòng \code{Application shutdown complete} trong log, tức server đã tắt hẳn; mã 7 của \code{curl} nghĩa là không kết nối được.

\hoanthanh{khối ở bước 4 in \code{(99441, 0, 0)} ở dòng cuối, và hai lần áp lại đều đổi 0 dòng; khối ở bước 5 in \code{(52177, 52177)} và \code{(0, 0)}; sau bước 6, \code{curl -s localhost:8010/health} trả mã 7.}

% =====================================================================
\section{Lab 3.5: Mart cho năm câu hỏi của repo}

Mục 11 của \repo{docs/data-model.md} liệt kê năm câu hỏi mà Gold của repo phải trả lời, cùng các bảng \repo{core} cần đọc. Lab này thiết kế và dựng năm mart đó trên DuckDB, từ một bản \repo{core} thu nhỏ dựng từ CSV, với đúng tên bảng và tên cột của repo. Đây là bản chạy thử. Ở M7, bạn viết lại các mart này bằng Spark SQL trên \repo{core} thật, và phương trình đối soát của lab này là tiêu chí hoàn thành ở đó.

\textbf{Bước 1: dựng core thu nhỏ và runner.} File \code{00\_\allowbreak{}core.sql} được cho sẵn, vì core là việc của tầng Silver, mà repo đã làm bằng Spark trong \repo{src/\allowbreak{}lakehouse/\allowbreak{}silver/\allowbreak{}core.py}. Runner chạy mọi file \code{.sql} trong \code{lab\_\allowbreak{}m03/gold/} theo thứ tự tên, in số dòng của từng bảng, rồi chạy năm test. Mỗi test so mart với CSV chứ không so với core, để bắt được cả lỗi của core lẫn lỗi của mart.

% chạy: P1
\begin{codeblock}
mkdir -p lab_m03/gold
cat > lab_m03/gold/00_core.sql <<'EOF'
/* Bản thu nhỏ của semantic và core trong repo, dựng từ CSV; tên bảng, tên cột
   theo docs/data-model.md mục 5 và 6, chỉ giữ cột mà năm câu hỏi ở mục 11 cần. */
CREATE SCHEMA IF NOT EXISTS semantic;
CREATE SCHEMA IF NOT EXISTS core;
/* Toạ độ mỗi zip: trung bình các điểm khác nhau nằm trong khung Brazil */
CREATE OR REPLACE TABLE semantic.zip_geo AS
SELECT zip_code_prefix, avg(lat) AS lat, avg(lng) AS lng FROM (
  SELECT DISTINCT lpad(CAST(geolocation_zip_code_prefix AS VARCHAR), 5, '0')
           AS zip_code_prefix, geolocation_lat AS lat, geolocation_lng AS lng,
         geolocation_city, geolocation_state
  FROM geolocation)
WHERE lat BETWEEN -33.75 AND 5.27 AND lng BETWEEN -73.99 AND -34.79
GROUP BY 1;
CREATE OR REPLACE TABLE core.products AS
SELECT p.product_id, coalesce(t.product_category_name_english,
         p.product_category_name, 'unknown') AS product_category_name_english
FROM products p LEFT JOIN category_translation t USING (product_category_name);
CREATE OR REPLACE TABLE core.sellers AS
SELECT s.seller_id, s.zip AS seller_zip_code_prefix, s.seller_state,
       g.lat AS seller_lat, g.lng AS seller_lng
FROM (SELECT *, lpad(CAST(seller_zip_code_prefix AS VARCHAR), 5, '0') AS zip
      FROM sellers) s
LEFT JOIN semantic.zip_geo g ON g.zip_code_prefix = s.zip;
CREATE OR REPLACE TABLE core.customers AS
WITH b AS (  /* địa chỉ ở đơn mới nhất; đặt cùng lúc thì lấy customer_id lớn hơn */
  SELECT c.customer_unique_id,
         arg_max(lpad(CAST(c.customer_zip_code_prefix AS VARCHAR), 5, '0'),
                 (o.order_purchase_timestamp, c.customer_id)) AS zip,
         arg_max(c.customer_state, (o.order_purchase_timestamp, c.customer_id))
           AS customer_state,
         min(o.order_purchase_timestamp) AS first_order_purchase_timestamp,
         max(o.order_purchase_timestamp) AS last_order_purchase_timestamp,
         count(DISTINCT c.customer_id) AS customer_id_count
  FROM customers c JOIN orders o USING (customer_id) GROUP BY 1)
SELECT sha256(customer_unique_id) AS customer_key, customer_unique_id,
       zip AS customer_zip_code_prefix, customer_state,
       g.lat AS customer_lat, g.lng AS customer_lng,
       first_order_purchase_timestamp, last_order_purchase_timestamp,
       customer_id_count
FROM b LEFT JOIN semantic.zip_geo g ON g.zip_code_prefix = b.zip;
CREATE OR REPLACE TABLE core.orders AS
SELECT o.order_id, sha256(c.customer_unique_id) AS customer_key, o.order_status,
       o.order_purchase_timestamp, o.order_delivered_customer_date, r.review_score,
       CAST(o.order_delivered_customer_date AS DATE)
         - CAST(o.order_purchase_timestamp AS DATE) AS delivery_days,
       CAST(o.order_delivered_customer_date AS DATE)
         > CAST(o.order_estimated_delivery_date AS DATE) AS is_delivered_late
FROM orders o JOIN customers c USING (customer_id)
LEFT JOIN (SELECT order_id, review_score FROM order_reviews  /* review mới nhất */
           QUALIFY row_number() OVER (PARTITION BY order_id ORDER BY
             review_creation_date DESC, review_answer_timestamp DESC,
             review_id DESC) = 1) r USING (order_id);
CREATE OR REPLACE TABLE core.order_items AS
SELECT order_id, order_item_id, product_id, seller_id,
       CAST(price AS DECIMAL(12,2)) AS price,
       CAST(freight_value AS DECIMAL(12,2)) AS freight_value
FROM order_items;
EOF
cat > lab_m03/gold/build_gold.py <<'EOF'
import duckdb, pathlib, sys
D = 'lakehouse/source_sim/data/'
G = pathlib.Path('lab_m03/gold')
con = duckdb.connect(str(G / 'gold.duckdb'))
for t in ['orders', 'order_items', 'order_reviews', 'customers', 'products',
          'sellers', 'geolocation']:
    con.sql(f"CREATE OR REPLACE TEMP VIEW {t} AS "
            f"SELECT * FROM '{D}olist_{t}_dataset.csv'")
con.sql(f"CREATE OR REPLACE TEMP VIEW category_translation AS "
        f"SELECT * FROM '{D}product_category_name_translation.csv'")
for f in sorted(G.glob('*.sql')):  # 00_core.sql, rồi 10_marts.sql khi bạn đã viết
    con.execute(f.read_text())
for s, t in con.sql("SELECT schema_name, table_name FROM duckdb_tables() "
                    "WHERE database_name = 'gold' ORDER BY ALL").fetchall():
    n = con.sql(f'SELECT count(*) FROM {s}.{t}').fetchall()[0][0]
    print(f'{s + "." + t:28} {n:>7} dòng')
GRAIN = {'revenue_daily': 'purchase_date, category, seller_state',
         'revenue_excluded': 'purchase_date, order_status',
         'delivery_state_month': 'purchase_month, customer_state',
         'repeat_buyers_month': 'first_purchase_month',
         'review_seller': 'seller_id', 'map_zip': 'zip_code_prefix'}
# Mỗi test trả về dòng lỗi, 0 dòng là đạt; so mart với CSV, không so với core
TESTS = {
  'grain sáu mart': ' UNION ALL '.join(
      f"SELECT '{m}' FROM marts.{m} GROUP BY {k} HAVING count(*) > 1"
      for m, k in GRAIN.items()),
  'doanh thu: giữ + loại = nguồn': """SELECT 1 WHERE
     (SELECT sum(revenue) FROM marts.revenue_daily)
     + (SELECT sum(revenue) FROM marts.revenue_excluded)
     IS DISTINCT FROM (SELECT sum(CAST(price AS DECIMAL(12,2))) FROM order_items)""",
  'giao hàng: đã tới + chưa tới = số đơn': """SELECT 1 WHERE
     (SELECT sum(delivered_orders) FROM marts.delivery_state_month)
     + (SELECT count(*) FROM orders WHERE order_delivered_customer_date IS NULL)
     IS DISTINCT FROM (SELECT count(*) FROM orders)""",
  'người mua: hai mart = nguồn': """SELECT n FROM (
     SELECT sum(buyers) AS n FROM marts.repeat_buyers_month
     UNION ALL SELECT sum(buyers) FROM marts.map_zip)
     WHERE n IS DISTINCT FROM (SELECT count(DISTINCT customer_unique_id)
                               FROM customers)""",
  'review: mỗi cặp đơn, người bán một lần': """SELECT 1 WHERE
     (SELECT sum(reviewed_orders) FROM marts.review_seller)
     IS DISTINCT FROM (SELECT count(*) FROM (SELECT DISTINCT order_id, seller_id
       FROM order_items WHERE order_id IN (SELECT order_id FROM order_reviews)))""",
}
failed = 0
for name, sql in TESTS.items():
    try:
        n = len(con.sql(sql).fetchall())
        msg = f'{n} dòng lỗi'
    except duckdb.Error as e:
        n, msg = 1, str(e).splitlines()[0].split(': ', 1)[-1]
    failed += n > 0
    print(f"{'ĐẠT' if n == 0 else 'LỖI'}  {name}: {msg}")
if not failed:
    k, x = con.sql("SELECT (SELECT sum(revenue) FROM marts.revenue_daily), "
                   "(SELECT sum(revenue) FROM marts.revenue_excluded)").fetchall()[0]
    print(f'đối soát doanh thu: {k} + {x} = {k + x}')
sys.exit(1 if failed else 0)
EOF
lakehouse/.venv/bin/python lab_m03/gold/build_gold.py
\end{codeblock}
\textit{Kết quả:}
% kết quả: P1
\begin{codeblock}
core.customers                 96096 dòng
core.order_items              112650 dòng
core.orders                    99441 dòng
core.products                  32951 dòng
core.sellers                    3095 dòng
semantic.zip_geo               19010 dòng
LỖI  grain sáu mart: Table with name "marts.revenue_daily" does not exist because schema "marts" does not exist.
LỖI  doanh thu: giữ + loại = nguồn: Table with name "marts.revenue_daily" does not exist because schema "marts" does not exist.
LỖI  giao hàng: đã tới + chưa tới = số đơn: Table with name "marts.delivery_state_month" does not exist because schema "marts" does not exist.
LỖI  người mua: hai mart = nguồn: Table with name "marts.repeat_buyers_month" does not exist because schema "marts" does not exist.
LỖI  review: mỗi cặp đơn, người bán một lần: Table with name "marts.review_seller" does not exist because schema "marts" does not exist.
\end{codeblock}

Số dòng của sáu bảng khớp mục 10 của \repo{docs/data-model.md}: 19.010 zip, 96.096 người mua, 3.095 người bán, 32.951 sản phẩm, 99.441 đơn và 112.650 dòng hàng. Năm test đỏ vì schema \code{marts} chưa có. Nếu số dòng khác, kiểm biến \code{D} trỏ đúng thư mục CSV.

\textbf{Bước 2: đọc đặc tả.} Mỗi mart có grain, cột và luật dưới đây; nếu đã làm Lab 2.2, so bảng này với thẻ yêu cầu của bạn. Mỗi \code{CREATE TABLE} trong file của bạn phải mở đầu bằng comment ghi câu grain và định nghĩa một câu của từng measure, theo mẫu định nghĩa metric ở Bài 12.

\begin{xltabular}{\textwidth}{@{}P{3.8cm} P{4.2cm} L@{}}
\toprule
\thd{mart và grain} & \thd{cột} & \thd{luật} \\
\midrule
\endhead
\bottomrule
\endfoot
\code{revenue\_\allowbreak{}daily}: một dòng mỗi ngày đặt, danh mục, bang người bán & \code{purchase\_\allowbreak{}date}, \code{category}, \code{seller\_\allowbreak{}state}, \code{item\_\allowbreak{}count}, \code{revenue}, \code{freight} & định nghĩa doanh thu ở Bài 12; \code{category} là \code{product\_\allowbreak{}category\_\allowbreak{}name\_\allowbreak{}english} \\
\code{revenue\_\allowbreak{}excluded}: một dòng mỗi ngày đặt và trạng thái đơn & \code{purchase\_\allowbreak{}date}, \code{order\_\allowbreak{}status}, \code{item\_\allowbreak{}count}, \code{revenue} & chỉ đơn \code{canceled} và \code{unavailable} \\
\code{delivery\_\allowbreak{}state\_\allowbreak{}month}: một dòng mỗi tháng đặt và bang người mua & \code{purchase\_\allowbreak{}month}, \code{customer\_\allowbreak{}state}, \code{delivered\_\allowbreak{}orders}, \code{late\_\allowbreak{}orders}, \code{delivery\_\allowbreak{}days\_\allowbreak{}sum} & chỉ đơn đã tới tay khách; lưu tử số và mẫu số, không lưu tỉ lệ hay trung bình \\
\code{repeat\_\allowbreak{}buyers\_\allowbreak{}month}: một dòng mỗi tháng mua lần đầu & \code{first\_\allowbreak{}purchase\_\allowbreak{}month}, \code{buyers}, \code{repeat\_\allowbreak{}buyers} & người mua quay lại có \code{customer\_\allowbreak{}id\_\allowbreak{}count} lớn hơn 1 \\
\code{review\_\allowbreak{}seller}: một dòng mỗi người bán & \code{seller\_\allowbreak{}id}, \code{reviewed\_\allowbreak{}orders}, \code{review\_\allowbreak{}score\_\allowbreak{}sum} & mỗi cặp đơn và người bán đếm một lần, dù đơn có nhiều dòng hàng của người bán đó \\
\code{map\_\allowbreak{}zip}: một dòng mỗi zip của người mua hoặc người bán & \code{zip\_\allowbreak{}code\_\allowbreak{}prefix}, \code{lat}, \code{lng}, \code{buyers}, \code{sellers} & toạ độ từ \code{semantic.zip\_\allowbreak{}geo}; zip không có toạ độ vẫn giữ, với \code{lat} là NULL \\
\end{xltabular}

\textbf{Bước 3: viết \code{lab\_\allowbreak{}m03/gold/10\_\allowbreak{}marts.sql}.} File bắt đầu bằng \code{CREATE SCHEMA IF NOT EXISTS marts;} và chỉ đọc \code{core} cùng \code{semantic}. Bài 3, Bài 4, Bài 7 và Bài 12 có đủ các mảnh bạn cần.

\textbf{Bước 4: chạy lại runner cho tới khi xanh.} Khi cả sáu bảng đúng, runner in ra:

% chạy: P2 cần-lời-giải
\begin{codeblock}
lakehouse/.venv/bin/python lab_m03/gold/build_gold.py
\end{codeblock}
\textit{Kết quả:}
% kết quả: P2
\begin{codeblock}
core.customers                 96096 dòng
core.order_items              112650 dòng
core.orders                    99441 dòng
core.products                  32951 dòng
core.sellers                    3095 dòng
marts.delivery_state_month       556 dòng
marts.map_zip                  15068 dòng
marts.repeat_buyers_month         25 dòng
marts.revenue_daily            35548 dòng
marts.revenue_excluded           286 dòng
marts.review_seller             3090 dòng
semantic.zip_geo               19010 dòng
ĐẠT  grain sáu mart: 0 dòng lỗi
ĐẠT  doanh thu: giữ + loại = nguồn: 0 dòng lỗi
ĐẠT  giao hàng: đã tới + chưa tới = số đơn: 0 dòng lỗi
ĐẠT  người mua: hai mart = nguồn: 0 dòng lỗi
ĐẠT  review: mỗi cặp đơn, người bán một lần: 0 dòng lỗi
đối soát doanh thu: 13494400.74 + 97242.96 = 13591643.70
\end{codeblock}

Nếu một test đỏ, tìm tên test trong danh sách dưới:
\begin{itemize}
  \item \textbf{Grain:} chạy riêng câu \code{GROUP BY} của từng mart trong \code{GRAIN} để tìm mart có khoá trùng. Ví dụ, nhóm theo \code{order\_\allowbreak{}purchase\_\allowbreak{}timestamp} thay vì theo ngày sẽ cho nhiều dòng cùng một ngày.
  \item \textbf{Doanh thu:} cộng \code{revenue} của hai mart. Thiếu đúng 2.007,69 BRL so với 13.591.643,70 nghĩa là bạn quên \code{unavailable} trong \code{revenue\_\allowbreak{}excluded}, như Bài 12.
  \item \textbf{Giao hàng:} tính tổng \code{delivered\_\allowbreak{}orders}. 96.478 thay vì 96.476 nghĩa là bạn lọc theo trạng thái \code{delivered} thay vì theo \code{order\_\allowbreak{}delivered\_\allowbreak{}customer\_\allowbreak{}date IS NOT NULL}.
  \item \textbf{Người mua:} tính tổng \code{buyers} của \code{map\_\allowbreak{}zip}. 95.826 thay vì 96.096 nghĩa là bạn dùng \code{INNER JOIN} với \code{zip\_\allowbreak{}geo} và làm rơi 270 người mua không có toạ độ.
  \item \textbf{Review:} tính tổng \code{reviewed\_\allowbreak{}orders}. 111.708 thay vì 99.247 nghĩa là bạn đang đếm mỗi dòng hàng thay vì mỗi cặp đơn và người bán.
\end{itemize}
% Kiểm: lời giải tham khảo với INNER JOIN zip_geo cho sum(buyers) = 95.826; đếm dòng hàng có review cho 111.708; DuckDB 1.5.6, 06/11/2026.
% Kiểm: olist_orders_dataset.csv, DuckDB 1.5.6, 06/11/2026: order_status = 'delivered' cho 96.478 đơn, order_delivered_customer_date IS NOT NULL cho 96.476 đơn, NULL cho 2.965 đơn; 2.007,69 BRL là phần lệch của khối R ở Bài 12.
% Kiểm: runner với lời giải tham khảo bị sửa từng chỗ, DuckDB 1.5.6, 06/11/2026: lọc order_status = 'delivered' cho tổng delivered_orders 96.478 và chỉ test giao hàng đỏ; revenue_excluded chỉ có canceled cho tổng hai mart 13.589.636,01 và chỉ test doanh thu đỏ; revenue_daily GROUP BY order_purchase_timestamp cho test grain 16.147 dòng lỗi.

\hoanthanh{runner in năm dòng \code{ĐẠT} và dòng \code{đối soát doanh thu: 13494400.74 + 97242.96 = 13591643.70}; mỗi \code{CREATE TABLE} có câu grain và định nghĩa measure trong comment.}

Thêm nếu còn giờ: đổi định nghĩa doanh thu sang chỉ đơn \code{delivered}. Phương trình vẫn đúng, với 13.221.498,11 BRL giữ lại và 370.145,59 BRL bị loại.

% =====================================================================
\section{Thử thách kỹ năng}

\begin{thuthach}
Không có hướng dẫn từng bước. Thêm quy trình đánh giá vào Gold của Lab 3.5.

Thêm vào \code{lab\_\allowbreak{}m03/gold/} một bảng \code{fct\_\allowbreak{}reviews} với grain do bạn tự xác định, sau khi kiểm xem \code{review\_\allowbreak{}id} có duy nhất trong dữ liệu không. Từ đó dựng \code{marts.review\_\allowbreak{}category\_\allowbreak{}month}: điểm review theo danh mục và tháng tạo review, lưu tử số và mẫu số. Tự quyết và ghi lại ba điều: đơn có nhiều review thì tính review nào, đơn có hàng của hai danh mục thì điểm tính cho danh mục nào, và review của đơn bị huỷ có được tính không. Runner chỉ kiểm grain của bảng trong schema \code{marts}, nên đặt \code{fct\_\allowbreak{}reviews} ở đó, hoặc sửa runner để \code{GRAIN} nhận tên có schema.
\hoanthanh{bảng và mart mới có câu grain trong comment và có mặt trong \code{GRAIN} của runner; phương trình đối soát cho review đúng: số review được tính cộng số review bị loại, mỗi phần có lý do, bằng 99.224; runner in 0 dòng \code{LỖI}.}

\textbf{Mang vào capstone.} Mốc capstone 1 là bus matrix, grain và tầng Gold trên DuckDB cho năm câu hỏi ở mục 11 của \repo{docs/\allowbreak{}data-model.md}, có phương trình đối soát. Đem bus matrix của Lab 3.1, star của Lab 3.2, snapshot của Lab 3.4, mart của Lab 3.5 và phần review này vào repo capstone, kèm lệnh chạy runner. Dùng thẻ yêu cầu của Lab 2.2 làm đặc tả cho năm mart. Chạy lại năm câu của thử thách Module 1 trên star của Lab 3.2 và snapshot của Lab 3.4, rồi so với \code{lab\_\allowbreak{}m01/\allowbreak{}thu\_\allowbreak{}thach.sql}.
\end{thuthach}
% Nguồn: lakehouse/docs/data-model.md mục 9 hàng 3 (99.224 review, 98.410 review_id); mốc 1 của lộ trình.
% Nguồn: roadmap/modules/m01_sql.tex, thử thách kỹ năng, dòng Mang vào capstone (mốc 1 khớp lab_m01/thu_thach.sql); roadmap/modules/m02_vong_doi.tex, Lab 2.2 và dòng Mang vào capstone (thẻ yêu cầu là đầu vào của mốc 1); grep ngày 06/11/2026.

% =====================================================================
\section{Tự kiểm tra}

Trả lời bằng lời của bạn trước, rồi mới đọc gợi ý.

\begin{enumerate}
  \item Vì sao grain phải được viết trước khi chọn cột? Join thanh toán với dòng hàng rồi cộng \code{payment\_\allowbreak{}value} cho 20.308.134,71 BRL; sửa thế nào để ra 16.008.872,12 BRL, và vì sao \code{sum(DISTINCT ...)} không phải là cách sửa?\par
  \textit{Gợi ý:} grain là hợp đồng mà mọi cột phải khớp; một dòng hàng và một lần trả tiền là hai grain. Gộp mỗi bảng về grain đơn rồi \code{LEFT JOIN}; \code{DISTINCT} bỏ cả những khoản tiền trùng số của các đơn khác nhau.
  \item Loại fact nào hợp với câu hỏi ``mỗi bước giao hàng mất bao lâu''? Vì sao tỉ lệ giao trễ của cả nước không phải trung bình tỉ lệ của 27 bang?\par
  \textit{Gợi ý:} accumulating snapshot, vì mọi mốc của một đơn nằm trên cùng một dòng. Tỉ lệ không cộng được: cộng tử số và mẫu số rồi mới chia, ra 6,77\% thay vì 9,15\%.
  \item Vì sao tỉ lệ khách quay lại tính theo \code{customer\_\allowbreak{}id} của Olist luôn bằng 0? Khoá tự nhiên, khoá thay thế và khoá bền khác nhau thế nào trong một dimension Type 2?\par
  \textit{Gợi ý:} mỗi đơn có một \code{customer\_\allowbreak{}id} mới; người mua thật là \code{customer\_\allowbreak{}unique\_\allowbreak{}id}. Nghĩ tới S3 Versioning: khoá bền giống key của object, khoá thay thế giống version ID, còn khoá tự nhiên là thứ hệ nguồn đưa tới và có thể đổi.
  \item Một người bán chuyển từ SP sang RJ. Viết các dòng dimension sau thay đổi theo Type 1, Type 2 và Type 3. Vì sao khoảng hiệu lực của Type 2 là nửa mở?\par
  \textit{Gợi ý:} Type 1 còn một dòng RJ; Type 2 có hai dòng nối nhau tại ngày đổi; Type 3 có một dòng với bang hiện tại là RJ và bang trước đó là SP. Với \code{BETWEEN}, 256 đơn Olist khớp hai version và bị nhân dòng.
  \item Vì sao tách \code{dim\_\allowbreak{}category} không tiết kiệm chỗ trên Parquet? Cái giá thật của OBT là gì?\par
  \textit{Gợi ý:} dictionary encoding lưu mỗi tên danh mục một lần, cả cột \code{category} chỉ tốn 30.231 byte; năm cột chép vào OBT chỉ chiếm 5,6\% file. Cái giá là cập nhật: đổi tên \code{bed\_\allowbreak{}bath\_\allowbreak{}table} phải ghi lại 11.115 dòng, và dễ lỡ tay trộn hai grain trong một bảng.
  \item Trong Kimball, Inmon và Data Vault, dữ liệu chi tiết được lưu ở dạng nào, và người phân tích đọc bảng nào?\par
  \textit{Gợi ý:} Kimball lưu ngay ở dạng dimension và fact, người phân tích đọc star. Inmon lưu trong kho trung tâm 3NF, người phân tích đọc data mart dạng dimension nạp từ kho. Data Vault lưu ở hub, link và satellite, người phân tích đọc star dựng bên trên.
  \item Khi áp các trang version của \repo{source\_\allowbreak{}sim} vào snapshot bằng \code{MERGE}, hai chi tiết nào giữ cho bảng đúng, và bỏ mỗi chi tiết thì hỏng thế nào? Phần nào của các version là thật?\par
  \textit{Gợi ý:} giữ version mới nhất của mỗi đơn trong trang, nếu không DuckDB 1.5.6 cho 257.365 dòng thay vì 99.441; chỉ ghi đè bằng cursor lớn hơn, nếu không áp lại hay áp lộn thứ tự sẽ làm dòng lùi về version cũ. Các mốc thời gian và trạng thái ở version cuối là thật; trạng thái trung gian là mock giả lập.
  \item Viết phương trình đối soát của doanh thu Olist theo định nghĩa ở Bài 12. Nó chứng minh điều gì, và không chứng minh điều gì?\par
  \textit{Gợi ý:} 13.494.400,74 cộng 97.242,96 bằng 13.591.643,70. Nó chứng minh không đồng nào bị mất hay bị cộng hai lần, nhưng không chứng minh bỏ đơn huỷ là đúng: định nghĩa chỉ giữ đơn \code{delivered} cũng thoả phương trình.
\end{enumerate}

% =====================================================================
\section{Đọc thêm}

Đọc sau khi đã làm xong các lab, để gắn tên chuẩn cho những gì bạn đã tự làm. Thời gian ở mỗi mục là ước lượng.

\begin{itemize}
  \item \href{https://www.kimballgroup.com/wp-content/uploads/2013/08/2013.09-Kimball-Dimensional-Modeling-Techniques11.pdf}{Kimball Dimensional Modeling Techniques}: tài liệu PDF 24 trang của Kimball Group, bản 2013. Đọc trang 1 tới 15, khoảng 60 phút: mỗi kỹ thuật trong module có ở đó một đoạn ngắn, kèm tên chuẩn.
  \item \href{https://dataversity.net/data-vault-series-1-data-vault-overview/5054}{Dan Linstedt: Data Vault Series 1}: bài viết năm 2002 của người tạo ra Data Vault. Đọc phần định nghĩa hub, link và satellite, khoảng 20 phút; đây là bản đầu, chưa có khoá băm của bản 2.0.
  \item \href{https://www.fivetran.com/blog/star-schema-vs-obt}{Michael Kaminsky: Star Schema vs. OBT}: bài viết ngày 16/08/2022, đo hai cách trên Redshift, Snowflake và BigQuery. Đọc phần kết quả và phần chỗ lưu, khoảng 15 phút; số chỗ lưu là của Redshift, không phải của Parquet.
  \item \href{https://decisionworks.com/2004/03/differences-of-opinion/}{Margy Ross: Differences of Opinion}: bài viết năm 2004 so cách của Kimball với Corporate Information Factory của Inmon, nhìn từ phía Kimball. Đọc cả bài, khoảng 15 phút.
  \item \href{https://github.com/DataTalksClub/data-engineering-zoomcamp/tree/main/04-analytics-engineering}{DE Zoomcamp, Module 4 Analytics Engineering}: miễn phí, nhánh main của repo, xem ngày 06/11/2026. Unit 1 tóm tắt fact và dimension theo Kimball, khoảng 30 phút; các unit sau thuộc nhánh T1.
  \item \href{https://www.coursera.org/learn/data-modeling-transformation-serving}{DeepLearning.AI, khoá 4: Data Modeling, Transformation, and Serving}: ghi danh miễn phí. Xem tuần 1, khoảng 10 giờ theo trang khoá: Inmon so với Kimball, Data Vault và OBT, ứng với Bài 8 và Bài 11.
\end{itemize}
% Kiểm: trang Coursera khoá 4 xem ngày 06/11/2026: "3 weeks to complete", "10 hours a week"; tuần 1 có video "Inmon VS Kimball", "Data Vault", "One Big Table".
% Kiểm: raw.githubusercontent.com/DataTalksClub/data-engineering-zoomcamp/main/04-analytics-engineering/01-analytics-engineering-basics.md có mục "Dimensional Modeling" theo Kimball, 06/11/2026.

% =====================================================================
\section{Kế hoạch theo tuần}

Mỗi tuần 9 giờ cho module; giờ thứ mười dành cho chương nghề nghiệp. Tuần 4 chỉ cần 3 giờ, nên 6 giờ module còn lại của tuần đó dành cho Module 4.

\begin{xltabular}{\textwidth}{@{}P{1.0cm} L P{4.4cm} P{0.8cm}@{}}
\toprule
\thd{tuần} & \thd{việc lõi} & \thd{thêm nếu còn giờ} & \thd{giờ} \\
\midrule
\endhead
\bottomrule
\endfoot
1 & Bài 1 tới Bài 7, chạy mọi khối code và đọc kỹ từng output; Lab 3.1 & chạy lại khối Làm thử của Bài 1 với \code{SET threads = 1} & 9 \\
2 & Bài 8 tới Bài 10; Lab 3.2 và Lab 3.3 & bài viết của Kaminsky & 9 \\
3 & Bài 11 và Bài 12; Lab 3.4 và Lab 3.5; Tự kiểm tra; bắt đầu thử thách kỹ năng & phần thêm nếu còn giờ của Lab 3.5; bài viết của Linstedt và của Margy Ross & 9 \\
4 & hết thử thách kỹ năng; mốc capstone 1 & DE Zoomcamp Module 4 unit 1; khoá 4 của DeepLearning.AI & 3 \\
\end{xltabular}
% Ước lượng ngày 06/11/2026, chưa đo trên người học. Tuần 1 khoảng 8 giờ: 23 trang và 10 khối code từ đầu module tới hết Bài 7, cộng khoảng 2 giờ cho Lab 3.1. Tuần 2 khoảng 9 giờ: Bài 8 tới Bài 10 khoảng 2,5 giờ, Lab 3.2 khoảng 4 giờ, Lab 3.3 khoảng 2,5 giờ. Tuần 3 khoảng 9 giờ: Bài 11 và Bài 12 khoảng 1,25 giờ, Lab 3.4 khoảng 2 giờ, Lab 3.5 khoảng 4 giờ, Tự kiểm tra 0,5 giờ, phần còn lại cho thử thách. Tuần 4 khoảng 3 giờ. Tổng khoảng 29 tới 31 giờ. Số trang lấy từ bản dựng tectonic cùng ngày.

Nếu thiếu giờ, đọc nhanh các phần sau thay vì bỏ hẳn:
\begin{itemize}
  \item Bài 1 và Bài 10, khi bạn qua được bài kiểm 15 phút ở thẻ module.
  \item Bài 11: đọc phần Giải thích và sơ đồ; phần còn lại để dành cho lúc phải chọn mô hình cho một kho mới.
  \item Khối thứ hai của Bài 6 và bước 3 của Lab 3.4: chỉ đọc output in sẵn, không chạy.
\end{itemize}
